% Options for packages loaded elsewhere
\PassOptionsToPackage{unicode}{hyperref}
\PassOptionsToPackage{hyphens}{url}
\PassOptionsToPackage{dvipsnames,svgnames,x11names}{xcolor}
%
\documentclass[
  10pt,
]{article}
\usepackage{amsmath,amssymb}
\usepackage{iftex}
\ifPDFTeX
  \usepackage[T1]{fontenc}
  \usepackage[utf8]{inputenc}
  \usepackage{textcomp} % provide euro and other symbols
\else % if luatex or xetex
  \usepackage{unicode-math} % this also loads fontspec
  \defaultfontfeatures{Scale=MatchLowercase}
  \defaultfontfeatures[\rmfamily]{Ligatures=TeX,Scale=1}
\fi
\usepackage{lmodern}
\ifPDFTeX\else
  % xetex/luatex font selection
\fi
% Use upquote if available, for straight quotes in verbatim environments
\IfFileExists{upquote.sty}{\usepackage{upquote}}{}
\IfFileExists{microtype.sty}{% use microtype if available
  \usepackage[]{microtype}
  \UseMicrotypeSet[protrusion]{basicmath} % disable protrusion for tt fonts
}{}
\makeatletter
\@ifundefined{KOMAClassName}{% if non-KOMA class
  \IfFileExists{parskip.sty}{%
    \usepackage{parskip}
  }{% else
    \setlength{\parindent}{0pt}
    \setlength{\parskip}{6pt plus 2pt minus 1pt}}
}{% if KOMA class
  \KOMAoptions{parskip=half}}
\makeatother
\usepackage{xcolor}
\usepackage{longtable,booktabs,array}
\usepackage{calc} % for calculating minipage widths
% Correct order of tables after \paragraph or \subparagraph
\usepackage{etoolbox}
\makeatletter
\patchcmd\longtable{\par}{\if@noskipsec\mbox{}\fi\par}{}{}
\makeatother
% Allow footnotes in longtable head/foot
\IfFileExists{footnotehyper.sty}{\usepackage{footnotehyper}}{\usepackage{footnote}}
\makesavenoteenv{longtable}
\setlength{\emergencystretch}{3em} % prevent overfull lines
\providecommand{\tightlist}{%
  \setlength{\itemsep}{0pt}\setlength{\parskip}{0pt}}
\setcounter{secnumdepth}{5}

\usepackage{amsmath,amssymb}
\usepackage{graphicx}
\usepackage{longtable,booktabs,array}
\usepackage{xcolor}
\usepackage[margin=2.0cm]{geometry}
\usepackage{microtype}
\setlength{\emergencystretch}{3em}
\setcounter{MaxMatrixCols}{20}
\pdfinfoomitdate=1
\pdftrailerid{}
\pdfsuppressptexinfo=-1
\renewcommand{\arraystretch}{1.12}
\usepackage{fvextra}
\RecustomVerbatimEnvironment{verbatim}{Verbatim}{breaklines,breakanywhere,fontsize=\small}
\fvset{breaklines,breakanywhere}
\allowdisplaybreaks
\sloppy
\ifLuaTeX
  \usepackage{selnolig}  % disable illegal ligatures
\fi
\IfFileExists{bookmark.sty}{\usepackage{bookmark}}{\usepackage{hyperref}}
\IfFileExists{xurl.sty}{\usepackage{xurl}}{} % add URL line breaks if available
\urlstyle{same}
\hypersetup{
  pdftitle={UM-4 model},
  pdfauthor={Underlying theory: Patrick L. Nash; Formulated and verified by Claude (Anthropic); Model of UM-4, 2026-10-25; Repository once-ere/Dirac\_claude, folder universe\_model\_4},
  colorlinks=true,
  linkcolor={blue},
  filecolor={Maroon},
  citecolor={Blue},
  urlcolor={blue},
  pdfcreator={LaTeX via pandoc}}

\title{UM-4 model}
\usepackage{etoolbox}
\makeatletter
\providecommand{\subtitle}[1]{% add subtitle to \maketitle
  \apptocmd{\@title}{\par {\large #1 \par}}{}{}
}
\makeatother
\subtitle{The universe model UM-4: the author's spacetime of eight
dimensions with the time warped like every other direction of the fold,
the fields and where they live, the complete action, and every
parameter}
\author{Underlying theory: Patrick L. Nash \and Formulated and verified
by Claude (Anthropic) \and Model of UM-4, 2026-10-25 \and Repository
once-ere/Dirac\_claude, folder universe\_model\_4}
\date{}

\begin{document}
\maketitle

{
\hypersetup{linkcolor=}
\setcounter{tocdepth}{3}
\tableofcontents
}
\hypertarget{the-fold-definition-and-examples}{%
\section{The fold: definition and
examples}\label{the-fold-definition-and-examples}}

\textbf{Definition.} The \emph{fold} is the seven-dimensional surface of
the author's eight-dimensional spacetime on which the cosine of the
angle \(z = 6Hx_0\) of the hidden direction vanishes:

\[
\mathcal{F} = \big\{X = (x_0, x_1, \dots, x_7):\ \cos(6Hx_0) = 0\big\},\qquad
\text{in one period } 0 < z < \pi\text{ of the metric: } z = \frac{\pi}{2},\quad x_0 = \frac{\pi}{12H} \tag{FD-1}
\] \textbf{Symbols in (FD-1).}

\begin{itemize}
\tightlist
\item
  \(\mathcal{F}\): The fold: the seven-dimensional surface
  \(\cos(6Hx_0) = 0\) of the eight-dimensional spacetime (FD-1), in one
  period \(0 < z < \pi\) the surface \(z = \pi/2\), \(x_0 = \pi/(12H)\),
  where \(\sigma = 0\); \(\vert_{\mathcal{F}}\) means restricted to the
  fold. Status: DERIVED.
\item
  \(X = \{x_0, \dots, x_7\}\): The author's coordinates of the
  eight-dimensional spacetime, \(X = \{x_0, x_1, \dots, x_7\}\) with the
  hidden direction first (Dirac\_Einstein-Lovelock\_4+4\_Nash.nb,
  section DEFCHART); \(X^M\) is the coordinate with the index \(M\) and
  \(\mathrm{d}^8X = \mathrm{d}x_0\,\mathrm{d}x_1\cdots\mathrm{d}x_7\)
  the coordinate volume element. Unit: length (each coordinate). Status:
  AUTHOR.
\item
  \(x_0\): The hidden space-like coordinate (listed first); the metric
  depends on it through \(z = 6Hx_0\), periodically with the period
  \(\pi/(6H)\). Unit: length. Status: AUTHOR.
\item
  \(x_1\): The first coordinate of ordinary 3-space. Unit: length.
  Status: AUTHOR.
\item
  \(x_7\): The third of the three extra time-like coordinates. Unit:
  length. Status: AUTHOR.
\item
  \(H\): The author's constant (an inverse length) that fixes the period
  \(\pi/(6H)\) of \(x_0\), the warp factor
  \(\mathrm{e}^{-2H\lvert\sigma\rvert}\) of the metric, the cosmological
  constant \(\Lambda = -21H^2\) and the tension
  \(\kappa\lambda_{\mathrm{F}} = 12H\) of the fold; in UM-4 a free
  parameter, bounded from below by the propagation of the gauge fields
  in the bulk. Unit: GeV (1/length). Value: \(H \gtrsim 2.5\) TeV (from
  the invisible width of the Z boson). Status: AUTHOR.
\item
  \(z\): The angle variable of the hidden direction, \(z = 6Hx_0\); one
  period of \(x_0\) is \(0 < z < \pi\), the fold is \(z = \pi/2\). Unit:
  dimensionless. Status: AUTHOR.
\item
  \(\pi\): The number pi, the ratio of the circumference of a circle to
  its diameter. Unit: dimensionless. Value: 3.141592654. Status:
  STANDARD.
\item
  \(\Lambda\): The eight-dimensional cosmological constant of UM-4,
  \(\Lambda = -21H^2\): the value for which the static block is a vacuum
  of Einstein gravity (its Einstein tensor is \(21H^2\) in every
  direction). The author writes the pure number \(\Lambda_{\mathrm{N}}\)
  in
  \(0 = -\Lambda_{\mathrm{N}} + \sum_p H^{-2p}w_{(p)}\mathcal{E}_{(p)}\),
  so that \(\Lambda = -\Lambda_{\mathrm{N}}H^2\). Unit: GeV\(^2\).
  Value: \(-21H^2\); the author's original \(-10H^2\); UM-3: \(-7H^2\).
  Status: AUTHOR (the term), DERIVED (the value).
\item
  \(\Lambda_{\mathrm{N}}\): The author's pure number in the vacuum
  equations of his notebook, \(\Lambda = -\Lambda_{\mathrm{N}}H^2\): his
  original value 10; 21 in UM-4 (Einstein gravity), 7 in UM-3. Unit:
  dimensionless. Value: 21 (the author's original 10). Status: AUTHOR.
\item
  \(w_{(p)}\): The Lovelock weights of the author's Einstein-Lovelock
  action, \(w_{(1)} = 1\) and \(w_{(2)}\), \(w_{(3)}\) of the
  Gauss-Bonnet and the cubic densities. UM-4 sets
  \(w_{(2)} = w_{(3)} = 0\) (Einstein gravity): the only weights for
  which gravitons propagate on the light cone on every background, not
  only around the static block. The author's original weights
  \(w_{(2)} = 1/32\), \(w_{(3)} = 1/1152\)
  (Dirac\_Einstein-Lovelock\_4+4\_Nash.nb, In{[}134{]}) keep the
  gravitons on the light cone only around the static block (with
  \(\Lambda = -161H^2/16\)); UM-3 used \((1, 1/20, 1/360)\). Unit:
  dimensionless. Value: \((1, 0, 0)\) in UM-4; the author's original
  \((1, 1/32, 1/1152)\). Status: AUTHOR (the form), DERIVED (the values
  that the order of 2026-10-25 selects).
\end{itemize}

\(X = \{x_0, \dots, x_7\}\) are the author's coordinates: \(x_0\) the
hidden space-like direction, \(x_1, x_2, x_3\) ordinary space, \(x_4\)
the time, \(x_5, x_6, x_7\) the three extra times; \(H\) is the author's
constant (an inverse length). The metric of UM-4 depends on \(x_0\) only
through \(\cot^2z\) and \((\sin^2z)^{1/6}\), so it repeats with the
period \(\pi/(6H)\) of \(x_0\), and every period \(0 < z < \pi\)
contains exactly one surface of \(\mathcal{F}\), at \(z = \pi/2\):
\(\mathcal{F}\) is the union of these copies, one per period (example
2), and ``the fold'' without a qualifier is the copy in the period that
contains our world. A point of the fold is labelled by the seven
coordinates \(x_1, \dots, x_7\), and the metric restricted to the fold
(\(\mathrm{d}x_0 = 0\), \(z = \pi/2\)) is

\[
g_{MN}\,\mathrm{d}X^M\mathrm{d}X^N\big\vert_{\mathcal{F}} = \mathrm{e}^{2a_4}\big(\mathrm{d}x_1^2 + \mathrm{d}x_2^2 + \mathrm{d}x_3^2\big)
 - \mathrm{d}x_4^2 - \mathrm{e}^{-2a_4}b_4^2\big(\mathrm{d}x_5^2 + \mathrm{d}x_6^2 + \mathrm{d}x_7^2\big) \tag{FD-2}
\] \textbf{Symbols in (FD-2).}

\begin{itemize}
\tightlist
\item
  \(g_{MN}\): The metric tensor of the eight-dimensional spacetime
  (signature \(+,+,+,+,-,-,-,-\) in the order \(x_0, \dots, x_7\)):
  \(g_{MN}\,\mathrm{d}X^M\mathrm{d}X^N\) is the squared interval,
  \(g^{MN}\) its inverse, \(\lvert g\rvert\) the absolute value of its
  determinant (\(\sqrt{\lvert g\rvert}\,\mathrm{d}^8X\) is the proper
  volume element). Unit: dimensionless. Status: AUTHOR.
\item
  \(M\): A coordinate index of eight dimensions (also with a numbered
  copy \(M_1, M_2, \dots\)): it runs over \(x_0, x_1, \dots, x_7\); a
  repeated index is summed. Status: STANDARD.
\item
  \(N\): A coordinate index of eight dimensions (also
  \(N_1, N_2, \dots\)), running over \(x_0, \dots, x_7\); a repeated
  index is summed. Status: STANDARD.
\item
  \(\mathrm{d}\): The differential: \(\mathrm{d}\) in front of a
  quantity is an infinitesimal change of it (as \(\mathrm{d}\sigma\), a
  small step of \(\sigma\)); the quotient of two such changes is a
  derivative; \(\mathrm{d}\) with an exponent, in front of several
  variables, is their volume element, the product of their
  differentials. Status: STANDARD.
\item
  \(X = \{x_0, \dots, x_7\}\): The author's coordinates of the
  eight-dimensional spacetime, \(X = \{x_0, x_1, \dots, x_7\}\) with the
  hidden direction first (Dirac\_Einstein-Lovelock\_4+4\_Nash.nb,
  section DEFCHART); \(X^M\) is the coordinate with the index \(M\) and
  \(\mathrm{d}^8X = \mathrm{d}x_0\,\mathrm{d}x_1\cdots\mathrm{d}x_7\)
  the coordinate volume element. Unit: length (each coordinate). Status:
  AUTHOR.
\item
  \(\mathcal{F}\): The fold: the seven-dimensional surface
  \(\cos(6Hx_0) = 0\) of the eight-dimensional spacetime (FD-1), in one
  period \(0 < z < \pi\) the surface \(z = \pi/2\), \(x_0 = \pi/(12H)\),
  where \(\sigma = 0\); \(\vert_{\mathcal{F}}\) means restricted to the
  fold. Status: DERIVED.
\item
  \(\mathrm{e}\): Euler's number, the base of the natural logarithm;
  \(\mathrm{e}^{X}\) is the exponential of \(X\). Unit: dimensionless.
  Value: 2.718281828. Status: STANDARD.
\item
  \(a_4\): The author's metric function of the time \(x_4\) that scales
  ordinary 3-space: on the fold \(\mathrm{e}^{a_4}\) is the scale factor
  of cosmology and \(a_4'\) the Hubble rate. Unit: dimensionless.
  Status: AUTHOR.
\item
  \(b_4\): The author's metric function of the time \(x_4\) in the
  extra-time components, \(b_4 > 0\), never identically 1 (the author's
  order of 2026-10-24); the scale factor (the factor of the lengths) of
  the extra times is \(\mathrm{e}^{-a_4}b_4\) on the fold and
  \(\mathrm{e}^{-a_4}b_4\,\mathrm{e}^{-H\lvert\sigma\rvert}\) at the
  distance \(\sigma\) from it (the square root of the warp factor times
  \(\mathrm{e}^{-a_4}b_4\)). Unit: dimensionless. Status: AUTHOR.
\item
  \(x_1\): The first coordinate of ordinary 3-space. Unit: length.
  Status: AUTHOR.
\item
  \(x_2\): The second coordinate of ordinary 3-space. Unit: length.
  Status: AUTHOR.
\item
  \(x_3\): The third coordinate of ordinary 3-space. Unit: length.
  Status: AUTHOR.
\item
  \(x_4\): The time; a prime on a function of \(x_4\) alone is
  \(\mathrm{d}/\mathrm{d}x_4\). Unit: length (= time, \(c = 1\)).
  Status: AUTHOR.
\item
  \(x_5\): The first of the three extra time-like coordinates. Unit:
  length. Status: AUTHOR.
\item
  \(x_6\): The second of the three extra time-like coordinates. Unit:
  length. Status: AUTHOR.
\item
  \(x_7\): The third of the three extra time-like coordinates. Unit:
  length. Status: AUTHOR.
\item
  \({}'\): The derivative of a function of the time \(x_4\) alone with
  respect to \(x_4\) (\(a_4' = \mathrm{d}a_4/\mathrm{d}x_4\); two primes
  the second derivative); on a function of \(\sigma\) alone (a profile)
  the derivative with respect to \(\sigma\). Status: STANDARD.
\end{itemize}

with three space-like and four time-like directions; \(a_4\) and \(b_4\)
are the author's two functions of the time \(x_4\). Notation: \(\ln\) is
the natural logarithm, \(\mathrm{e}\) Euler's number,
\(\exp(y) = \mathrm{e}^y\), \(\sin\), \(\cos\) and \(\cot = \cos/\sin\)
the trigonometric functions, \(\sinh\) and \(\tanh\) the hyperbolic
functions, \(\lvert y\rvert\) the absolute value; \(f(0^+)\) and
\(f(0^-)\) are the limits of \(f(\sigma)\) as \(\sigma\) approaches 0
from above and from below.

\textbf{Why this surface is called a fold.} Measured by the proper
distance along the hidden direction,

\[
\sigma = \mathrm{sgn}\Big(\frac{\pi}{2} - z\Big)\frac{\ln\sin z}{6H},\qquad \mathrm{d}\sigma = \lvert\cot z\rvert\,\mathrm{d}x_0,\qquad
\sin z = \mathrm{e}^{-6H\lvert\sigma\rvert} \tag{FD-3}
\] \textbf{Symbols in (FD-3).}

\begin{itemize}
\tightlist
\item
  \(\sigma\): The proper coordinate of the hidden direction:
  \(\mathrm{d}\sigma = \lvert\cot z\rvert\,\mathrm{d}x_0\),
  \(\sigma = \mathrm{sgn}(\pi/2 - z)\ln(\sin z)/(6H)\) on one period
  \(0 < z < \pi\); \(\sigma = 0\) is the fold \(z = \pi/2\),
  \(\sigma \to \mp\infty\) the tips \(z \to 0, \pi\); one period of
  \(x_0\) is the whole line \(-\infty < \sigma < \infty\). Unit: length.
  Status: DERIVED.
\item
  \(\mathrm{sgn}\): The sign function: \(+1\) for a positive argument
  and \(-1\) for a negative one. Status: STANDARD.
\item
  \(\pi\): The number pi, the ratio of the circumference of a circle to
  its diameter. Unit: dimensionless. Value: 3.141592654. Status:
  STANDARD.
\item
  \(z\): The angle variable of the hidden direction, \(z = 6Hx_0\); one
  period of \(x_0\) is \(0 < z < \pi\), the fold is \(z = \pi/2\). Unit:
  dimensionless. Status: AUTHOR.
\item
  \(H\): The author's constant (an inverse length) that fixes the period
  \(\pi/(6H)\) of \(x_0\), the warp factor
  \(\mathrm{e}^{-2H\lvert\sigma\rvert}\) of the metric, the cosmological
  constant \(\Lambda = -21H^2\) and the tension
  \(\kappa\lambda_{\mathrm{F}} = 12H\) of the fold; in UM-4 a free
  parameter, bounded from below by the propagation of the gauge fields
  in the bulk. Unit: GeV (1/length). Value: \(H \gtrsim 2.5\) TeV (from
  the invisible width of the Z boson). Status: AUTHOR.
\item
  \(\mathrm{d}\): The differential: \(\mathrm{d}\) in front of a
  quantity is an infinitesimal change of it (as \(\mathrm{d}\sigma\), a
  small step of \(\sigma\)); the quotient of two such changes is a
  derivative; \(\mathrm{d}\) with an exponent, in front of several
  variables, is their volume element, the product of their
  differentials. Status: STANDARD.
\item
  \(x_0\): The hidden space-like coordinate (listed first); the metric
  depends on it through \(z = 6Hx_0\), periodically with the period
  \(\pi/(6H)\). Unit: length. Status: AUTHOR.
\item
  \(\mathrm{e}\): Euler's number, the base of the natural logarithm;
  \(\mathrm{e}^{X}\) is the exponential of \(X\). Unit: dimensionless.
  Value: 2.718281828. Status: STANDARD.
\item
  \(\infty\): Infinity: as a limit of an integral or a sum it means that
  the range has no bound. Status: STANDARD.
\end{itemize}

(the fold is \(\sigma = 0\); the half-period \(0 < z < \pi/2\) is
\(\sigma < 0\), the half-period \(\pi/2 < z < \pi\) is \(\sigma > 0\);
the two ends of the period, the tips \(z \to 0\) and \(z \to \pi\), lie
at \(\sigma \to -\infty\) and \(\sigma \to +\infty\)), the metric of
UM-4 is

\[
g_{MN}\,\mathrm{d}X^M\mathrm{d}X^N = \mathrm{d}\sigma^2 + \mathrm{e}^{-2H\lvert\sigma\rvert}\Big[\mathrm{e}^{2a_4}\big(\mathrm{d}x_1^2
 + \mathrm{d}x_2^2 + \mathrm{d}x_3^2\big) - \mathrm{d}x_4^2 - \mathrm{e}^{-2a_4}b_4^2\big(\mathrm{d}x_5^2 + \mathrm{d}x_6^2 + \mathrm{d}x_7^2\big)\Big] \tag{FD-4}
\] \textbf{Symbols in (FD-4).}

\begin{itemize}
\tightlist
\item
  \(g_{MN}\): The metric tensor of the eight-dimensional spacetime
  (signature \(+,+,+,+,-,-,-,-\) in the order \(x_0, \dots, x_7\)):
  \(g_{MN}\,\mathrm{d}X^M\mathrm{d}X^N\) is the squared interval,
  \(g^{MN}\) its inverse, \(\lvert g\rvert\) the absolute value of its
  determinant (\(\sqrt{\lvert g\rvert}\,\mathrm{d}^8X\) is the proper
  volume element). Unit: dimensionless. Status: AUTHOR.
\item
  \(M\): A coordinate index of eight dimensions (also with a numbered
  copy \(M_1, M_2, \dots\)): it runs over \(x_0, x_1, \dots, x_7\); a
  repeated index is summed. Status: STANDARD.
\item
  \(N\): A coordinate index of eight dimensions (also
  \(N_1, N_2, \dots\)), running over \(x_0, \dots, x_7\); a repeated
  index is summed. Status: STANDARD.
\item
  \(\mathrm{d}\): The differential: \(\mathrm{d}\) in front of a
  quantity is an infinitesimal change of it (as \(\mathrm{d}\sigma\), a
  small step of \(\sigma\)); the quotient of two such changes is a
  derivative; \(\mathrm{d}\) with an exponent, in front of several
  variables, is their volume element, the product of their
  differentials. Status: STANDARD.
\item
  \(X = \{x_0, \dots, x_7\}\): The author's coordinates of the
  eight-dimensional spacetime, \(X = \{x_0, x_1, \dots, x_7\}\) with the
  hidden direction first (Dirac\_Einstein-Lovelock\_4+4\_Nash.nb,
  section DEFCHART); \(X^M\) is the coordinate with the index \(M\) and
  \(\mathrm{d}^8X = \mathrm{d}x_0\,\mathrm{d}x_1\cdots\mathrm{d}x_7\)
  the coordinate volume element. Unit: length (each coordinate). Status:
  AUTHOR.
\item
  \(\sigma\): The proper coordinate of the hidden direction:
  \(\mathrm{d}\sigma = \lvert\cot z\rvert\,\mathrm{d}x_0\),
  \(\sigma = \mathrm{sgn}(\pi/2 - z)\ln(\sin z)/(6H)\) on one period
  \(0 < z < \pi\); \(\sigma = 0\) is the fold \(z = \pi/2\),
  \(\sigma \to \mp\infty\) the tips \(z \to 0, \pi\); one period of
  \(x_0\) is the whole line \(-\infty < \sigma < \infty\). Unit: length.
  Status: DERIVED.
\item
  \(\mathrm{e}\): Euler's number, the base of the natural logarithm;
  \(\mathrm{e}^{X}\) is the exponential of \(X\). Unit: dimensionless.
  Value: 2.718281828. Status: STANDARD.
\item
  \(H\): The author's constant (an inverse length) that fixes the period
  \(\pi/(6H)\) of \(x_0\), the warp factor
  \(\mathrm{e}^{-2H\lvert\sigma\rvert}\) of the metric, the cosmological
  constant \(\Lambda = -21H^2\) and the tension
  \(\kappa\lambda_{\mathrm{F}} = 12H\) of the fold; in UM-4 a free
  parameter, bounded from below by the propagation of the gauge fields
  in the bulk. Unit: GeV (1/length). Value: \(H \gtrsim 2.5\) TeV (from
  the invisible width of the Z boson). Status: AUTHOR.
\item
  \(a_4\): The author's metric function of the time \(x_4\) that scales
  ordinary 3-space: on the fold \(\mathrm{e}^{a_4}\) is the scale factor
  of cosmology and \(a_4'\) the Hubble rate. Unit: dimensionless.
  Status: AUTHOR.
\item
  \(b_4\): The author's metric function of the time \(x_4\) in the
  extra-time components, \(b_4 > 0\), never identically 1 (the author's
  order of 2026-10-24); the scale factor (the factor of the lengths) of
  the extra times is \(\mathrm{e}^{-a_4}b_4\) on the fold and
  \(\mathrm{e}^{-a_4}b_4\,\mathrm{e}^{-H\lvert\sigma\rvert}\) at the
  distance \(\sigma\) from it (the square root of the warp factor times
  \(\mathrm{e}^{-a_4}b_4\)). Unit: dimensionless. Status: AUTHOR.
\item
  \(x_1\): The first coordinate of ordinary 3-space. Unit: length.
  Status: AUTHOR.
\item
  \(x_2\): The second coordinate of ordinary 3-space. Unit: length.
  Status: AUTHOR.
\item
  \(x_3\): The third coordinate of ordinary 3-space. Unit: length.
  Status: AUTHOR.
\item
  \(x_4\): The time; a prime on a function of \(x_4\) alone is
  \(\mathrm{d}/\mathrm{d}x_4\). Unit: length (= time, \(c = 1\)).
  Status: AUTHOR.
\item
  \(x_5\): The first of the three extra time-like coordinates. Unit:
  length. Status: AUTHOR.
\item
  \(x_6\): The second of the three extra time-like coordinates. Unit:
  length. Status: AUTHOR.
\item
  \(x_7\): The third of the three extra time-like coordinates. Unit:
  length. Status: AUTHOR.
\end{itemize}

It depends on \(\sigma\) only through \(\lvert\sigma\rvert\). The
function \(\mathcal{W}(\sigma) = \mathrm{e}^{-2H\lvert\sigma\rvert}\),
which multiplies the squared steps of all seven directions of the fold
(ordinary space, the time and the extra times), is the warp factor of
the metric (its square root \(\mathrm{e}^{-H\lvert\sigma\rvert}\)
multiplies the lengths). The time \(x_4\) is warped like the other
directions (the author's order of 2026-10-25, 13:15 PDT; in UM-3 it was
not), so that every surface \(\sigma = \text{const}\) has the light cone
of the fold. Four properties make the surface \(\sigma = 0\) a fold, and
together they single it out within the period:

\begin{enumerate}
\def\labelenumi{\arabic{enumi}.}
\tightlist
\item
  \emph{It is a mirror.} The reflection \(\sigma \to -\sigma\) (that is
  \(z \to \pi - z\), or \(x_0 \to \pi/(6H) - x_0\)) maps the metric to
  itself, and the fold is the only surface of the period that the
  reflection leaves in place: the two halves of the period are mirror
  images of each other, joined along the fold, as the two halves of a
  folded sheet of paper are joined along the crease.
\item
  \emph{The warp factor is largest on it.} The warp factor
  \((\sin^2z)^{1/6} = \mathrm{e}^{-2H\lvert\sigma\rvert}\) equals 1 on
  the fold and decreases on both sides.
\item
  \emph{The metric has a kink on it.} The metric is continuous across
  the fold, but its first derivative along \(\sigma\) jumps. The
  extrinsic curvature of the surfaces \(\sigma = \text{const}\) changes
  sign across the fold,
\end{enumerate}

\[
K^a{}_b = \frac12\,g^{ac}\,\partial_\sigma g_{cb} = -H\,\mathrm{sgn}(\sigma)\,\delta^a{}_b\quad(\text{all seven directions of the fold}) \tag{FD-5}
\] \textbf{Symbols in (FD-5).}

\begin{itemize}
\tightlist
\item
  \(K^a{}_b\): The extrinsic curvature of a surface of constant
  \(\sigma\) of the hidden direction,
  \(K^a{}_b = \frac12 g^{ac}\partial_\sigma g_{cb}\) for \(a, b\) along
  the surface (FD-5): \(-H\,\mathrm{sgn}(\sigma)\) along all seven
  directions of the fold for the author's form of the metric of UM-4;
  \(K = K^c{}_c\) its trace. Unit: GeV (1/length). Status: DERIVED.
\item
  \(a\): A frame (orthonormal) index running over \(0, 1, \dots, 7\)
  (frame direction \((0)\) along \(\sigma\), \((1), (2), (3)\) along
  ordinary space, \((4)\) along the time, \((5), (6), (7)\) along the
  extra times); raised and lowered with the frame metric
  \(\eta_{ab} = \mathrm{diag}(1, 1, 1, 1, -1, -1, -1, -1)\); on the fold
  (\(a, b \in \{1, \dots, 7\}\)) the frame of the fold; in the tensors
  of the fold written with its coordinates \(x^a\) (its metric
  \(\hat g_{ab}\) and the tensors built from it) a coordinate index of
  the fold, \(a \in \{1, \dots, 7\}\); a repeated index is summed.
  Status: STANDARD.
\item
  \(b\): A frame index (as \(a\)). Status: STANDARD.
\item
  \(c\): A frame index (as \(a\)). Status: STANDARD.
\item
  \(g_{MN}\): The metric tensor of the eight-dimensional spacetime
  (signature \(+,+,+,+,-,-,-,-\) in the order \(x_0, \dots, x_7\)):
  \(g_{MN}\,\mathrm{d}X^M\mathrm{d}X^N\) is the squared interval,
  \(g^{MN}\) its inverse, \(\lvert g\rvert\) the absolute value of its
  determinant (\(\sqrt{\lvert g\rvert}\,\mathrm{d}^8X\) is the proper
  volume element). Unit: dimensionless. Status: AUTHOR.
\item
  \(\partial\): A partial derivative: the derivative with respect to one
  variable at fixed values of the others (\(\partial_\sigma\) with
  respect to \(\sigma\); \(\partial_M\) and \(\partial_\mu\) with
  respect to the coordinate that carries the same index, of the eight
  dimensions or of the fold). Status: STANDARD.
\item
  \(\sigma\): The label \(\sigma\) in \(\partial_\sigma\): the
  derivative with respect to the proper hidden coordinate \(\sigma\).
  Status: STANDARD.
\item
  \(\sigma\): The proper coordinate of the hidden direction:
  \(\mathrm{d}\sigma = \lvert\cot z\rvert\,\mathrm{d}x_0\),
  \(\sigma = \mathrm{sgn}(\pi/2 - z)\ln(\sin z)/(6H)\) on one period
  \(0 < z < \pi\); \(\sigma = 0\) is the fold \(z = \pi/2\),
  \(\sigma \to \mp\infty\) the tips \(z \to 0, \pi\); one period of
  \(x_0\) is the whole line \(-\infty < \sigma < \infty\). Unit: length.
  Status: DERIVED.
\item
  \(H\): The author's constant (an inverse length) that fixes the period
  \(\pi/(6H)\) of \(x_0\), the warp factor
  \(\mathrm{e}^{-2H\lvert\sigma\rvert}\) of the metric, the cosmological
  constant \(\Lambda = -21H^2\) and the tension
  \(\kappa\lambda_{\mathrm{F}} = 12H\) of the fold; in UM-4 a free
  parameter, bounded from below by the propagation of the gauge fields
  in the bulk. Unit: GeV (1/length). Value: \(H \gtrsim 2.5\) TeV (from
  the invisible width of the Z boson). Status: AUTHOR.
\item
  \(\mathrm{sgn}\): The sign function: \(+1\) for a positive argument
  and \(-1\) for a negative one. Status: STANDARD.
\item
  \(\delta\): The Kronecker delta: with two indices \(\delta^M_N = 1\)
  if \(M = N\) and 0 otherwise (likewise \(\delta_{X\mathrm{L}}\) is 1
  if the label \(X\) of a half of dirac16complex is L and 0 otherwise);
  with two index strings of equal length the generalised Kronecker delta
  \(\delta^{M_1M_2\cdots}_{N_1N_2\cdots}\), the determinant of the
  square matrix of the Kronecker deltas \(\delta^{M_a}_{N_b}\) (\(+1\)
  or \(-1\) if the upper indices are an even or odd permutation of the
  lower ones, 0 otherwise). Status: STANDARD.
\item
  \(\eta_{ab}\): The flat metric of the frame,
  \(\eta = \mathrm{diag}(1, 1, 1, 1, -1, -1, -1, -1)\) (eta4488 of
  4+4.nb); \(\eta^{ab}\) is its inverse (the same matrix). Unit:
  dimensionless. Status: AUTHOR.
\item
  \(\hat g_{ab}\): The metric of the seven-dimensional fold (FD-2) in
  its coordinates \(x_1, \dots, x_7\) (or in its frame), \(\hat g^{ab}\)
  its inverse; \(\hat h_{\mu\nu}\) is its part along the four observable
  directions. Unit: dimensionless. Status: AUTHOR.
\item
  \(\mathrm{diag}\): A diagonal matrix with the listed entries on its
  diagonal (in the order of the coordinates \(x_0, x_1, \dots, x_7\) for
  a tensor of eight dimensions, \(x_1, \dots, x_7\) for a tensor of the
  fold, \(x_1, \dots, x_4\) for a four-dimensional tensor). Status:
  STANDARD.
\item
  \(x\): The coordinates of the fold: \(x_1, x_2, x_3\) of ordinary
  space, \(x_4\) the time, \(x_5, x_6, x_7\) the extra times; a field of
  the fold is a function of the four observable coordinates
  \(x = (x_1, x_2, x_3, x_4)\), and
  \(\mathrm{d}^4x = \mathrm{d}x_1\,\mathrm{d}x_2\,\mathrm{d}x_3\,\mathrm{d}x_4\).
  Unit: length (each coordinate). Status: AUTHOR.
\end{itemize}

(\(a, b\) frame indices along the fold); it jumps by \(-2H\) in every
direction of the fold. The field equations of gravity of UM-4,
\(\mathcal{E}_{(1)}{}^M{}_N + \Lambda\delta^M_N = \kappa\mathcal{T}^M{}_N\)
(Einstein gravity: the author's Einstein-Lovelock equations with the
weights of the Gauss-Bonnet and cubic terms set to zero;
\(\mathcal{E}_{(1)}\) the Einstein tensor, \(\Lambda = -21H^2\) the
cosmological constant, \(\kappa = 8\pi G_8\) the gravitational coupling
of eight dimensions, \(\mathcal{T}^M{}_N\) the energy-momentum tensor,
\(\delta^M_N\) the Kronecker delta), are linear in the second
\(\sigma\)-derivatives of the metric, so the jump puts a Dirac delta
function \(\delta_{\mathrm{D}}(\sigma)\) into them: the metric solves
them only if the fold carries a thin source \(\mathcal{S}^a{}_b\), an
energy-momentum tensor concentrated on \(\sigma = 0\), given by the
integral of the equations across the fold (the junction condition of
Israel),

\[
\kappa\,\mathcal{S}^a{}_b = \int_{0^-}^{0^+}\mathrm{d}\sigma\,\mathcal{E}_{(1)}{}^a{}_b
 = -\Big(\big(K^a{}_b\big\vert_{0^+} - K^a{}_b\big\vert_{0^-}\big) - \delta^a{}_b\big(K\big\vert_{0^+} - K\big\vert_{0^-}\big)\Big) \tag{FD-11}
\] \textbf{Symbols in (FD-11).}

\begin{itemize}
\tightlist
\item
  \(\kappa\): The eight-dimensional gravitational coupling,
  \(\kappa = 8\pi G_8\). Unit: GeV\(^{-6}\). Status: AUTHOR.
\item
  \(\mathcal{S}^a{}_b\): The thin source of the fold: the
  energy-momentum tensor concentrated on the fold, integrated across it,
  \(\kappa\mathcal{S}^a{}_b = \int_{0^-}^{0^+}\mathrm{d}\sigma\,\mathcal{E}_{(1)}{}^a{}_b\)
  (the junction condition of Israel); \(\mathcal{S}^4{}_4\) is minus its
  energy density, \(\mathcal{S}^1{}_1\) and \(\mathcal{S}^5{}_5\) its
  pressures along ordinary space and along the extra times (frame
  components along the fold, eight-dimensional normalisation: an energy
  per unit volume of the six other directions of the fold). Unit:
  GeV\(^7\). Value: \(\kappa\mathcal{S}^a{}_b = -12H\delta^a{}_b\) for
  every \(a_4\), \(b_4\) of the author's form. Status: DERIVED.
\item
  \(a\): A frame (orthonormal) index running over \(0, 1, \dots, 7\)
  (frame direction \((0)\) along \(\sigma\), \((1), (2), (3)\) along
  ordinary space, \((4)\) along the time, \((5), (6), (7)\) along the
  extra times); raised and lowered with the frame metric
  \(\eta_{ab} = \mathrm{diag}(1, 1, 1, 1, -1, -1, -1, -1)\); on the fold
  (\(a, b \in \{1, \dots, 7\}\)) the frame of the fold; in the tensors
  of the fold written with its coordinates \(x^a\) (its metric
  \(\hat g_{ab}\) and the tensors built from it) a coordinate index of
  the fold, \(a \in \{1, \dots, 7\}\); a repeated index is summed.
  Status: STANDARD.
\item
  \(b\): A frame index (as \(a\)). Status: STANDARD.
\item
  \(\mathrm{d}\): The differential: \(\mathrm{d}\) in front of a
  quantity is an infinitesimal change of it (as \(\mathrm{d}\sigma\), a
  small step of \(\sigma\)); the quotient of two such changes is a
  derivative; \(\mathrm{d}\) with an exponent, in front of several
  variables, is their volume element, the product of their
  differentials. Status: STANDARD.
\item
  \(\sigma\): The proper coordinate of the hidden direction:
  \(\mathrm{d}\sigma = \lvert\cot z\rvert\,\mathrm{d}x_0\),
  \(\sigma = \mathrm{sgn}(\pi/2 - z)\ln(\sin z)/(6H)\) on one period
  \(0 < z < \pi\); \(\sigma = 0\) is the fold \(z = \pi/2\),
  \(\sigma \to \mp\infty\) the tips \(z \to 0, \pi\); one period of
  \(x_0\) is the whole line \(-\infty < \sigma < \infty\). Unit: length.
  Status: DERIVED.
\item
  \(\mathcal{E}_{(p)}{}^M{}_N\): The Lovelock tensor of order \(p\),
  \(\mathcal{E}_{(p)}{}^M{}_N = -2^{-(p+1)}\,\delta^{MP_1\cdots P_{2p}}_{NQ_1\cdots Q_{2p}}\,R^{Q_1Q_2}{}_{P_1P_2}\cdots R^{Q_{2p-1}Q_{2p}}{}_{P_{2p-1}P_{2p}}\)
  (with both indices down it is \(1/\sqrt{\lvert g\rvert}\) times the
  variation of \(\sqrt{\lvert g\rvert}\,\mathcal{L}_{(p)}\) with respect
  to \(g^{MN}\); \(\mathcal{E}_{(1)}\) the Einstein tensor,
  \(\mathcal{E}_{(2)}\) the Gauss-Bonnet tensor, \(\mathcal{E}_{(3)}\)
  the cubic Lovelock tensor); with frame indices its frame components.
  Unit: 1/length\(^{2p}\). Status: STANDARD.
\item
  \(p\): The Lovelock order \(p = 1, 2, 3\) (Einstein, Gauss-Bonnet,
  cubic), as an index and as a number (UM-4 keeps only \(p = 1\)); in
  the Kohn-Sham equations also a second label of a Kohn-Sham orbital,
  with the Kronecker delta of the two labels. Status: STANDARD.
\item
  \(K^a{}_b\): The extrinsic curvature of a surface of constant
  \(\sigma\) of the hidden direction,
  \(K^a{}_b = \frac12 g^{ac}\partial_\sigma g_{cb}\) for \(a, b\) along
  the surface (FD-5): \(-H\,\mathrm{sgn}(\sigma)\) along all seven
  directions of the fold for the author's form of the metric of UM-4;
  \(K = K^c{}_c\) its trace. Unit: GeV (1/length). Status: DERIVED.
\item
  \(\delta\): The Kronecker delta: with two indices \(\delta^M_N = 1\)
  if \(M = N\) and 0 otherwise (likewise \(\delta_{X\mathrm{L}}\) is 1
  if the label \(X\) of a half of dirac16complex is L and 0 otherwise);
  with two index strings of equal length the generalised Kronecker delta
  \(\delta^{M_1M_2\cdots}_{N_1N_2\cdots}\), the determinant of the
  square matrix of the Kronecker deltas \(\delta^{M_a}_{N_b}\) (\(+1\)
  or \(-1\) if the upper indices are an even or odd permutation of the
  lower ones, 0 otherwise). Status: STANDARD.
\item
  \(G_8\): The eight-dimensional Newton constant
  (\(\kappa = 8\pi G_8\)). Unit: GeV\(^{-6}\). Status: AUTHOR.
\item
  \(R^{MN}{}_{PQ},\ R_{MN},\ R\): The curvature of the metric
  \(g_{MN}\): with four indices the Riemann tensor \(R^M{}_{NPQ}\) (in
  the convention in which the commutator of two covariant derivatives,
  acting on a vector, is the Riemann tensor, with its last two indices
  those of the two derivatives, contracted with the vector), with two
  indices the Ricci tensor \(R_{MN} = R^P{}_{MPN}\), without indices the
  scalar curvature \(R = g^{MN}R_{MN}\); with frame indices
  \(R^{(a)(b)}{}_{(c)(d)}\) its frame components,
  \(R^{(a)(b)}{}_{(a)(b)}\) (no sum) the sectional curvature of the
  frame plane \((a, b)\). Unit: 1/length\(^2\). Status: STANDARD.
\item
  \(\mathcal{L}_{(p)}\): The Lovelock density of order \(p\),
  \(\mathcal{L}_{(p)} = 2^{-p}\,\delta^{P_1\cdots P_{2p}}_{Q_1\cdots Q_{2p}}\,R^{Q_1Q_2}{}_{P_1P_2}\cdots R^{Q_{2p-1}Q_{2p}}{}_{P_{2p-1}P_{2p}}\),
  a scalar built from \(p\) curvature tensors: \(\mathcal{L}_{(1)} = R\)
  (Einstein-Hilbert), \(\mathcal{L}_{(2)}\) the Gauss-Bonnet density,
  \(\mathcal{L}_{(3)}\) the cubic Lovelock density. Unit:
  1/length\(^{2p}\). Status: STANDARD.
\end{itemize}

(\(K = K^c{}_c\) the trace of the extrinsic curvature), with
\(\mathcal{S}^4{}_4 = -(\text{its energy density})\),
\(\mathcal{S}^1{}_1 = \mathcal{S}^2{}_2 = \mathcal{S}^3{}_3\) its
pressure along ordinary space and
\(\mathcal{S}^5{}_5 = \mathcal{S}^6{}_6 = \mathcal{S}^7{}_7\) its
pressure along the extra times; \(\mathcal{S}\) is \(\mathcal{T}\)
integrated across the fold, an energy per unit volume of the six other
directions of the fold (unit GeV\(^7\)). For the static block (\(a_4\)
and \(b_4\) constant) and, to the first order in the curvature of the
fold divided by \(H^2\), for every fold

\[
\begin{aligned}
&\text{static block:}\quad \kappa\mathcal{S}^a{}_b = -12H\,\delta^a{}_b\quad\text{(a pure tension)},\qquad \kappa\lambda_{\mathrm{F}} = 12H,\\
&\text{any fold:}\quad \kappa\mathcal{S}^a{}_b = -12H\,\delta^a{}_b + \frac{2}{5H}\,\hat G^a{}_b + \dots
\end{aligned} \tag{FD-12}
\] \textbf{Symbols in (FD-12).}

\begin{itemize}
\tightlist
\item
  \(\kappa\): The eight-dimensional gravitational coupling,
  \(\kappa = 8\pi G_8\). Unit: GeV\(^{-6}\). Status: AUTHOR.
\item
  \(\mathcal{S}^a{}_b\): The thin source of the fold: the
  energy-momentum tensor concentrated on the fold, integrated across it,
  \(\kappa\mathcal{S}^a{}_b = \int_{0^-}^{0^+}\mathrm{d}\sigma\,\mathcal{E}_{(1)}{}^a{}_b\)
  (the junction condition of Israel); \(\mathcal{S}^4{}_4\) is minus its
  energy density, \(\mathcal{S}^1{}_1\) and \(\mathcal{S}^5{}_5\) its
  pressures along ordinary space and along the extra times (frame
  components along the fold, eight-dimensional normalisation: an energy
  per unit volume of the six other directions of the fold). Unit:
  GeV\(^7\). Value: \(\kappa\mathcal{S}^a{}_b = -12H\delta^a{}_b\) for
  every \(a_4\), \(b_4\) of the author's form. Status: DERIVED.
\item
  \(a\): A frame (orthonormal) index running over \(0, 1, \dots, 7\)
  (frame direction \((0)\) along \(\sigma\), \((1), (2), (3)\) along
  ordinary space, \((4)\) along the time, \((5), (6), (7)\) along the
  extra times); raised and lowered with the frame metric
  \(\eta_{ab} = \mathrm{diag}(1, 1, 1, 1, -1, -1, -1, -1)\); on the fold
  (\(a, b \in \{1, \dots, 7\}\)) the frame of the fold; in the tensors
  of the fold written with its coordinates \(x^a\) (its metric
  \(\hat g_{ab}\) and the tensors built from it) a coordinate index of
  the fold, \(a \in \{1, \dots, 7\}\); a repeated index is summed.
  Status: STANDARD.
\item
  \(b\): A frame index (as \(a\)). Status: STANDARD.
\item
  \(H\): The author's constant (an inverse length) that fixes the period
  \(\pi/(6H)\) of \(x_0\), the warp factor
  \(\mathrm{e}^{-2H\lvert\sigma\rvert}\) of the metric, the cosmological
  constant \(\Lambda = -21H^2\) and the tension
  \(\kappa\lambda_{\mathrm{F}} = 12H\) of the fold; in UM-4 a free
  parameter, bounded from below by the propagation of the gauge fields
  in the bulk. Unit: GeV (1/length). Value: \(H \gtrsim 2.5\) TeV (from
  the invisible width of the Z boson). Status: AUTHOR.
\item
  \(\delta\): The Kronecker delta: with two indices \(\delta^M_N = 1\)
  if \(M = N\) and 0 otherwise (likewise \(\delta_{X\mathrm{L}}\) is 1
  if the label \(X\) of a half of dirac16complex is L and 0 otherwise);
  with two index strings of equal length the generalised Kronecker delta
  \(\delta^{M_1M_2\cdots}_{N_1N_2\cdots}\), the determinant of the
  square matrix of the Kronecker deltas \(\delta^{M_a}_{N_b}\) (\(+1\)
  or \(-1\) if the upper indices are an even or odd permutation of the
  lower ones, 0 otherwise). Status: STANDARD.
\item
  \(\lambda_{\mathrm{F}}\): The tension that the static block requires
  of the fold, \(\kappa\lambda_{\mathrm{F}} = 12H\): the total vacuum
  energy of the seven-dimensional fold, made of the tension of the fold
  fluid and the vacuum energy of the fields of the fold (both have the
  pressure minus their energy density along ordinary space and along the
  extra times), per unit volume of the six other directions of the fold
  (eight-dimensional normalisation); in the four-dimensional
  normalisation it is \(\lambda_{\mathrm{F}}V_*\), and the excess of the
  total over \(\lambda_{\mathrm{F}}V_*\) is the dark energy
  \(\rho_\Lambda\). Unit: GeV\(^7\). Value: \(12H/\kappa\);
  \(\lambda_{\mathrm{F}}V_* = 30H^2M_{\mathrm{red}}^2 = 1.1\times10^{45}\)
  GeV\(^4\) for \(H = 2.5\) TeV. Status: DERIVED.
\item
  \(\hat G^\mu{}_\nu\): The Einstein tensor of the metric
  \(\hat g_{ab}\) of the fold (FD-2), with frame indices along its seven
  directions; with four-dimensional indices, \(\hat G^\mu{}_\nu\), its
  components along the four observable directions, which for frozen
  extra times are those of the Einstein tensor of the four-dimensional
  metric \(\hat h_{\mu\nu}\). Unit: 1/length\(^2\). Status: STANDARD.
\item
  \(M_{\mathrm{red}}\): The reduced Planck mass,
  \(M_{\mathrm{red}} = (8\pi G_{\mathrm{N}})^{-1/2}\). Unit: GeV. Value:
  \(2.435\times10^{18}\) GeV. Status: STANDARD.
\item
  \(V_*\): The physical volume of the extra times \(x_5, x_6, x_7\) on
  the fold in the late universe,
  \(V_* = V_{\mathrm{t}}(\mathrm{e}^{-a_4}b_4)^3\) (with
  \(V_{\mathrm{t}} = \int\mathrm{d}x_5\,\mathrm{d}x_6\,\mathrm{d}x_7\)
  the coordinate volume of the extra times) at a time when
  \(\mathrm{e}^{-a_4}b_4\) is constant; it converts the
  eight-dimensional normalisation into the four-dimensional one. Unit:
  length\(^3\) (GeV\(^{-3}\)). Status: ASSUMED.
\item
  \(\hat h_{\mu\nu}\): The four-dimensional metric of the fold,
  \(\hat h_{\mu\nu} = \mathrm{diag}(\mathrm{e}^{2a_4}, \mathrm{e}^{2a_4}, \mathrm{e}^{2a_4}, -1)\)
  in the coordinates \(x_1, \dots, x_4\) (the induced metric at
  \(\sigma = 0\) without the extra times); \(\hat h^{\mu\nu}\) its
  inverse, \(\hat h\) its determinant,
  \(\sqrt{-\hat h} = \mathrm{e}^{3a_4}\); \(\hat h^{(1)}_{\mu\nu}\) its
  first-order cosmological perturbation. Unit: dimensionless. Status:
  DERIVED.
\item
  \(\rho_\Lambda\): The vacuum energy density of the fold above its
  static tension, the dark energy of UM-4:
  \(\rho_\Lambda = \rho_{\mathrm{ff}} + \rho_{\mathrm{vac}} - \lambda_{\mathrm{F}}V_*\)
  (four-dimensional normalisation). Unit: GeV\(^4\). Value:
  \(0.685\,\rho_{\mathrm{c}}^{(0)} = 2.5\times10^{-47}\) GeV\(^4\)
  (Planck 2018). Status: DERIVED.
\item
  \(G_{\mathrm{N}}\): Newton's constant of four-dimensional gravity (the
  gravity of our world, on the fold),
  \(G_{\mathrm{N}} = 1/M_{\mathrm{Pl}}^2\) with the Planck mass
  \(M_{\mathrm{Pl}} = 1.221\times10^{19}\) GeV. Unit: GeV\(^{-2}\).
  Value: \(6.708\times10^{-39}\) GeV\(^{-2}\). Status: STANDARD.
\item
  \(V_{\mathrm{t}}\): The coordinate volume of the (compact) extra
  times,
  \(V_{\mathrm{t}} = \int\mathrm{d}x_5\,\mathrm{d}x_6\,\mathrm{d}x_7\).
  Unit: length\(^3\). Status: ASSUMED.
\item
  \(\rho_{\mathrm{c}}\): The critical density, the energy density of a
  spatially flat fold at the expansion rate \(H_4\);
  \(\rho_{\mathrm{c}}^{(0)}\) its value today. Unit: GeV\(^4\). Value:
  \(\rho_{\mathrm{c}}^{(0)} = 1.0537\times10^{-5}\,h_{100}^2\)
  GeV/cm\(^3\). Status: STANDARD.
\item
  \(\rho_{\mathrm{ff}}\): The energy density of the fold fluid. Unit:
  GeV\(^4\). Value:
  \(\lambda_{\mathrm{F}}V_* + \rho_\Lambda - \rho_{\mathrm{vac}}\):
  together with the vacuum energy \(\rho_{\mathrm{vac}}\) of the fields
  of the fold it makes up the tension \(\lambda_{\mathrm{F}}V_*\) that
  the static block requires plus the dark energy \(\rho_\Lambda\).
  Status: ASSUMED (the fluid), DERIVED (the value).
\item
  \(\rho_{\mathrm{vac}}\): The vacuum energy of the fields of the fold:
  the minimum of their potential (for the Standard Model the minimum of
  the Higgs potential with its loop corrections), a tension of the
  seven-dimensional fold. Unit: GeV\(^4\). Status: STANDARD.
\item
  \(H_4\): The expansion rate (Hubble rate) of the fold, \(H_4 = a_4'\),
  as a function of the time or of the redshift; \(H_4^{(0)}\) its value
  today, the Hubble constant. Unit: GeV (1/time). Value:
  \(H_4^{(0)} = 100\,h_{100}\) km/s/Mpc with \(h_{100} = 0.674\) (Planck
  2018): 67.4 km/s/Mpc \(= 1.44\times10^{-42}\) GeV. Status: STANDARD.
\item
  \(M_{\mathrm{Pl}}\): The Planck mass,
  \(M_{\mathrm{Pl}} = G_{\mathrm{N}}^{-1/2}\). Unit: GeV. Value:
  \(1.221\times10^{19}\) GeV. Status: STANDARD.
\item
  \(h_{100}\): The present expansion rate of the fold in units of 100
  km/s/Mpc. Unit: dimensionless. Value: 0.674 (Planck 2018). Status:
  STANDARD.
\end{itemize}

(\(\hat G^a{}_b\) the Einstein tensor of the metric FD-2 of the fold):
the tension plus matter that obeys the Einstein equations of seven
dimensions, \(\hat G^a{}_b = \kappa_7\mathsf{t}^a{}_b\) with
\(\kappa_7 = 5\kappa H/2\) (\(\mathsf{t}^a{}_b\) the matter of the
fold), which for frozen extra times are the Friedmann equations of
general relativity. The fold is therefore not only a place but a
physical object: a thin layer that carries energy and momentum (a
brane). 4. \emph{The author's coordinate \(x_0\) degenerates on it.} The
metric component \(g_{x_0x_0} = \cot^2(6Hx_0)\) vanishes on the fold and
nowhere else in the period; near the fold the proper distance grows with
the square of the coordinate distance,

\[
\lvert\sigma\rvert = 3H\Big(x_0 - \frac{\pi}{12H}\Big)^2\Big[1 + 6H^2\Big(x_0 - \frac{\pi}{12H}\Big)^2 + \dots\Big] \tag{FD-6}
\] \textbf{Symbols in (FD-6).}

\begin{itemize}
\tightlist
\item
  \(\sigma\): The proper coordinate of the hidden direction:
  \(\mathrm{d}\sigma = \lvert\cot z\rvert\,\mathrm{d}x_0\),
  \(\sigma = \mathrm{sgn}(\pi/2 - z)\ln(\sin z)/(6H)\) on one period
  \(0 < z < \pi\); \(\sigma = 0\) is the fold \(z = \pi/2\),
  \(\sigma \to \mp\infty\) the tips \(z \to 0, \pi\); one period of
  \(x_0\) is the whole line \(-\infty < \sigma < \infty\). Unit: length.
  Status: DERIVED.
\item
  \(H\): The author's constant (an inverse length) that fixes the period
  \(\pi/(6H)\) of \(x_0\), the warp factor
  \(\mathrm{e}^{-2H\lvert\sigma\rvert}\) of the metric, the cosmological
  constant \(\Lambda = -21H^2\) and the tension
  \(\kappa\lambda_{\mathrm{F}} = 12H\) of the fold; in UM-4 a free
  parameter, bounded from below by the propagation of the gauge fields
  in the bulk. Unit: GeV (1/length). Value: \(H \gtrsim 2.5\) TeV (from
  the invisible width of the Z boson). Status: AUTHOR.
\item
  \(x_0\): The hidden space-like coordinate (listed first); the metric
  depends on it through \(z = 6Hx_0\), periodically with the period
  \(\pi/(6H)\). Unit: length. Status: AUTHOR.
\item
  \(\pi\): The number pi, the ratio of the circumference of a circle to
  its diameter. Unit: dimensionless. Value: 3.141592654. Status:
  STANDARD.
\end{itemize}

so the two coordinate values on either side of \(x_0 = \pi/(12H)\) at
the same coordinate distance lie at the same proper distance from the
fold: the coordinate line of \(x_0\) is folded onto the line of proper
distances (a fold map, example 5).

\textbf{What lives on the fold.} In UM-4 the Higgs doublet (with the
author's ad hoc vacuum, a time-dependent choice of its vacuum value),
three generations of quarks and leptons, three right-handed neutrinos
and the fold fluid (its tension and the pressure along the extra times
that keeps their volume fixed) live on the fold: they are fields at
\(\sigma = 0\), functions of \(x_1, \dots, x_4\) (no field depends on
the extra times). The gauge fields of
\(SU(3)_C \times SU(2)_L \times U(1)_Y\) live in the bulk; their
massless modes are bound to the fold, so that the photon and the gluons
of the four-dimensional world are waves of the bulk that the fold holds,
and photons that leave the fold travel through the bulk at the speed of
light. dirac16complex (the author's spinor field of eight dimensions)
and gravity live in the bulk \(\sigma \ne 0\) as well as on the fold;
the states \(\chi\) of dirac16complex that the fold holds (in UM-4
long-lived resonances) are the part of it near the fold, and the portals
(couplings between dirac16complex and the fields of the fold) act at
\(\sigma = 0\). Integrated over the extra times, the fold is the
four-dimensional universe that we observe, with the metric
\(\mathrm{diag}(\mathrm{e}^{2a_4}, \mathrm{e}^{2a_4}, \mathrm{e}^{2a_4}, -1)\)
in \(x_1, \dots, x_4\). In UM-4 files ``the fold'' means this surface;
the examples below use the word for other surfaces with the same
properties. The word is not used in its other senses: a set of equal
energy levels is called degenerate, and a growth by the factor
\(\mathrm{e}\) is called an e-folding.

\textbf{Examples of folds that are meaningful for UM-4.}

\begin{enumerate}
\def\labelenumi{\arabic{enumi}.}
\tightlist
\item
  \emph{Our fold.} The fold \(z = \pi/2\) of the period that contains
  our world (FD-1): the seven-dimensional world-volume of the quarks,
  leptons and the Higgs field, and after the integration over the extra
  times the four-dimensional universe that we observe.
\item
  \emph{The copies of the fold.} Because the metric repeats with the
  period \(\pi/(6H)\), every surface
  \(x_0 = \pi/(12H) + (\text{an integer multiple of } \pi/(6H))\) is a
  copy of the fold, one per period. Neighbouring copies are separated by
  a tip (\(z\) a multiple of \(\pi\)), where the warp factor
  \((\sin^2z)^{1/6}\) vanishes and \(g_{x_0x_0} = \cot^2z\) diverges,
  and which lies at an infinite proper distance from both: in proper
  distance every copy has its own complete line
  \(-\infty < \sigma < \infty\). A compact direction of length \(L\)
  would give every field a tower of particles with the masses
  \(2\pi n/L\), \(n = 1, 2, \dots\) (a Kaluza-Klein tower); because the
  copies are infinitely far apart, the period \(\pi/(6H)\) of \(x_0\)
  produces no such tower with the spacing \(12H\): every field of the
  bulk has instead a continuum of modes that starts at zero mass, and
  the fold holds its massless mode.
\item
  \emph{A folded sheet of paper.} The crease of a folded sheet is a fold
  in this sense: the two halves are mirror images (property 1) joined
  along the crease, the sheet is continuous there and its direction
  changes abruptly (property 3). In UM-4 the direction across the crease
  is the hidden direction \(\sigma\), the seven directions along the
  crease are \(x_1, \dots, x_7\), and the crease is the fold.
\item
  \emph{A folded function.} The warp factor
  \(\mathrm{e}^{-2H\lvert\sigma\rvert}\) is the simplest folded
  function: continuous, mirror symmetric, largest at \(\sigma = 0\),
  with a jump of its derivative,
\end{enumerate}

\[
\frac{\mathrm{d}}{\mathrm{d}\sigma}\mathrm{e}^{-2H\lvert\sigma\rvert} = -2H\,\mathrm{sgn}(\sigma)\,\mathrm{e}^{-2H\lvert\sigma\rvert},\qquad
\frac{\mathrm{d}^2}{\mathrm{d}\sigma^2}\mathrm{e}^{-2H\lvert\sigma\rvert} = 4H^2\,\mathrm{e}^{-2H\lvert\sigma\rvert} - 4H\,\delta_{\mathrm{D}}(\sigma) \tag{FD-7}
\] \textbf{Symbols in (FD-7).}

\begin{itemize}
\tightlist
\item
  \(\mathrm{d}\): The differential: \(\mathrm{d}\) in front of a
  quantity is an infinitesimal change of it (as \(\mathrm{d}\sigma\), a
  small step of \(\sigma\)); the quotient of two such changes is a
  derivative; \(\mathrm{d}\) with an exponent, in front of several
  variables, is their volume element, the product of their
  differentials. Status: STANDARD.
\item
  \(\sigma\): The proper coordinate of the hidden direction:
  \(\mathrm{d}\sigma = \lvert\cot z\rvert\,\mathrm{d}x_0\),
  \(\sigma = \mathrm{sgn}(\pi/2 - z)\ln(\sin z)/(6H)\) on one period
  \(0 < z < \pi\); \(\sigma = 0\) is the fold \(z = \pi/2\),
  \(\sigma \to \mp\infty\) the tips \(z \to 0, \pi\); one period of
  \(x_0\) is the whole line \(-\infty < \sigma < \infty\). Unit: length.
  Status: DERIVED.
\item
  \(\mathrm{e}\): Euler's number, the base of the natural logarithm;
  \(\mathrm{e}^{X}\) is the exponential of \(X\). Unit: dimensionless.
  Value: 2.718281828. Status: STANDARD.
\item
  \(H\): The author's constant (an inverse length) that fixes the period
  \(\pi/(6H)\) of \(x_0\), the warp factor
  \(\mathrm{e}^{-2H\lvert\sigma\rvert}\) of the metric, the cosmological
  constant \(\Lambda = -21H^2\) and the tension
  \(\kappa\lambda_{\mathrm{F}} = 12H\) of the fold; in UM-4 a free
  parameter, bounded from below by the propagation of the gauge fields
  in the bulk. Unit: GeV (1/length). Value: \(H \gtrsim 2.5\) TeV (from
  the invisible width of the Z boson). Status: AUTHOR.
\item
  \(\mathrm{sgn}\): The sign function: \(+1\) for a positive argument
  and \(-1\) for a negative one. Status: STANDARD.
\item
  \(\delta_{\mathrm{D}}\): The Dirac delta function of \(\sigma\): zero
  for \(\sigma \ne 0\), with the integral 1 over any range that contains
  \(\sigma = 0\); the integral of \(\delta_{\mathrm{D}}(\sigma)\) times
  a smooth function of \(\sigma\) is the value of that function at
  \(\sigma = 0\). Unit: 1/length (for \(\delta_{\mathrm{D}}(\sigma)\)).
  Status: STANDARD.
\end{itemize}

whose delta function is the origin of the thin source of property 3.
Every metric function of FD-4 depends on \(\sigma\) through
\(\lvert\sigma\rvert\) and is folded at \(\sigma = 0\) in this sense. 5.
\emph{A fold map.} In the mathematics of smooth maps the fold (Whitney's
fold singularity) is the map \(u \mapsto u^2\) of a line onto a
half-line: it folds the line at \(u = 0\), two points \(\pm u\) land on
the same point, and the derivative of the map vanishes on the fold. Near
the fold the map from the author's coordinate \(x_0\) to the distance
\(\lvert\sigma\rvert\) from the fold is of this kind (FD-6), and
\(g_{x_0x_0} = \cot^2z\), the square of the derivative
\(\mathrm{d}\sigma/\mathrm{d}x_0\), vanishes there. 6. \emph{The brane
of Randall and Sundrum.} The five-dimensional metric

\[
\mathrm{d}y_{\mathrm{RS}}^2 + \mathrm{e}^{-2k_{\mathrm{RS}}\lvert y_{\mathrm{RS}}\rvert}\big(\mathrm{d}x_1^2 + \mathrm{d}x_2^2 + \mathrm{d}x_3^2
 - \mathrm{d}x_4^2\big),\qquad \lambda_{\mathrm{RS}} = \frac{6k_{\mathrm{RS}}}{\kappa_5},\qquad \Lambda_{\mathrm{RS}} = -\frac{6k_{\mathrm{RS}}^2}{\kappa_5} \tag{FD-8}
\] \textbf{Symbols in (FD-8).}

\begin{itemize}
\tightlist
\item
  \(\mathrm{d}\): The differential: \(\mathrm{d}\) in front of a
  quantity is an infinitesimal change of it (as \(\mathrm{d}\sigma\), a
  small step of \(\sigma\)); the quotient of two such changes is a
  derivative; \(\mathrm{d}\) with an exponent, in front of several
  variables, is their volume element, the product of their
  differentials. Status: STANDARD.
\item
  \(y_{\mathrm{RS}}\): The proper coordinate across the brane of the
  five-dimensional model of Randall and Sundrum (FD-8, an example, not
  UM-4). Unit: length. Status: STANDARD.
\item
  \(\mathrm{e}\): Euler's number, the base of the natural logarithm;
  \(\mathrm{e}^{X}\) is the exponential of \(X\). Unit: dimensionless.
  Value: 2.718281828. Status: STANDARD.
\item
  \(k_{\mathrm{RS}}\): The inverse curvature length of the bulk of the
  model of Randall and Sundrum: the warp factor of the brane is
  \(\mathrm{e}^{-2k_{\mathrm{RS}}\lvert y_{\mathrm{RS}}\rvert}\) (FD-8,
  an example, not UM-4). Unit: GeV (1/length). Status: STANDARD.
\item
  \(x_1\): The first coordinate of ordinary 3-space. Unit: length.
  Status: AUTHOR.
\item
  \(x_2\): The second coordinate of ordinary 3-space. Unit: length.
  Status: AUTHOR.
\item
  \(x_3\): The third coordinate of ordinary 3-space. Unit: length.
  Status: AUTHOR.
\item
  \(x_4\): The time; a prime on a function of \(x_4\) alone is
  \(\mathrm{d}/\mathrm{d}x_4\). Unit: length (= time, \(c = 1\)).
  Status: AUTHOR.
\item
  \(\lambda_{\mathrm{RS}}\): The tension (energy per unit volume of
  ordinary space) of the brane of the model of Randall and Sundrum,
  \(\lambda_{\mathrm{RS}} = 6k_{\mathrm{RS}}/\kappa_5\) (FD-8, an
  example, not UM-4). Unit: GeV\(^4\). Status: STANDARD.
\item
  \(\kappa_5\): The gravitational coupling of five dimensions in the
  model of Randall and Sundrum, \(8\pi\) times the Newton constant of
  five dimensions (FD-8, an example, not UM-4). Unit: GeV\(^{-3}\).
  Status: STANDARD.
\item
  \(\Lambda_{\mathrm{RS}}\): The vacuum energy density of the bulk of
  the model of Randall and Sundrum,
  \(\Lambda_{\mathrm{RS}} = -6k_{\mathrm{RS}}^2/\kappa_5\) (FD-8, an
  example, not UM-4). Unit: GeV\(^5\). Status: STANDARD.
\end{itemize}

has a fold at \(y_{\mathrm{RS}} = 0\) with the properties 1 to 3 (its
coordinate \(y_{\mathrm{RS}}\) is already the proper distance, so
property 4 does not arise); Einstein's equations require there a brane
of the tension \(\lambda_{\mathrm{RS}}\) (Israel's junction condition)
and in the bulk the vacuum energy \(\Lambda_{\mathrm{RS}}\), and the
bulk is otherwise empty. The fold of UM-4 is of the same kind, with
every direction of the fold warped by the same factor as in the model of
Randall and Sundrum: its static block is a space of constant curvature
\(-H^2\) in all eight directions, a vacuum of Einstein gravity with
\(\Lambda = -21H^2\), with a fold of pure tension (FD-12), and the
graviton is bound to the fold as there. It differs in that the fold has
seven dimensions, three of them extra times, in that the space is not
cut at the brane (Randall and Sundrum identify the two sides of their
brane, one physical side; the fold of UM-4 has two distinct sides,
example 7), and in that the gauge fields live in the bulk as well. 7.
\emph{A particle held by a thin layer.} A particle held by an attractive
thin layer has a wave function that falls off exponentially with the
distance on both sides, and the layer makes its derivative jump (for the
first-order Dirac equation, the wave function itself). In UM-4 the
states \(\chi\) of dirac16complex near the fold are of this kind only
near the fold: at rest and without their self-interaction (the field is
\(\Psi = \mathrm{e}^{(7/2)H\lvert\sigma\rvert}\chi(\sigma)\mathrm{e}^{-\mathrm{i}\omega x_4}\))

\[
\begin{aligned}
&\chi(\sigma) = \mathrm{e}^{-q_{\mathrm{b}}\lvert\sigma\rvert}\chi(0^+)\ \ (0 < \sigma \ll \sigma_{\mathrm{t}}),\qquad
\chi(\sigma) = \mathrm{e}^{-q_{\mathrm{b}}\lvert\sigma\rvert}\chi(0^-)\ \ (0 < -\sigma \ll \sigma_{\mathrm{t}}),\\
&\chi(0^+) = \exp\Big(\gamma^{(0)}\big(\mu_{\mathrm{L}}P_{\mathrm{L}} + \mu_{\mathrm{R}}P_{\mathrm{R}}\big)\Big)\chi(0^-),\qquad
\omega_0 = \omega_{\mathrm{r}} - \frac{\mathrm{i}}{2}\,\Gamma_\chi,\qquad \omega_{\mathrm{r}} \to \omega_{\mathrm{b}} = \sqrt{m_\Psi^2 - q_{\mathrm{b}}^2}
\end{aligned} \tag{FD-9}
\] \textbf{Symbols in (FD-9).}

\begin{itemize}
\tightlist
\item
  \(\chi(\sigma)\): The profile of a mode of dirac16complex along the
  hidden direction,
  \(\Psi = \mathrm{e}^{\frac72H\lvert\sigma\rvert}\chi(\sigma)\mathrm{e}^{\mathrm{i}(\vec k\cdot\vec x - \omega x_4)}\)
  (a 16-component spinor function of \(\sigma\)); as a four-dimensional
  field, \(\chi\) denotes the particle of the level of dirac16complex at
  the fold (a long-lived resonance) and \(\bar\chi\) its antiparticle.
  Unit: GeV\(^{2}\) (with the \(V_*\) normalisation). Status: DERIVED.
\item
  \(\sigma\): The proper coordinate of the hidden direction:
  \(\mathrm{d}\sigma = \lvert\cot z\rvert\,\mathrm{d}x_0\),
  \(\sigma = \mathrm{sgn}(\pi/2 - z)\ln(\sin z)/(6H)\) on one period
  \(0 < z < \pi\); \(\sigma = 0\) is the fold \(z = \pi/2\),
  \(\sigma \to \mp\infty\) the tips \(z \to 0, \pi\); one period of
  \(x_0\) is the whole line \(-\infty < \sigma < \infty\). Unit: length.
  Status: DERIVED.
\item
  \(\mathrm{e}\): Euler's number, the base of the natural logarithm;
  \(\mathrm{e}^{X}\) is the exponential of \(X\). Unit: dimensionless.
  Value: 2.718281828. Status: STANDARD.
\item
  \(q_{\mathrm{b}}\): The inverse width of a state of dirac16complex at
  the fold: inside the turning point its profile falls as
  \(\mathrm{e}^{-q_{\mathrm{b}}\lvert\sigma\rvert}\). Unit: GeV. Status:
  DERIVED.
\item
  \(\sigma_{\mathrm{t}}\): The turning point of a state of
  dirac16complex of the frequency \(\omega\): the depth at which the
  warped frequency equals the bulk mass,
  \(\lvert\omega\rvert\mathrm{e}^{H\sigma_{\mathrm{t}}} = m_\Psi\);
  beyond it every solution oscillates. Unit: length. Status: DERIVED.
\item
  \(\gamma^M,\ \gamma^{(a)}\): The Dirac matrices of dirac16complex
  along the coordinates, \(\gamma^M = \mathbf{e}_a{}^MT^a\),
  \(\gamma_M = g_{MN}\gamma^N\); \(\gamma^{(a)} = T^a\) along the frame
  direction \((a)\) (\(\gamma^{(0)}\) along \(\sigma\), \(\gamma^{(4)}\)
  along the time). Unit: dimensionless. Status: AUTHOR.
\item
  \(\mu_X,\ \mu_{\mathrm{L}},\ \mu_{\mathrm{R}}\): The dimensionless
  strengths of the fold for the left-handed and right-handed halves of
  dirac16complex, set by the mean fields of the Standard Model at the
  fold: the jump of the profile across the fold is the exponential of
  \(\gamma^{(0)}(\mu_{\mathrm{L}}P_{\mathrm{L}} + \mu_{\mathrm{R}}P_{\mathrm{R}})\);
  negative values attract. Unit: dimensionless. Status: DERIVED.
\item
  \(P_{\mathrm{L}}\): The author's left-handed chiral projector,
  \(P_{\mathrm{L}} = (\mathbf{1} - \Gamma_8)/2\). Unit: dimensionless.
  Status: AUTHOR.
\item
  \(P_{\mathrm{R}}\): The author's right-handed chiral projector,
  \(P_{\mathrm{R}} = (\mathbf{1} + \Gamma_8)/2\). Unit: dimensionless.
  Status: AUTHOR.
\item
  \(\omega_0\): The energy at rest (zero momentum along ordinary space
  and along the extra times) of a state of dirac16complex at the fold;
  complex for a resonance,
  \(\omega_0 = \omega_{\mathrm{r}} - \frac{\mathrm{i}}{2}\Gamma_\chi\).
  Unit: GeV. Status: DERIVED.
\item
  \(\omega_{\mathrm{r}}\): The mass of the particle \(\chi\): the real
  part of the energy at rest of the resonance of dirac16complex at the
  fold,
  \(\omega_{\mathrm{r}} = \omega_{\mathrm{b}}(1 - H/(2q_{\mathrm{b}})) + \dots\).
  Unit: GeV. Status: DERIVED.
\item
  \(\mathrm{i}\): The imaginary unit, \(\mathrm{i}^2 = -1\). Unit:
  dimensionless. Status: STANDARD.
\item
  \(\Gamma_\chi\): The width of the particle \(\chi\) at rest: the rate
  of its decay into the waves of dirac16complex in the bulk. Unit: GeV.
  Value: below \(1/t_0 = 1.5\times10^{-42}\) GeV for dark matter.
  Status: DERIVED.
\item
  \(\omega_{\mathrm{b}}\): The energy at zero momentum of the state of
  dirac16complex at the fold in the limit \(H \to 0\) (the bound state
  of UM-3): the particle \(\chi\) of UM-4 has the mass
  \(\omega_{\mathrm{r}} = \omega_{\mathrm{b}}(1 - H/(2q_{\mathrm{b}})) + \dots\),
  equal to \(\omega_{\mathrm{b}}\) up to the relative correction
  \(H/(2q_{\mathrm{b}})\). Unit: GeV. Status: DERIVED.
\item
  \(m_\Psi\): The mass parameter of dirac16complex in the bulk. Unit:
  GeV. Status: RECORD.
\item
  \(\Gamma_8\): The chirality matrix of eight dimensions,
  \(\Gamma_8 = T_0T_1\cdots T_7\) (the author's T16\^{}A{[}8{]})
  \(= \mathrm{diag}(-\mathbf{1}_8, \mathbf{1}_8)\);
  \(\Gamma_8^2 = \mathbf{1}\); it anticommutes with every \(T_a\). Unit:
  dimensionless. Status: AUTHOR.
\item
  \(\Psi,\ \bar\Psi\): dirac16complex: the author's 16-component complex
  spinor field of eight dimensions;
  \(\bar\Psi = \Psi^\dagger\sigma_{16}\) is its Dirac adjoint; in a term
  of the fold, \(\Psi\) is its value at \(\sigma = 0\). Unit:
  GeV\(^{7/2}\). Status: AUTHOR.
\item
  \(\mathbf{1},\ \mathbf{1}_8\): The identity matrix of the size fixed
  by the context (\(16 \times 16\) on dirac16complex, \(4 \times 4\) on
  a four-dimensional spinor, \(3 \times 3\) in generation or flavour
  space); \(\mathbf{1}_8\) is the \(8 \times 8\) identity matrix. Unit:
  dimensionless. Status: STANDARD.
\item
  \(\mathbf{e}_a{}^M,\ \mathbf{e}^a{}_M\): The vielbein (frame field) of
  the metric: \(g_{MN} = \eta_{ab}\,\mathbf{e}^a{}_M\mathbf{e}^b{}_N\);
  \(\mathbf{e}_a{}^M\) is its inverse;
  \(\mathbf{e}_{bN} = \eta_{bc}\mathbf{e}^c{}_N\); \(\mathbf{e}^{(a)}\)
  is the frame one-form \(\mathbf{e}^a{}_M\,\mathrm{d}X^M\). Unit:
  dimensionless. Status: STANDARD.
\item
  \(\omega\): A frequency: the energy of a mode
  \(\propto \mathrm{e}^{-\mathrm{i}\omega x_4}\) (\(\omega_o\) of the
  orbital \(o\)). Unit: GeV. Status: STANDARD.
\item
  \(\vec k\): The (comoving) momentum along ordinary space, \(\vec k^2\)
  its square. Unit: GeV. Status: STANDARD.
\item
  \(\vec x\): The position in ordinary space,
  \(\vec x = (x_1, x_2, x_3)\). Unit: length. Status: AUTHOR.
\item
  \(t_0\): The age of the universe. Unit: time. Value: 13.8 Gyr. Status:
  STANDARD.
\item
  \(\sigma_{16}\): The matrix of the Dirac adjoint,
  \(\sigma_{16} = T_0T_1T_2T_3\), real and symmetric; \(\sigma_{16}T_a\)
  is real and antisymmetric. Unit: dimensionless. Status: AUTHOR.
\item
  \({}^\dagger\): Hermitian conjugation: the complex conjugate of the
  transpose of a matrix or spinor; for an operator on the space of
  states, its adjoint with respect to the scalar product of that space
  (for the quantised dirac16complex the indefinite Krein scalar
  product). Status: STANDARD.
\item
  \(o\): The label of a single-particle state of dirac16complex: a
  Kohn-Sham orbital, or a mode of the quantised field; a sum over \(o\)
  runs over the orbitals or the modes. Status: STANDARD.
\end{itemize}

where \(\mu_{\mathrm{L}}\), \(\mu_{\mathrm{R}}\) are the strengths of
the fold for the left- and right-handed halves of dirac16complex (set by
the fields of the fold),
\(q_{\mathrm{b}} = -m_\Psi(\mu_{\mathrm{L}} + \mu_{\mathrm{R}})\tanh\varsigma/(2\varsigma)\),
\(\varsigma = \sqrt{\mu_{\mathrm{L}}\mu_{\mathrm{R}}}\)
(\(q_{\mathrm{b}} = -m_\Psi\mu_{\mathrm{L}}/2\) for
\(\mu_{\mathrm{R}} = 0\)), and \(\omega_{\mathrm{r}}\) tends to the
energy \(\omega_{\mathrm{b}}\) of the bound state of UM-3 when
\(H/q_{\mathrm{b}} \to 0\). Because the time is warped, the energy
\(\omega\) of the field enters its equation multiplied by
\(\mathrm{e}^{H\lvert\sigma\rvert}\): beyond the turning point
\(\sigma_{\mathrm{t}}\)
(\(\omega_{\mathrm{r}}\mathrm{e}^{H\sigma_{\mathrm{t}}} = m_\Psi\)) the
profile oscillates as a wave that leaves the fold, so the state is a
resonance of the width \(\Gamma_\chi\), which falls like
\(\mathrm{e}^{-2F(\upsilon_{\mathrm{b}})m_\Psi/H}\) with
\(F(\upsilon) = \ln((1 + \sqrt{1 - \upsilon^2})/\upsilon) - \sqrt{1 - \upsilon^2}\)
and \(\upsilon_{\mathrm{b}} = \omega_{\mathrm{b}}/m_\Psi\) (the
tunnelling through the region between the fold and
\(\sigma_{\mathrm{t}}\)). The profile is not a function of
\(\lvert\sigma\rvert\): it jumps at the fold,
\(\chi(0^+) \ne \chi(0^-)\). 8. \emph{The fold as a mirror for
dirac16complex.} The reflection across the fold acts on the profiles
\(\chi(\sigma)\) in two ways,

\[
(\mathsf{R}_1\chi)(\sigma) = T_0T_1T_2T_3\,\chi(-\sigma),\qquad (\mathsf{R}_2\chi)(\sigma) = T_0T_5T_6T_7\,\chi(-\sigma) \tag{FD-10}
\] \textbf{Symbols in (FD-10).}

\begin{itemize}
\tightlist
\item
  \(\mathsf{R}_1\): The reflection across the fold acting on a profile
  of dirac16complex together with the reflection of ordinary space,
  \((\mathsf{R}_1\chi)(\sigma) = T_0T_1T_2T_3\,\chi(-\sigma)\) (FD-10):
  a symmetry of dirac16complex, of gravity and of the mean field of the
  gas of the particles \(\chi\); the states at the fold of Krein charge
  \(-1\) are even, those of Krein charge \(+1\) odd. Unit:
  dimensionless. Status: DERIVED.
\item
  \(\mathsf{R}_2\): The reflection across the fold acting on a profile
  of dirac16complex together with the reflection of the extra times,
  \((\mathsf{R}_2\chi)(\sigma) = T_0T_5T_6T_7\,\chi(-\sigma)\) (FD-10):
  it maps the states at the fold of one Krein charge onto those of the
  other. Unit: dimensionless. Status: DERIVED.
\item
  \(\chi(\sigma)\): The profile of a mode of dirac16complex along the
  hidden direction,
  \(\Psi = \mathrm{e}^{\frac72H\lvert\sigma\rvert}\chi(\sigma)\mathrm{e}^{\mathrm{i}(\vec k\cdot\vec x - \omega x_4)}\)
  (a 16-component spinor function of \(\sigma\)); as a four-dimensional
  field, \(\chi\) denotes the particle of the level of dirac16complex at
  the fold (a long-lived resonance) and \(\bar\chi\) its antiparticle.
  Unit: GeV\(^{2}\) (with the \(V_*\) normalisation). Status: DERIVED.
\item
  \(\sigma\): The proper coordinate of the hidden direction:
  \(\mathrm{d}\sigma = \lvert\cot z\rvert\,\mathrm{d}x_0\),
  \(\sigma = \mathrm{sgn}(\pi/2 - z)\ln(\sin z)/(6H)\) on one period
  \(0 < z < \pi\); \(\sigma = 0\) is the fold \(z = \pi/2\),
  \(\sigma \to \mp\infty\) the tips \(z \to 0, \pi\); one period of
  \(x_0\) is the whole line \(-\infty < \sigma < \infty\). Unit: length.
  Status: DERIVED.
\item
  \(T_a,\ T^a\): The author's eight real \(16 \times 16\) Dirac matrices
  T16\^{}A{[}a{]} of 4+4.nb (\(a = 0, \dots, 7\); \(T_0, \dots, T_7\)),
  with \(\{T_a, T_b\} = 2\eta_{ab}\mathbf{1}\); \(T^a = \eta^{ab}T_b\).
  Unit: dimensionless. Status: AUTHOR.
\end{itemize}

(\(T_a\) the author's Dirac matrices; for states at rest
\(\mathsf{R}_1\) is the reflection across the fold combined with the
reflection of ordinary space, \(\mathsf{R}_2\) the reflection across the
fold combined with the reflection of the extra times). The Krein charge
of a state is its share of the conserved charge of dirac16complex, which
is indefinite: a state can carry the charge \(-1\) or \(+1\).
\(\mathsf{R}_1\) is a symmetry of dirac16complex, of gravity and of a
mean field of a gas of the particles \(\chi\) that is the same on both
sides of the fold: the four resonances of Krein charge \(-1\) (the
particles \(\chi\)) are even under it and the four of Krein charge
\(+1\) are odd, so such a mean field cannot mix the two signs of the
indefinite charge. \(\mathsf{R}_2\) maps the states of one Krein charge
onto those of the other and makes the two sets degenerate without the
self-interaction. 9. \emph{The planes of Horava and Witten.} In M-theory
(the eleven-dimensional theory whose limits include the superstring
theories) on an interval (a circle whose points are identified with
their mirror images), the two ten-dimensional planes left in place by
the reflection are mirror surfaces (property 1); each carries one
\(E_8\) gauge group while gravity fills the bulk. As in UM-4, gravity
fills the bulk; unlike UM-4, the gauge fields live on the mirror planes
(in UM-4 they live in the bulk), and the two sides of each plane are
identified (one physical side).

\textbf{What is not a fold.}

\begin{itemize}
\tightlist
\item
  \emph{The tips} \(z = 0\) and \(z = \pi\) (mod \(\pi\)): the mirror
  exchanges them within a period, but there the warp factor vanishes,
  \(g_{x_0x_0}\) diverges, and they lie at an infinite proper distance;
  they carry no thin source.
\item
  \emph{A smooth maximum of a warp factor}, such as
  \(1/\cosh^2(H\sigma)\): mirror symmetric and largest at
  \(\sigma = 0\), but without a kink and without a thin source.
\end{itemize}

\hypertarget{how-to-read-this-document}{%
\section{How to read this document}\label{how-to-read-this-document}}

\hypertarget{what-this-document-contains}{%
\subsection{What this document
contains}\label{what-this-document-contains}}

This document defines the universe model UM-4, ordered by the author on
2026-10-25 at 13:15 PDT: in UM-3 anything that propagated in the bulk
outran light on the fold, because the time \(x_4\) of the author's
metric was not warped while ordinary space and the extra times were
(bulk scalars moved 1.22 times, bulk photons 1.41 times as fast as light
on the fold). The author ordered ``a new metric model in which x4 is
warped like the other directions'', in which ``both photons and
gravitons actually propagate at c in the bulk'', photon waves are
transverse spin-1 waves and gravitational waves are spin-2 waves with
the ``plus'' and ``cross'' polarisations, and he ordered that UM-3 be
kept and UM-4 placed in new folders. This document gives the spacetime
and its metric, the fields and where they live, the complete action from
which every field equation follows, and every parameter with its value
and its origin. It is self-contained: every equation is written out (the
terms of the action carry the tags IN-1 to IN-16 and the metric the tag
FE-27, the same tags as in every UM-4 document), and every symbol of
every display equation is defined right after it; a symbol that one of
these definitions uses is defined in this document as well.

\hypertarget{the-status-labels}{%
\subsection{The status labels}\label{the-status-labels}}

Every equation and every parameter carries one of the labels
\textbf{AUTHOR} (from the author: his notebooks \texttt{4+4.nb} and
\texttt{Dirac\_Einstein-Lovelock\_4+4\_Nash.nb} and his orders of
2026-10-24 and 2026-10-25, the warped time among them), \textbf{RECORD}
(from the record of this repository: the Lagrangian of dirac16complex
with its canonical spin connection and its contact self-interaction, the
Kohn-Sham treatment), \textbf{STANDARD} (textbook physics or a published
measurement), \textbf{ASSUMED} (a choice of UM-4 without derivation),
\textbf{INVENTED} (a new interaction: the portals of dirac16complex),
\textbf{DERIVED} (computed and verified in UM-4), \textbf{ESTIMATE} (an
order of magnitude) and \textbf{OPEN} (not settled).

\hypertarget{coordinates-signature-and-units}{%
\subsection{Coordinates, signature and
units}\label{coordinates-signature-and-units}}

The author's coordinates are \(X = \{x_0, x_1, \dots, x_7\}\), always
with the hidden direction first: \(x_0\) is the hidden space-like
coordinate, \(x_1, x_2, x_3\) ordinary space, \(x_4\) the time, and
\(x_5, x_6, x_7\) three extra time-like coordinates. The signature is
\((+,+,+,+,-,-,-,-)\) in this order. Units are natural: the reduced
Planck constant \(\hbar\) and the speed of light \(c\) are 1. A prime on
a function of the time \(x_4\) alone is its derivative with respect to
\(x_4\).

\hypertarget{the-spacetime}{%
\section{The spacetime}\label{the-spacetime}}

\[
\begin{aligned}
&g_{MN} = \mathrm{diag}\Big(\cot^2z,\ \mathrm{e}^{2a_4}(\sin^2z)^{1/6},\ \mathrm{e}^{2a_4}(\sin^2z)^{1/6},
 \ \mathrm{e}^{2a_4}(\sin^2z)^{1/6},\ -(\sin^2z)^{1/6},\\
&\qquad\qquad -\mathrm{e}^{-2a_4}b_4^2(\sin^2z)^{1/6},\ -\mathrm{e}^{-2a_4}b_4^2(\sin^2z)^{1/6},
 \ -\mathrm{e}^{-2a_4}b_4^2(\sin^2z)^{1/6}\Big),\qquad z = 6Hx_0,\\
&\sigma = \mathrm{sgn}\Big(\frac{\pi}{2} - z\Big)\frac{\ln\sin z}{6H},\qquad \mathrm{d}\sigma = \lvert\cot z\rvert\,\mathrm{d}x_0,\\
&g_{MN}\,\mathrm{d}X^M\mathrm{d}X^N = \mathrm{d}\sigma^2 + \mathrm{e}^{-2H\lvert\sigma\rvert}\Big[\mathrm{e}^{2a_4}\big(\mathrm{d}x_1^2
 + \mathrm{d}x_2^2 + \mathrm{d}x_3^2\big) - \mathrm{d}x_4^2 - \mathrm{e}^{-2a_4}b_4^2\big(\mathrm{d}x_5^2 + \mathrm{d}x_6^2 + \mathrm{d}x_7^2\big)\Big]
\end{aligned} \tag{FE-27}
\] \textbf{Symbols in (FE-27).}

\begin{itemize}
\tightlist
\item
  \(g_{MN}\): The metric tensor of the eight-dimensional spacetime
  (signature \(+,+,+,+,-,-,-,-\) in the order \(x_0, \dots, x_7\)):
  \(g_{MN}\,\mathrm{d}X^M\mathrm{d}X^N\) is the squared interval,
  \(g^{MN}\) its inverse, \(\lvert g\rvert\) the absolute value of its
  determinant (\(\sqrt{\lvert g\rvert}\,\mathrm{d}^8X\) is the proper
  volume element). Unit: dimensionless. Status: AUTHOR.
\item
  \(M\): A coordinate index of eight dimensions (also with a numbered
  copy \(M_1, M_2, \dots\)): it runs over \(x_0, x_1, \dots, x_7\); a
  repeated index is summed. Status: STANDARD.
\item
  \(N\): A coordinate index of eight dimensions (also
  \(N_1, N_2, \dots\)), running over \(x_0, \dots, x_7\); a repeated
  index is summed. Status: STANDARD.
\item
  \(\mathrm{diag}\): A diagonal matrix with the listed entries on its
  diagonal (in the order of the coordinates \(x_0, x_1, \dots, x_7\) for
  a tensor of eight dimensions, \(x_1, \dots, x_7\) for a tensor of the
  fold, \(x_1, \dots, x_4\) for a four-dimensional tensor). Status:
  STANDARD.
\item
  \(z\): The angle variable of the hidden direction, \(z = 6Hx_0\); one
  period of \(x_0\) is \(0 < z < \pi\), the fold is \(z = \pi/2\). Unit:
  dimensionless. Status: AUTHOR.
\item
  \(\mathrm{e}\): Euler's number, the base of the natural logarithm;
  \(\mathrm{e}^{X}\) is the exponential of \(X\). Unit: dimensionless.
  Value: 2.718281828. Status: STANDARD.
\item
  \(a_4\): The author's metric function of the time \(x_4\) that scales
  ordinary 3-space: on the fold \(\mathrm{e}^{a_4}\) is the scale factor
  of cosmology and \(a_4'\) the Hubble rate. Unit: dimensionless.
  Status: AUTHOR.
\item
  \(b_4\): The author's metric function of the time \(x_4\) in the
  extra-time components, \(b_4 > 0\), never identically 1 (the author's
  order of 2026-10-24); the scale factor (the factor of the lengths) of
  the extra times is \(\mathrm{e}^{-a_4}b_4\) on the fold and
  \(\mathrm{e}^{-a_4}b_4\,\mathrm{e}^{-H\lvert\sigma\rvert}\) at the
  distance \(\sigma\) from it (the square root of the warp factor times
  \(\mathrm{e}^{-a_4}b_4\)). Unit: dimensionless. Status: AUTHOR.
\item
  \(H\): The author's constant (an inverse length) that fixes the period
  \(\pi/(6H)\) of \(x_0\), the warp factor
  \(\mathrm{e}^{-2H\lvert\sigma\rvert}\) of the metric, the cosmological
  constant \(\Lambda = -21H^2\) and the tension
  \(\kappa\lambda_{\mathrm{F}} = 12H\) of the fold; in UM-4 a free
  parameter, bounded from below by the propagation of the gauge fields
  in the bulk. Unit: GeV (1/length). Value: \(H \gtrsim 2.5\) TeV (from
  the invisible width of the Z boson). Status: AUTHOR.
\item
  \(x_0\): The hidden space-like coordinate (listed first); the metric
  depends on it through \(z = 6Hx_0\), periodically with the period
  \(\pi/(6H)\). Unit: length. Status: AUTHOR.
\item
  \(\sigma\): The proper coordinate of the hidden direction:
  \(\mathrm{d}\sigma = \lvert\cot z\rvert\,\mathrm{d}x_0\),
  \(\sigma = \mathrm{sgn}(\pi/2 - z)\ln(\sin z)/(6H)\) on one period
  \(0 < z < \pi\); \(\sigma = 0\) is the fold \(z = \pi/2\),
  \(\sigma \to \mp\infty\) the tips \(z \to 0, \pi\); one period of
  \(x_0\) is the whole line \(-\infty < \sigma < \infty\). Unit: length.
  Status: DERIVED.
\item
  \(\mathrm{sgn}\): The sign function: \(+1\) for a positive argument
  and \(-1\) for a negative one. Status: STANDARD.
\item
  \(\pi\): The number pi, the ratio of the circumference of a circle to
  its diameter. Unit: dimensionless. Value: 3.141592654. Status:
  STANDARD.
\item
  \(\mathrm{d}\): The differential: \(\mathrm{d}\) in front of a
  quantity is an infinitesimal change of it (as \(\mathrm{d}\sigma\), a
  small step of \(\sigma\)); the quotient of two such changes is a
  derivative; \(\mathrm{d}\) with an exponent, in front of several
  variables, is their volume element, the product of their
  differentials. Status: STANDARD.
\item
  \(X = \{x_0, \dots, x_7\}\): The author's coordinates of the
  eight-dimensional spacetime, \(X = \{x_0, x_1, \dots, x_7\}\) with the
  hidden direction first (Dirac\_Einstein-Lovelock\_4+4\_Nash.nb,
  section DEFCHART); \(X^M\) is the coordinate with the index \(M\) and
  \(\mathrm{d}^8X = \mathrm{d}x_0\,\mathrm{d}x_1\cdots\mathrm{d}x_7\)
  the coordinate volume element. Unit: length (each coordinate). Status:
  AUTHOR.
\item
  \(x_1\): The first coordinate of ordinary 3-space. Unit: length.
  Status: AUTHOR.
\item
  \(x_2\): The second coordinate of ordinary 3-space. Unit: length.
  Status: AUTHOR.
\item
  \(x_3\): The third coordinate of ordinary 3-space. Unit: length.
  Status: AUTHOR.
\item
  \(x_4\): The time; a prime on a function of \(x_4\) alone is
  \(\mathrm{d}/\mathrm{d}x_4\). Unit: length (= time, \(c = 1\)).
  Status: AUTHOR.
\item
  \(x_5\): The first of the three extra time-like coordinates. Unit:
  length. Status: AUTHOR.
\item
  \(x_6\): The second of the three extra time-like coordinates. Unit:
  length. Status: AUTHOR.
\item
  \(x_7\): The third of the three extra time-like coordinates. Unit:
  length. Status: AUTHOR.
\end{itemize}

The metric of UM-4: the author's metric with his two functions of the
time \(a_4(x_4)\) and \(b_4(x_4)\) (\(b_4\) is never identically 1, by
the author's order of 2026-10-24) and his constant \(H\), with the one
change that his order of 2026-10-25 13:15 requires: the time \(x_4\)
carries the warp factor \((\sin^2z)^{1/6}\) of the other six directions
of the fold (\(g_{x_4x_4} = -(\sin^2z)^{1/6}\); in UM-3 it was \(-1\)).
The metric depends on \(x_0\) only through \(z = 6Hx_0\); one period
\(0 < z < \pi\) of \(x_0\) is, in the proper coordinate \(\sigma\) of
the hidden direction, the whole line \(-\infty < \sigma < \infty\)
(DERIVED). In that coordinate all seven directions of the fold, the time
included, are warped by \(\mathrm{e}^{-H\lvert\sigma\rvert}\), which has
a kink at \(\sigma = 0\) (\(z = \pi/2\)): this seven-dimensional surface
is the fold, and the rest of spacetime is the bulk. Every slice
\(\sigma = \text{const}\) is then a copy of the fold scaled by the same
factor in all seven directions, so that it has the light cone of the
fold: whatever moves at the local speed of light in the bulk moves along
ordinary space at the speed of light of the fold (DERIVED,
\texttt{um4\_geometry.py}, G8). This is the smallest change of the
author's metric that meets the order: it keeps the coordinates, the
period, the fold, the functions \(a_4\), \(b_4\) and the warp of
ordinary space and of the extra times. Without the fold and the
functions \(a_4\), \(b_4\) the static block (\(a_4\), \(b_4\) constant)
is, on each side of the fold, a region of the anti-de Sitter space of
eight dimensions of radius \(1/H\) (constant curvature \(-H^2\);
DERIVED, \texttt{um4\_geometry.py}, G3). Status: AUTHOR (the metric and
the order), ASSUMED (the way the order is met), DERIVED (the chart and
the light cone).

\hypertarget{the-fields-and-where-they-live}{%
\section{The fields and where they
live}\label{the-fields-and-where-they-live}}

\begin{longtable}[]{@{}
  >{\raggedright\arraybackslash}p{(\columnwidth - 8\tabcolsep) * \real{0.1930}}
  >{\raggedright\arraybackslash}p{(\columnwidth - 8\tabcolsep) * \real{0.1404}}
  >{\raggedright\arraybackslash}p{(\columnwidth - 8\tabcolsep) * \real{0.3158}}
  >{\raggedright\arraybackslash}p{(\columnwidth - 8\tabcolsep) * \real{0.0877}}
  >{\raggedright\arraybackslash}p{(\columnwidth - 8\tabcolsep) * \real{0.2632}}@{}}
\toprule\noalign{}
\begin{minipage}[b]{\linewidth}\raggedright
field
\end{minipage} & \begin{minipage}[b]{\linewidth}\raggedright
symbol
\end{minipage} & \begin{minipage}[b]{\linewidth}\raggedright
what it is
\end{minipage} & \begin{minipage}[b]{\linewidth}\raggedright
lives in
\end{minipage} & \begin{minipage}[b]{\linewidth}\raggedright
couples to
\end{minipage} \\
\midrule\noalign{}
\endhead
\bottomrule\noalign{}
\endlastfoot
Einstein gravity & \(g_{MN}\) & the metric of eight dimensions; the
author's Einstein-Lovelock action with the weights of the Gauss-Bonnet
and cubic terms zero and \(\Lambda = -21H^2\) & bulk and fold & every
field \\
dirac16complex & \(\Psi\) & the author's 16-component complex spinor
field, with left- and right-handed halves \(\Psi_{\mathrm{L}}\),
\(\Psi_{\mathrm{R}}\) & the bulk & gravity (canonical spin connection),
itself, every Standard-Model field at the fold; only
\(\Psi_{\mathrm{L}}\) to the Higgs doublet \\
gluons & \(G^A\) & the gauge field of \(SU(3)_{\mathrm{C}}\) & the bulk;
its massless mode is bound to the fold & gravity, quarks,
\(\Psi_{\mathrm{L}}\), \(\Psi_{\mathrm{R}}\) (at the fold) \\
electroweak gauge fields & \(W^1\), \(W^2\), \(W^3\), \(B\) & the gauge
fields of \(SU(2)_{\mathrm{L}}\) and \(U(1)_Y\); the photon is the
combination of \(W^3\) and \(B\) that the Higgs vacuum leaves massless &
the bulk; their massless modes are bound to the fold & gravity,
\(\Phi\), the fermions, \(\Psi_{\mathrm{L}}\), \(\Psi_{\mathrm{R}}\) (at
the fold) \\
Higgs doublet & \(\Phi\) & the scalar doublet with
\(V = \mu^2\Phi^\dagger\Phi + \lambda(\Phi^\dagger\Phi)^2\),
\(\lambda > 0\), \(\mu^2 < 0\), expanded about the ad hoc vacuum
\((\mathcal{A}, v + h)/\sqrt2\) & the fold & gravity, \(W\), \(B\), the
fermions, \(\Psi_{\mathrm{L}}\) \\
lepton doublets & \(\ell_i\) & \((\nu_{\mathrm{L}i}, e_{\mathrm{L}i})\),
three generations & the fold & gravity, \(W\), \(B\), \(\Phi\),
\(\Psi_{\mathrm{L}}\), \(\Psi_{\mathrm{R}}\) \\
right-handed charged leptons & \(e_{\mathrm{R}i}\) & electron, muon, tau
& the fold & gravity, \(B\), \(\Phi\), \(\Psi_{\mathrm{L}}\),
\(\Psi_{\mathrm{R}}\) \\
right-handed neutrinos & \(\nu_{\mathrm{R}i}\) & three gauge singlets:
Dirac neutrino masses (accepted by the author, 2026-10-24 17:32) & the
fold & gravity, \(\Phi\), \(\Psi_{\mathrm{L}}\),
\(\Psi_{\mathrm{R}}\) \\
quarks & \(Q_i\), \(u_{\mathrm{R}i}\), \(d_{\mathrm{R}i}\) & three
generations of doublets and singlets & the fold & gravity, gluons,
\(W\), \(B\), \(\Phi\), \(\Psi_{\mathrm{L}}\), \(\Psi_{\mathrm{R}}\) \\
fold fluid & \(\mathcal{T}_{\mathrm{ff}}\) & the part of the thin source
of the fold not supplied by the other fields: a tension, which the
static block requires, and a pressure along the extra times that holds
their volume fixed & the fold & gravity \\
induced gravity & \(\kappa_{\mathrm{ind}}\) & an Einstein-Hilbert term
that the fields of the fold induce; negligible for a cutoff of the order
of \(H\) & the fold & gravity \\
\end{longtable}

The order places the photon in the bulk, and with it the electroweak
gauge fields of which it is a combination; UM-4 places the gluons there
as well, so that all gauge fields of the Standard Model live in one
place. In the metric of UM-4 every massless field of the bulk has a mode
bound to the fold that travels at the speed of light at every depth and
every wavelength, and these modes are the photon, the gluons and the
electroweak gauge fields of the four-dimensional world, with the
measured couplings. The Higgs doublet, the quarks and the leptons stay
on the fold, where their couplings and their flavour physics are those
of the Standard Model with three Dirac neutrinos (in the bulk they could
also travel at the speed of light, but their Yukawa couplings would be
set by the overlaps of eighteen profiles). dirac16complex, the author's
field of eight dimensions, lives in the bulk and couples to the fields
of the fold at the fold. No field depends on the extra times
\(x_5, x_6, x_7\): a field with a momentum along them is tachyonic (its
frequency squared falls with that momentum, because the extra times are
time-like); the restriction is consistent but unstable, and UM-4 needs a
gauge principle that removes the momentum along the extra times of every
field. Status: AUTHOR (the list of fields; the photon and the graviton
in the bulk, by the order), ASSUMED (the additions, the other places and
the restriction), DERIVED (the speeds of the fields of the bulk, the
tachyons), OPEN (the gauge principle).

\hypertarget{the-action}{%
\section{The action}\label{the-action}}

All field equations of UM-4 follow from one action by varying it with
respect to every field.

\hypertarget{the-sum-of-all-terms}{%
\subsection{The sum of all terms}\label{the-sum-of-all-terms}}

\[
\begin{aligned}
S_{\mathrm{tot}} &= S_{\mathrm{EL}} + S_\Psi + S_{\mathrm{gauge}} + S_{\mathrm{fold}},\\
S_{\mathrm{fold}} &= \int\mathrm{d}^8X\,\sqrt{\lvert g\rvert}\;\delta_{\mathrm{D}}(\sigma)\,\frac{1}{V_*}\,
\big(\mathcal{L}_{\mathrm{IG}} + \mathcal{L}_{\mathrm{ff}} + \mathcal{L}_{\mathrm{H}} + \mathcal{L}_{\mathrm{f}}
+ \mathcal{L}_{\mathrm{Y}} + \mathcal{L}_{\mathrm{P}}\big)
\end{aligned} \tag{IN-1}
\] \textbf{Symbols in (IN-1).}

\begin{itemize}
\tightlist
\item
  \(S_{\mathrm{tot}}\): The action of UM-4 (the integral of the
  Lagrangian density over spacetime); its variation with respect to
  every field gives the field equations. Unit: dimensionless (in units
  in which the reduced Planck constant is 1). Status: ASSUMED.
\item
  \(S_{\mathrm{EL}}\): The Einstein-Lovelock action of gravity in eight
  dimensions; in UM-4 the Einstein-Hilbert action with the cosmological
  constant. Unit: dimensionless. Status: AUTHOR.
\item
  \(S_\Psi\): The action of dirac16complex in the bulk (kinetic term
  with the canonical spin connection, mass term, contact
  self-interaction). Unit: dimensionless. Status: RECORD.
\item
  \(S_{\mathrm{gauge}}\): The action of the gauge fields in the bulk.
  Unit: dimensionless. Status: STANDARD.
\item
  \(S_{\mathrm{fold}}\): The action of the fold: the Standard Model, the
  portals of dirac16complex, the fold fluid and the Einstein-Hilbert
  term that the fields of the fold induce, all at \(\sigma = 0\). Unit:
  dimensionless. Status: ASSUMED.
\item
  \(\mathrm{d}\): The differential: \(\mathrm{d}\) in front of a
  quantity is an infinitesimal change of it (as \(\mathrm{d}\sigma\), a
  small step of \(\sigma\)); the quotient of two such changes is a
  derivative; \(\mathrm{d}\) with an exponent, in front of several
  variables, is their volume element, the product of their
  differentials. Status: STANDARD.
\item
  \(X = \{x_0, \dots, x_7\}\): The author's coordinates of the
  eight-dimensional spacetime, \(X = \{x_0, x_1, \dots, x_7\}\) with the
  hidden direction first (Dirac\_Einstein-Lovelock\_4+4\_Nash.nb,
  section DEFCHART); \(X^M\) is the coordinate with the index \(M\) and
  \(\mathrm{d}^8X = \mathrm{d}x_0\,\mathrm{d}x_1\cdots\mathrm{d}x_7\)
  the coordinate volume element. Unit: length (each coordinate). Status:
  AUTHOR.
\item
  \(g_{MN}\): The metric tensor of the eight-dimensional spacetime
  (signature \(+,+,+,+,-,-,-,-\) in the order \(x_0, \dots, x_7\)):
  \(g_{MN}\,\mathrm{d}X^M\mathrm{d}X^N\) is the squared interval,
  \(g^{MN}\) its inverse, \(\lvert g\rvert\) the absolute value of its
  determinant (\(\sqrt{\lvert g\rvert}\,\mathrm{d}^8X\) is the proper
  volume element). Unit: dimensionless. Status: AUTHOR.
\item
  \(\delta_{\mathrm{D}}\): The Dirac delta function of \(\sigma\): zero
  for \(\sigma \ne 0\), with the integral 1 over any range that contains
  \(\sigma = 0\); the integral of \(\delta_{\mathrm{D}}(\sigma)\) times
  a smooth function of \(\sigma\) is the value of that function at
  \(\sigma = 0\). Unit: 1/length (for \(\delta_{\mathrm{D}}(\sigma)\)).
  Status: STANDARD.
\item
  \(\sigma\): The proper coordinate of the hidden direction:
  \(\mathrm{d}\sigma = \lvert\cot z\rvert\,\mathrm{d}x_0\),
  \(\sigma = \mathrm{sgn}(\pi/2 - z)\ln(\sin z)/(6H)\) on one period
  \(0 < z < \pi\); \(\sigma = 0\) is the fold \(z = \pi/2\),
  \(\sigma \to \mp\infty\) the tips \(z \to 0, \pi\); one period of
  \(x_0\) is the whole line \(-\infty < \sigma < \infty\). Unit: length.
  Status: DERIVED.
\item
  \(V_*\): The physical volume of the extra times \(x_5, x_6, x_7\) on
  the fold in the late universe,
  \(V_* = V_{\mathrm{t}}(\mathrm{e}^{-a_4}b_4)^3\) (with
  \(V_{\mathrm{t}} = \int\mathrm{d}x_5\,\mathrm{d}x_6\,\mathrm{d}x_7\)
  the coordinate volume of the extra times) at a time when
  \(\mathrm{e}^{-a_4}b_4\) is constant; it converts the
  eight-dimensional normalisation into the four-dimensional one. Unit:
  length\(^3\) (GeV\(^{-3}\)). Status: ASSUMED.
\item
  \(\mathcal{L}_{\mathrm{IG}}\): The Lagrangian of the Einstein-Hilbert
  term of the four-dimensional metric of the fold that the quantum
  fluctuations of the fields of the fold induce; it adds to the
  Einstein-Hilbert term that the graviton of the bulk gives the fold and
  does not change the speed of gravitational waves. Unit: GeV\(^4\).
  Status: ESTIMATE.
\item
  \(\mathcal{L}_{\mathrm{ff}}\): The Lagrangian of the fold fluid (only
  its energy-momentum tensor \(\mathcal{T}_{\mathrm{ff}}\) enters the
  equations). Unit: GeV\(^4\). Status: ASSUMED.
\item
  \(\mathcal{L}_{\mathrm{H}}\): The Lagrangian of the Higgs doublet on
  the fold. Unit: GeV\(^4\). Status: AUTHOR.
\item
  \(\mathcal{L}_{\mathrm{f}}\): The kinetic Lagrangian of the quarks and
  leptons on the fold. Unit: GeV\(^4\). Status: STANDARD.
\item
  \(\mathcal{L}_{\mathrm{Y}}\): The Yukawa Lagrangian of the quarks and
  leptons with the Higgs doublet. Unit: GeV\(^4\). Status: STANDARD.
\item
  \(\mathcal{L}_{\mathrm{P}}\): The Lagrangian of the portals of
  dirac16complex at the fold. Unit: GeV\(^4\). Status: INVENTED.
\end{itemize}

The action has four parts: Einstein gravity in eight dimensions,
dirac16complex in the bulk, the gauge fields in the bulk, and the fold,
whose terms are confined to \(\sigma = 0\) by the Dirac delta of the
proper hidden coordinate; the bulk needs no other source (the static
block is a vacuum of Einstein gravity with \(\Lambda = -21H^2\)).
Status: ASSUMED.

\hypertarget{gravity}{%
\subsection{Gravity}\label{gravity}}

\[
\begin{aligned}
S_{\mathrm{EL}} &= \frac{1}{2\kappa}\int\mathrm{d}^8X\,\sqrt{\lvert g\rvert}\,\Big(\sum_{p=1}^{3}\alpha_{(p)}\,\mathcal{L}_{(p)}
 - 2\Lambda\Big) = \frac{1}{2\kappa}\int\mathrm{d}^8X\,\sqrt{\lvert g\rvert}\,\big(R - 2\Lambda\big),\qquad \kappa = 8\pi G_8,\\
\alpha_{(p)} &= w_{(p)}\,H^{2-2p},\qquad w_{(1)} = 1,\quad w_{(2)} = w_{(3)} = 0,\qquad \Lambda = -21H^2
\end{aligned} \tag{IN-2}
\] \textbf{Symbols in (IN-2).}

\begin{itemize}
\tightlist
\item
  \(S_{\mathrm{EL}}\): The Einstein-Lovelock action of gravity in eight
  dimensions; in UM-4 the Einstein-Hilbert action with the cosmological
  constant. Unit: dimensionless. Status: AUTHOR.
\item
  \(\kappa\): The eight-dimensional gravitational coupling,
  \(\kappa = 8\pi G_8\). Unit: GeV\(^{-6}\). Status: AUTHOR.
\item
  \(\mathrm{d}\): The differential: \(\mathrm{d}\) in front of a
  quantity is an infinitesimal change of it (as \(\mathrm{d}\sigma\), a
  small step of \(\sigma\)); the quotient of two such changes is a
  derivative; \(\mathrm{d}\) with an exponent, in front of several
  variables, is their volume element, the product of their
  differentials. Status: STANDARD.
\item
  \(X = \{x_0, \dots, x_7\}\): The author's coordinates of the
  eight-dimensional spacetime, \(X = \{x_0, x_1, \dots, x_7\}\) with the
  hidden direction first (Dirac\_Einstein-Lovelock\_4+4\_Nash.nb,
  section DEFCHART); \(X^M\) is the coordinate with the index \(M\) and
  \(\mathrm{d}^8X = \mathrm{d}x_0\,\mathrm{d}x_1\cdots\mathrm{d}x_7\)
  the coordinate volume element. Unit: length (each coordinate). Status:
  AUTHOR.
\item
  \(g_{MN}\): The metric tensor of the eight-dimensional spacetime
  (signature \(+,+,+,+,-,-,-,-\) in the order \(x_0, \dots, x_7\)):
  \(g_{MN}\,\mathrm{d}X^M\mathrm{d}X^N\) is the squared interval,
  \(g^{MN}\) its inverse, \(\lvert g\rvert\) the absolute value of its
  determinant (\(\sqrt{\lvert g\rvert}\,\mathrm{d}^8X\) is the proper
  volume element). Unit: dimensionless. Status: AUTHOR.
\item
  \(p\): The Lovelock order \(p = 1, 2, 3\) (Einstein, Gauss-Bonnet,
  cubic), as an index and as a number (UM-4 keeps only \(p = 1\)); in
  the Kohn-Sham equations also a second label of a Kohn-Sham orbital,
  with the Kronecker delta of the two labels. Status: STANDARD.
\item
  \(\alpha_{(p)}\): The Lovelock coupling of order \(p\),
  \(\alpha_{(p)} = w_{(p)}H^{2-2p}\): in UM-4 \(\alpha_{(1)} = 1\) and
  \(\alpha_{(2)} = \alpha_{(3)} = 0\) (Einstein gravity). Unit:
  GeV\(^{2-2p}\). Status: AUTHOR.
\item
  \(\mathcal{L}_{(p)}\): The Lovelock density of order \(p\),
  \(\mathcal{L}_{(p)} = 2^{-p}\,\delta^{P_1\cdots P_{2p}}_{Q_1\cdots Q_{2p}}\,R^{Q_1Q_2}{}_{P_1P_2}\cdots R^{Q_{2p-1}Q_{2p}}{}_{P_{2p-1}P_{2p}}\),
  a scalar built from \(p\) curvature tensors: \(\mathcal{L}_{(1)} = R\)
  (Einstein-Hilbert), \(\mathcal{L}_{(2)}\) the Gauss-Bonnet density,
  \(\mathcal{L}_{(3)}\) the cubic Lovelock density. Unit:
  1/length\(^{2p}\). Status: STANDARD.
\item
  \(\Lambda\): The eight-dimensional cosmological constant of UM-4,
  \(\Lambda = -21H^2\): the value for which the static block is a vacuum
  of Einstein gravity (its Einstein tensor is \(21H^2\) in every
  direction). The author writes the pure number \(\Lambda_{\mathrm{N}}\)
  in
  \(0 = -\Lambda_{\mathrm{N}} + \sum_p H^{-2p}w_{(p)}\mathcal{E}_{(p)}\),
  so that \(\Lambda = -\Lambda_{\mathrm{N}}H^2\). Unit: GeV\(^2\).
  Value: \(-21H^2\); the author's original \(-10H^2\); UM-3: \(-7H^2\).
  Status: AUTHOR (the term), DERIVED (the value).
\item
  \(R^{MN}{}_{PQ},\ R_{MN},\ R\): The curvature of the metric
  \(g_{MN}\): with four indices the Riemann tensor \(R^M{}_{NPQ}\) (in
  the convention in which the commutator of two covariant derivatives,
  acting on a vector, is the Riemann tensor, with its last two indices
  those of the two derivatives, contracted with the vector), with two
  indices the Ricci tensor \(R_{MN} = R^P{}_{MPN}\), without indices the
  scalar curvature \(R = g^{MN}R_{MN}\); with frame indices
  \(R^{(a)(b)}{}_{(c)(d)}\) its frame components,
  \(R^{(a)(b)}{}_{(a)(b)}\) (no sum) the sectional curvature of the
  frame plane \((a, b)\). Unit: 1/length\(^2\). Status: STANDARD.
\item
  \(\pi\): The number pi, the ratio of the circumference of a circle to
  its diameter. Unit: dimensionless. Value: 3.141592654. Status:
  STANDARD.
\item
  \(G_8\): The eight-dimensional Newton constant
  (\(\kappa = 8\pi G_8\)). Unit: GeV\(^{-6}\). Status: AUTHOR.
\item
  \(w_{(p)}\): The Lovelock weights of the author's Einstein-Lovelock
  action, \(w_{(1)} = 1\) and \(w_{(2)}\), \(w_{(3)}\) of the
  Gauss-Bonnet and the cubic densities. UM-4 sets
  \(w_{(2)} = w_{(3)} = 0\) (Einstein gravity): the only weights for
  which gravitons propagate on the light cone on every background, not
  only around the static block. The author's original weights
  \(w_{(2)} = 1/32\), \(w_{(3)} = 1/1152\)
  (Dirac\_Einstein-Lovelock\_4+4\_Nash.nb, In{[}134{]}) keep the
  gravitons on the light cone only around the static block (with
  \(\Lambda = -161H^2/16\)); UM-3 used \((1, 1/20, 1/360)\). Unit:
  dimensionless. Value: \((1, 0, 0)\) in UM-4; the author's original
  \((1, 1/32, 1/1152)\). Status: AUTHOR (the form), DERIVED (the values
  that the order of 2026-10-25 selects).
\item
  \(H\): The author's constant (an inverse length) that fixes the period
  \(\pi/(6H)\) of \(x_0\), the warp factor
  \(\mathrm{e}^{-2H\lvert\sigma\rvert}\) of the metric, the cosmological
  constant \(\Lambda = -21H^2\) and the tension
  \(\kappa\lambda_{\mathrm{F}} = 12H\) of the fold; in UM-4 a free
  parameter, bounded from below by the propagation of the gauge fields
  in the bulk. Unit: GeV (1/length). Value: \(H \gtrsim 2.5\) TeV (from
  the invisible width of the Z boson). Status: AUTHOR.
\end{itemize}

\[
\begin{aligned}
\mathcal{L}_{(p)} &= \frac{1}{2^{p}}\,\delta^{M_1N_1\cdots M_pN_p}_{P_1Q_1\cdots P_pQ_p}\,
 R^{P_1Q_1}{}_{M_1N_1}\cdots R^{P_pQ_p}{}_{M_pN_p},\\
\mathcal{L}_{(1)} &= R,\qquad \mathcal{L}_{(2)} = R^2 - 4R_{MN}R^{MN} + R_{MNPQ}R^{MNPQ}
\end{aligned} \tag{IN-3}
\] \textbf{Symbols in (IN-3).}

\begin{itemize}
\tightlist
\item
  \(\mathcal{L}_{(p)}\): The Lovelock density of order \(p\),
  \(\mathcal{L}_{(p)} = 2^{-p}\,\delta^{P_1\cdots P_{2p}}_{Q_1\cdots Q_{2p}}\,R^{Q_1Q_2}{}_{P_1P_2}\cdots R^{Q_{2p-1}Q_{2p}}{}_{P_{2p-1}P_{2p}}\),
  a scalar built from \(p\) curvature tensors: \(\mathcal{L}_{(1)} = R\)
  (Einstein-Hilbert), \(\mathcal{L}_{(2)}\) the Gauss-Bonnet density,
  \(\mathcal{L}_{(3)}\) the cubic Lovelock density. Unit:
  1/length\(^{2p}\). Status: STANDARD.
\item
  \(p\): The Lovelock order \(p = 1, 2, 3\) (Einstein, Gauss-Bonnet,
  cubic), as an index and as a number (UM-4 keeps only \(p = 1\)); in
  the Kohn-Sham equations also a second label of a Kohn-Sham orbital,
  with the Kronecker delta of the two labels. Status: STANDARD.
\item
  \(\delta\): The Kronecker delta: with two indices \(\delta^M_N = 1\)
  if \(M = N\) and 0 otherwise (likewise \(\delta_{X\mathrm{L}}\) is 1
  if the label \(X\) of a half of dirac16complex is L and 0 otherwise);
  with two index strings of equal length the generalised Kronecker delta
  \(\delta^{M_1M_2\cdots}_{N_1N_2\cdots}\), the determinant of the
  square matrix of the Kronecker deltas \(\delta^{M_a}_{N_b}\) (\(+1\)
  or \(-1\) if the upper indices are an even or odd permutation of the
  lower ones, 0 otherwise). Status: STANDARD.
\item
  \(M\): A coordinate index of eight dimensions (also with a numbered
  copy \(M_1, M_2, \dots\)): it runs over \(x_0, x_1, \dots, x_7\); a
  repeated index is summed. Status: STANDARD.
\item
  \(N\): A coordinate index of eight dimensions (also
  \(N_1, N_2, \dots\)), running over \(x_0, \dots, x_7\); a repeated
  index is summed. Status: STANDARD.
\item
  \(P\): A coordinate index of eight dimensions (also
  \(P_1, P_2, \dots\)), running over \(x_0, \dots, x_7\); a repeated
  index is summed. Status: STANDARD.
\item
  \(Q\): A coordinate index of eight dimensions (also
  \(Q_1, Q_2, \dots\)), running over \(x_0, \dots, x_7\); a repeated
  index is summed. Status: STANDARD.
\item
  \(R^{MN}{}_{PQ},\ R_{MN},\ R\): The curvature of the metric
  \(g_{MN}\): with four indices the Riemann tensor \(R^M{}_{NPQ}\) (in
  the convention in which the commutator of two covariant derivatives,
  acting on a vector, is the Riemann tensor, with its last two indices
  those of the two derivatives, contracted with the vector), with two
  indices the Ricci tensor \(R_{MN} = R^P{}_{MPN}\), without indices the
  scalar curvature \(R = g^{MN}R_{MN}\); with frame indices
  \(R^{(a)(b)}{}_{(c)(d)}\) its frame components,
  \(R^{(a)(b)}{}_{(a)(b)}\) (no sum) the sectional curvature of the
  frame plane \((a, b)\). Unit: 1/length\(^2\). Status: STANDARD.
\end{itemize}

Gravity is the author's Einstein-Lovelock theory (his notebook,
In{[}134{]}) with the weights of the Gauss-Bonnet density
\(\mathcal{L}_{(2)}\) and of the cubic Lovelock density
\(\mathcal{L}_{(3)}\) set to zero: Einstein gravity with the
cosmological constant \(\Lambda = -21H^2\), for which the static block
is a vacuum (its Einstein tensor is \(21H^2\) in every direction;
DERIVED, \texttt{um4\_geometry.py}, G4). The order requires gravitons at
the speed of light in the bulk. Around the static block every
non-degenerate choice of the weights gives that (the graviton equation
is, for any weights, the wave equation of the light cone of the fold
times the factor \(1 - 40w_{(2)} + 360w_{(3)}\); DERIVED,
\texttt{um4\_gravity.py}, L2), but on every background that is not of
constant curvature (an expanding fold, a star, a gravitational wave of
finite amplitude) the Gauss-Bonnet and cubic terms give the coefficients
of the time and space derivatives of a graviton different factors, and
gravitons leave the light cone; only \(w_{(2)} = w_{(3)} = 0\) keeps
them on it on every background (DERIVED, L3). The author's original
weights \(w_{(2)} = 1/32\), \(w_{(3)} = 1/1152\) would make the static
block a vacuum only for \(\Lambda = -161H^2/16\) (not his \(-10H^2\)),
with gravitons at the speed of light around it but not around an
expanding fold; UM-3's weights \(1/20\), \(1/360\) make the static block
a degenerate vacuum without linear gravitons. Lovelock's densities
contain no derivative of the metric higher than the second in their
variation, so a thin source at the fold produces a finite jump of first
derivatives (the junction condition). Status: AUTHOR (the action),
DERIVED (the weights and \(\Lambda\) that the order selects).

\hypertarget{dirac16complex}{%
\subsection{dirac16complex}\label{dirac16complex}}

\[
S_\Psi = \int\mathrm{d}^8X\,\sqrt{\lvert g\rvert}\,\Big[\frac12\big(\bar\Psi\gamma^MD_M\Psi - (D_M\bar\Psi)\gamma^M\Psi\big)
 - m_\Psi\,\bar\Psi\Psi - \frac{\lambda_\Psi}{2}\big(\bar\Psi\Psi\big)^2\Big] \tag{IN-4}
\] \textbf{Symbols in (IN-4).}

\begin{itemize}
\tightlist
\item
  \(S_\Psi\): The action of dirac16complex in the bulk (kinetic term
  with the canonical spin connection, mass term, contact
  self-interaction). Unit: dimensionless. Status: RECORD.
\item
  \(\mathrm{d}\): The differential: \(\mathrm{d}\) in front of a
  quantity is an infinitesimal change of it (as \(\mathrm{d}\sigma\), a
  small step of \(\sigma\)); the quotient of two such changes is a
  derivative; \(\mathrm{d}\) with an exponent, in front of several
  variables, is their volume element, the product of their
  differentials. Status: STANDARD.
\item
  \(X = \{x_0, \dots, x_7\}\): The author's coordinates of the
  eight-dimensional spacetime, \(X = \{x_0, x_1, \dots, x_7\}\) with the
  hidden direction first (Dirac\_Einstein-Lovelock\_4+4\_Nash.nb,
  section DEFCHART); \(X^M\) is the coordinate with the index \(M\) and
  \(\mathrm{d}^8X = \mathrm{d}x_0\,\mathrm{d}x_1\cdots\mathrm{d}x_7\)
  the coordinate volume element. Unit: length (each coordinate). Status:
  AUTHOR.
\item
  \(g_{MN}\): The metric tensor of the eight-dimensional spacetime
  (signature \(+,+,+,+,-,-,-,-\) in the order \(x_0, \dots, x_7\)):
  \(g_{MN}\,\mathrm{d}X^M\mathrm{d}X^N\) is the squared interval,
  \(g^{MN}\) its inverse, \(\lvert g\rvert\) the absolute value of its
  determinant (\(\sqrt{\lvert g\rvert}\,\mathrm{d}^8X\) is the proper
  volume element). Unit: dimensionless. Status: AUTHOR.
\item
  \(\Psi,\ \bar\Psi\): dirac16complex: the author's 16-component complex
  spinor field of eight dimensions;
  \(\bar\Psi = \Psi^\dagger\sigma_{16}\) is its Dirac adjoint; in a term
  of the fold, \(\Psi\) is its value at \(\sigma = 0\). Unit:
  GeV\(^{7/2}\). Status: AUTHOR.
\item
  \(\gamma^M,\ \gamma^{(a)}\): The Dirac matrices of dirac16complex
  along the coordinates, \(\gamma^M = \mathbf{e}_a{}^MT^a\),
  \(\gamma_M = g_{MN}\gamma^N\); \(\gamma^{(a)} = T^a\) along the frame
  direction \((a)\) (\(\gamma^{(0)}\) along \(\sigma\), \(\gamma^{(4)}\)
  along the time). Unit: dimensionless. Status: AUTHOR.
\item
  \(M\): A coordinate index of eight dimensions (also with a numbered
  copy \(M_1, M_2, \dots\)): it runs over \(x_0, x_1, \dots, x_7\); a
  repeated index is summed. Status: STANDARD.
\item
  \(D_M,\ D_\mu\): The covariant derivative: on dirac16complex
  \(D_M = \partial_M + \Omega_M\) (with
  \(D_M\bar\Psi = \partial_M\bar\Psi - \bar\Psi\Omega_M\)); on the
  fields of the fold the gauge and spin covariant derivative \(D_\mu\)
  (the partial derivative plus the spin connection of the fold and the
  gauge fields times their couplings and generators);
  \(D^\mu = \hat h^{\mu\nu}D_\nu\); on an adjoint field
  \((D_\mu X)^A = \partial_\mu X^A + g_3f^{ABC}G^B_\mu X^C\) (likewise
  for \(SU(2)\) with \(g_2\epsilon^{IJK}\)). Unit: 1/length. Status:
  STANDARD.
\item
  \(m_\Psi\): The mass parameter of dirac16complex in the bulk. Unit:
  GeV. Status: RECORD.
\item
  \(\lambda_\Psi\): The coupling of the contact self-interaction of
  dirac16complex. Unit: GeV\(^{-6}\). Status: RECORD.
\end{itemize}

\[
\begin{aligned}
&\{T_a, T_b\} = 2\eta_{ab}\mathbf{1},\qquad \eta = \mathrm{diag}(1, 1, 1, 1, -1, -1, -1, -1),\qquad T^a = \eta^{ab}T_b,
 \qquad \gamma^M = \mathbf{e}_a{}^M\,T^a,\\
&\bar\Psi = \Psi^\dagger\sigma_{16},\qquad \sigma_{16} = T_0T_1T_2T_3,\qquad \Gamma_8 = T_0T_1T_2T_3T_4T_5T_6T_7,\\
&P_{\mathrm{L}} = \frac{\mathbf{1} - \Gamma_8}{2},\qquad P_{\mathrm{R}} = \frac{\mathbf{1} + \Gamma_8}{2},\qquad
 \Psi_{\mathrm{L}} = P_{\mathrm{L}}\Psi,\qquad \Psi_{\mathrm{R}} = P_{\mathrm{R}}\Psi
\end{aligned} \tag{IN-5}
\] \textbf{Symbols in (IN-5).}

\begin{itemize}
\tightlist
\item
  \(T_a,\ T^a\): The author's eight real \(16 \times 16\) Dirac matrices
  T16\^{}A{[}a{]} of 4+4.nb (\(a = 0, \dots, 7\); \(T_0, \dots, T_7\)),
  with \(\{T_a, T_b\} = 2\eta_{ab}\mathbf{1}\); \(T^a = \eta^{ab}T_b\).
  Unit: dimensionless. Status: AUTHOR.
\item
  \(a\): A frame (orthonormal) index running over \(0, 1, \dots, 7\)
  (frame direction \((0)\) along \(\sigma\), \((1), (2), (3)\) along
  ordinary space, \((4)\) along the time, \((5), (6), (7)\) along the
  extra times); raised and lowered with the frame metric
  \(\eta_{ab} = \mathrm{diag}(1, 1, 1, 1, -1, -1, -1, -1)\); on the fold
  (\(a, b \in \{1, \dots, 7\}\)) the frame of the fold; in the tensors
  of the fold written with its coordinates \(x^a\) (its metric
  \(\hat g_{ab}\) and the tensors built from it) a coordinate index of
  the fold, \(a \in \{1, \dots, 7\}\); a repeated index is summed.
  Status: STANDARD.
\item
  \(b\): A frame index (as \(a\)). Status: STANDARD.
\item
  \(\eta_{ab}\): The flat metric of the frame,
  \(\eta = \mathrm{diag}(1, 1, 1, 1, -1, -1, -1, -1)\) (eta4488 of
  4+4.nb); \(\eta^{ab}\) is its inverse (the same matrix). Unit:
  dimensionless. Status: AUTHOR.
\item
  \(\mathbf{1},\ \mathbf{1}_8\): The identity matrix of the size fixed
  by the context (\(16 \times 16\) on dirac16complex, \(4 \times 4\) on
  a four-dimensional spinor, \(3 \times 3\) in generation or flavour
  space); \(\mathbf{1}_8\) is the \(8 \times 8\) identity matrix. Unit:
  dimensionless. Status: STANDARD.
\item
  \(\mathrm{diag}\): A diagonal matrix with the listed entries on its
  diagonal (in the order of the coordinates \(x_0, x_1, \dots, x_7\) for
  a tensor of eight dimensions, \(x_1, \dots, x_7\) for a tensor of the
  fold, \(x_1, \dots, x_4\) for a four-dimensional tensor). Status:
  STANDARD.
\item
  \(\gamma^M,\ \gamma^{(a)}\): The Dirac matrices of dirac16complex
  along the coordinates, \(\gamma^M = \mathbf{e}_a{}^MT^a\),
  \(\gamma_M = g_{MN}\gamma^N\); \(\gamma^{(a)} = T^a\) along the frame
  direction \((a)\) (\(\gamma^{(0)}\) along \(\sigma\), \(\gamma^{(4)}\)
  along the time). Unit: dimensionless. Status: AUTHOR.
\item
  \(M\): A coordinate index of eight dimensions (also with a numbered
  copy \(M_1, M_2, \dots\)): it runs over \(x_0, x_1, \dots, x_7\); a
  repeated index is summed. Status: STANDARD.
\item
  \(\mathbf{e}_a{}^M,\ \mathbf{e}^a{}_M\): The vielbein (frame field) of
  the metric: \(g_{MN} = \eta_{ab}\,\mathbf{e}^a{}_M\mathbf{e}^b{}_N\);
  \(\mathbf{e}_a{}^M\) is its inverse;
  \(\mathbf{e}_{bN} = \eta_{bc}\mathbf{e}^c{}_N\); \(\mathbf{e}^{(a)}\)
  is the frame one-form \(\mathbf{e}^a{}_M\,\mathrm{d}X^M\). Unit:
  dimensionless. Status: STANDARD.
\item
  \(\Psi,\ \bar\Psi\): dirac16complex: the author's 16-component complex
  spinor field of eight dimensions;
  \(\bar\Psi = \Psi^\dagger\sigma_{16}\) is its Dirac adjoint; in a term
  of the fold, \(\Psi\) is its value at \(\sigma = 0\). Unit:
  GeV\(^{7/2}\). Status: AUTHOR.
\item
  \({}^\dagger\): Hermitian conjugation: the complex conjugate of the
  transpose of a matrix or spinor; for an operator on the space of
  states, its adjoint with respect to the scalar product of that space
  (for the quantised dirac16complex the indefinite Krein scalar
  product). Status: STANDARD.
\item
  \(\sigma_{16}\): The matrix of the Dirac adjoint,
  \(\sigma_{16} = T_0T_1T_2T_3\), real and symmetric; \(\sigma_{16}T_a\)
  is real and antisymmetric. Unit: dimensionless. Status: AUTHOR.
\item
  \(\Gamma_8\): The chirality matrix of eight dimensions,
  \(\Gamma_8 = T_0T_1\cdots T_7\) (the author's T16\^{}A{[}8{]})
  \(= \mathrm{diag}(-\mathbf{1}_8, \mathbf{1}_8)\);
  \(\Gamma_8^2 = \mathbf{1}\); it anticommutes with every \(T_a\). Unit:
  dimensionless. Status: AUTHOR.
\item
  \(P_{\mathrm{L}}\): The author's left-handed chiral projector,
  \(P_{\mathrm{L}} = (\mathbf{1} - \Gamma_8)/2\). Unit: dimensionless.
  Status: AUTHOR.
\item
  \(P_{\mathrm{R}}\): The author's right-handed chiral projector,
  \(P_{\mathrm{R}} = (\mathbf{1} + \Gamma_8)/2\). Unit: dimensionless.
  Status: AUTHOR.
\item
  \(\Psi_{\mathrm{L}}\): The left-handed half of dirac16complex,
  \(\Psi_{\mathrm{L}} = P_{\mathrm{L}}\Psi\) (the upper eight components
  in the basis of 4+4.nb). Unit: GeV\(^{7/2}\). Status: AUTHOR.
\item
  \(\Psi_{\mathrm{R}}\): The right-handed half of dirac16complex,
  \(\Psi_{\mathrm{R}} = P_{\mathrm{R}}\Psi\) (the lower eight
  components). Unit: GeV\(^{7/2}\). Status: AUTHOR.
\end{itemize}

\[
\begin{aligned}
&D_M\Psi = \partial_M\Psi + \Omega_M\Psi,\qquad D_M\bar\Psi = \partial_M\bar\Psi - \bar\Psi\,\Omega_M,\qquad
 \Omega_M = \frac12\,\varpi_{Mab}\,\Sigma^{ab},\qquad \Sigma^{ab} = \frac14\big[T^a, T^b\big],\\
&\varpi_{Mab} = \mathbf{e}_a{}^N\big(\partial_M\mathbf{e}_{bN} - \Gamma^P{}_{MN}\,\mathbf{e}_{bP}\big)
\end{aligned} \tag{IN-6}
\] \textbf{Symbols in (IN-6).}

\begin{itemize}
\tightlist
\item
  \(D_M,\ D_\mu\): The covariant derivative: on dirac16complex
  \(D_M = \partial_M + \Omega_M\) (with
  \(D_M\bar\Psi = \partial_M\bar\Psi - \bar\Psi\Omega_M\)); on the
  fields of the fold the gauge and spin covariant derivative \(D_\mu\)
  (the partial derivative plus the spin connection of the fold and the
  gauge fields times their couplings and generators);
  \(D^\mu = \hat h^{\mu\nu}D_\nu\); on an adjoint field
  \((D_\mu X)^A = \partial_\mu X^A + g_3f^{ABC}G^B_\mu X^C\) (likewise
  for \(SU(2)\) with \(g_2\epsilon^{IJK}\)). Unit: 1/length. Status:
  STANDARD.
\item
  \(M\): A coordinate index of eight dimensions (also with a numbered
  copy \(M_1, M_2, \dots\)): it runs over \(x_0, x_1, \dots, x_7\); a
  repeated index is summed. Status: STANDARD.
\item
  \(\partial\): A partial derivative: the derivative with respect to one
  variable at fixed values of the others (\(\partial_\sigma\) with
  respect to \(\sigma\); \(\partial_M\) and \(\partial_\mu\) with
  respect to the coordinate that carries the same index, of the eight
  dimensions or of the fold). Status: STANDARD.
\item
  \(\Psi,\ \bar\Psi\): dirac16complex: the author's 16-component complex
  spinor field of eight dimensions;
  \(\bar\Psi = \Psi^\dagger\sigma_{16}\) is its Dirac adjoint; in a term
  of the fold, \(\Psi\) is its value at \(\sigma = 0\). Unit:
  GeV\(^{7/2}\). Status: AUTHOR.
\item
  \(\Omega_M\): The spin connection acting on dirac16complex,
  \(\Omega_M = \frac12\varpi_{Mab}\Sigma^{ab}\). Unit: 1/length. Status:
  RECORD.
\item
  \(\varpi_{Mab}\): The canonical (Levi-Civita) spin connection of the
  vielbein, antisymmetric in \(a, b\). Unit: 1/length. Status: STANDARD.
\item
  \(a\): A frame (orthonormal) index running over \(0, 1, \dots, 7\)
  (frame direction \((0)\) along \(\sigma\), \((1), (2), (3)\) along
  ordinary space, \((4)\) along the time, \((5), (6), (7)\) along the
  extra times); raised and lowered with the frame metric
  \(\eta_{ab} = \mathrm{diag}(1, 1, 1, 1, -1, -1, -1, -1)\); on the fold
  (\(a, b \in \{1, \dots, 7\}\)) the frame of the fold; in the tensors
  of the fold written with its coordinates \(x^a\) (its metric
  \(\hat g_{ab}\) and the tensors built from it) a coordinate index of
  the fold, \(a \in \{1, \dots, 7\}\); a repeated index is summed.
  Status: STANDARD.
\item
  \(b\): A frame index (as \(a\)). Status: STANDARD.
\item
  \(\Sigma^{ab},\ \Sigma^{\mu\nu}\): The generators of Spin(4,4) on
  dirac16complex, \(\Sigma^{ab} = \frac14[T^a, T^b]\) (the scaled
  commutators of the author's matrices);
  \(\Sigma^{\mu\nu} = \mathbf{e}_a{}^\mu\mathbf{e}_b{}^\nu\Sigma^{ab}\)
  along the fold. Unit: dimensionless. Status: AUTHOR.
\item
  \(T_a,\ T^a\): The author's eight real \(16 \times 16\) Dirac matrices
  T16\^{}A{[}a{]} of 4+4.nb (\(a = 0, \dots, 7\); \(T_0, \dots, T_7\)),
  with \(\{T_a, T_b\} = 2\eta_{ab}\mathbf{1}\); \(T^a = \eta^{ab}T_b\).
  Unit: dimensionless. Status: AUTHOR.
\item
  \(\mathbf{e}_a{}^M,\ \mathbf{e}^a{}_M\): The vielbein (frame field) of
  the metric: \(g_{MN} = \eta_{ab}\,\mathbf{e}^a{}_M\mathbf{e}^b{}_N\);
  \(\mathbf{e}_a{}^M\) is its inverse;
  \(\mathbf{e}_{bN} = \eta_{bc}\mathbf{e}^c{}_N\); \(\mathbf{e}^{(a)}\)
  is the frame one-form \(\mathbf{e}^a{}_M\,\mathrm{d}X^M\). Unit:
  dimensionless. Status: STANDARD.
\item
  \(N\): A coordinate index of eight dimensions (also
  \(N_1, N_2, \dots\)), running over \(x_0, \dots, x_7\); a repeated
  index is summed. Status: STANDARD.
\item
  \(\Gamma^P{}_{MN}\): The Christoffel symbols of the metric \(g_{MN}\)
  (the Levi-Civita connection). Unit: 1/length. Status: STANDARD.
\item
  \(P\): A coordinate index of eight dimensions (also
  \(P_1, P_2, \dots\)), running over \(x_0, \dots, x_7\); a repeated
  index is summed. Status: STANDARD.
\end{itemize}

dirac16complex with the kinetic term in the symmetric form (real for the
author's real matrices), the mass \(m_\Psi\) and the contact
self-interaction \(\lambda_\Psi\) (RECORD); the author's eight real
\(16 \times 16\) Dirac matrices, the Dirac adjoint and the chiral
projectors (AUTHOR); the canonical spin connection, which couples
dirac16complex to gravity and is never identically zero for the metric
of UM-4:
\(\gamma^M\Omega_M = -\frac72H\,\mathrm{sgn}(\sigma)\gamma^{(0)} + \frac32(b_4'/b_4)\mathrm{e}^{H\lvert\sigma\rvert}\gamma^{(4)}\)
(DERIVED, \texttt{um4\_geometry.py}, G9; \(\gamma^{(0)} = T^0\) and
\(\gamma^{(4)} = T^4\) are the author's matrices along the hidden
direction and along the time, \(\sigma\) the proper coordinate of the
hidden direction with the fold at \(\sigma = 0\), \(\mathrm{sgn}\) the
sign function; in UM-3 the first term was
\(-3H\,\mathrm{sgn}(\sigma)\gamma^{(0)}\), because the time was not
warped). The factor \(\mathrm{e}^{(7/2)H\lvert\sigma\rvert}\) removes
the first term: a mode is
\(\Psi = \mathrm{e}^{(7/2)H\lvert\sigma\rvert}\chi(\sigma)\) times a
plane wave of the fold, and the warp multiplies the whole Dirac operator
of the fold, the frequency as well as the momentum, by
\(\mathrm{e}^{H\lvert\sigma\rvert}\). The states of dirac16complex at
the fold are therefore Lorentz invariant along the fold, and, because
the warped frequency exceeds the bulk mass far from the fold, they are
long-lived resonances rather than bound states (DERIVED,
\texttt{um4\_spectra.py}).

\hypertarget{the-gauge-fields-in-the-bulk}{%
\subsection{The gauge fields in the
bulk}\label{the-gauge-fields-in-the-bulk}}

\[
\begin{aligned}
S_{\mathrm{gauge}} &= \int\mathrm{d}^8X\,\sqrt{\lvert g\rvert}\;\mathcal{L}_{\mathrm{gauge}},\qquad
\mathcal{L}_{\mathrm{gauge}} = -\frac{1}{4\mathsf{g}_3^2}\,G^A_{MN}G^{AMN} - \frac{1}{4\mathsf{g}_2^2}\,W^I_{MN}W^{IMN}
 - \frac{1}{4\mathsf{g}_1^2}\,B_{MN}B^{MN},\\
G^A_{MN} &= \partial_MG^A_N - \partial_NG^A_M + f^{ABC}G^B_MG^C_N,\qquad
W^I_{MN} = \partial_MW^I_N - \partial_NW^I_M + \epsilon^{IJK}W^J_MW^K_N,\\
B_{MN} &= \partial_MB_N - \partial_NB_M
\end{aligned} \tag{IN-7}
\] \textbf{Symbols in (IN-7).}

\begin{itemize}
\tightlist
\item
  \(S_{\mathrm{gauge}}\): The action of the gauge fields in the bulk.
  Unit: dimensionless. Status: STANDARD.
\item
  \(\mathrm{d}\): The differential: \(\mathrm{d}\) in front of a
  quantity is an infinitesimal change of it (as \(\mathrm{d}\sigma\), a
  small step of \(\sigma\)); the quotient of two such changes is a
  derivative; \(\mathrm{d}\) with an exponent, in front of several
  variables, is their volume element, the product of their
  differentials. Status: STANDARD.
\item
  \(X = \{x_0, \dots, x_7\}\): The author's coordinates of the
  eight-dimensional spacetime, \(X = \{x_0, x_1, \dots, x_7\}\) with the
  hidden direction first (Dirac\_Einstein-Lovelock\_4+4\_Nash.nb,
  section DEFCHART); \(X^M\) is the coordinate with the index \(M\) and
  \(\mathrm{d}^8X = \mathrm{d}x_0\,\mathrm{d}x_1\cdots\mathrm{d}x_7\)
  the coordinate volume element. Unit: length (each coordinate). Status:
  AUTHOR.
\item
  \(g_{MN}\): The metric tensor of the eight-dimensional spacetime
  (signature \(+,+,+,+,-,-,-,-\) in the order \(x_0, \dots, x_7\)):
  \(g_{MN}\,\mathrm{d}X^M\mathrm{d}X^N\) is the squared interval,
  \(g^{MN}\) its inverse, \(\lvert g\rvert\) the absolute value of its
  determinant (\(\sqrt{\lvert g\rvert}\,\mathrm{d}^8X\) is the proper
  volume element). Unit: dimensionless. Status: AUTHOR.
\item
  \(\mathcal{L}_{\mathrm{gauge}}\): The Lagrangian of the gauge fields
  of \(SU(3)_C\times SU(2)_L\times U(1)_Y\) in the bulk. Unit:
  GeV\(^8\). Status: STANDARD.
\item
  \(\mathsf{g}_1,\ \mathsf{g}_2,\ \mathsf{g}_3\): The gauge couplings of
  eight dimensions of \(U(1)_Y\), \(SU(2)_L\) and \(SU(3)_C\): the
  coefficients of the kinetic terms of the gauge fields of the bulk in
  their geometric normalisation; the four-dimensional couplings are
  \(g_k^2 = 3H\mathsf{g}_k^2/(2V_*)\). Unit: GeV\(^{-2}\)
  (\(\mathsf{g}^2\): length\(^4\)). Status: ASSUMED.
\item
  \(G^A_M,\ G^A_\mu,\ G^A_{MN}\): The gluon field of \(SU(3)_C\)
  (\(A = 1, \dots, 8\)): in the bulk \(G^A_M\), a field of eight
  dimensions in its geometric normalisation (the covariant derivative
  \(\partial - \mathrm{i}G\)), and on the fold the canonically
  normalised four-dimensional field \(G^A_\mu(x)\) of its massless mode,
  \(g_3G^A_\mu(x) = G^A_\mu(\sigma = 0, x)\); \(G^A_{MN}\) and
  \(G^A_{\mu\nu}\) their field strengths. Unit: GeV (field), GeV\(^2\)
  (field strength). Status: AUTHOR.
\item
  \(A\): A colour (adjoint SU(3)) index running over \(1, \dots, 8\); a
  repeated index is summed. Status: STANDARD.
\item
  \(M\): A coordinate index of eight dimensions (also with a numbered
  copy \(M_1, M_2, \dots\)): it runs over \(x_0, x_1, \dots, x_7\); a
  repeated index is summed. Status: STANDARD.
\item
  \(N\): A coordinate index of eight dimensions (also
  \(N_1, N_2, \dots\)), running over \(x_0, \dots, x_7\); a repeated
  index is summed. Status: STANDARD.
\item
  \(W^I_M,\ W^I_\mu,\ W^I_{MN}\): The gauge field of \(SU(2)_L\)
  (\(I = 1, 2, 3\)): in the bulk \(W^I_M\) (geometric normalisation), on
  the fold the canonically normalised four-dimensional field
  \(W^I_\mu(x)\) of its massless mode,
  \(g_2W^I_\mu(x) = W^I_\mu(\sigma = 0, x)\); \(W^I_{MN}\),
  \(W^I_{\mu\nu}\) their field strengths;
  \(W^\pm_\mu = (W^1_\mu \mp \mathrm{i}W^2_\mu)/\sqrt2\). Unit: GeV
  (field), GeV\(^2\) (field strength). Status: AUTHOR.
\item
  \(I\): A weak-isospin (adjoint SU(2)) index running over \(1, 2, 3\);
  a repeated index is summed. Status: STANDARD.
\item
  \(B_M,\ B_\mu,\ B_{MN}\): The hypercharge gauge field of \(U(1)_Y\):
  in the bulk \(B_M\) (geometric normalisation), on the fold the
  canonically normalised four-dimensional field \(B_\mu(x)\) of its
  massless mode, \(g_1B_\mu(x) = B_\mu(\sigma = 0, x)\); \(B_{MN}\),
  \(B_{\mu\nu}\) their field strengths. Unit: GeV (field), GeV\(^2\)
  (field strength). Status: AUTHOR.
\item
  \(\partial\): A partial derivative: the derivative with respect to one
  variable at fixed values of the others (\(\partial_\sigma\) with
  respect to \(\sigma\); \(\partial_M\) and \(\partial_\mu\) with
  respect to the coordinate that carries the same index, of the eight
  dimensions or of the fold). Status: STANDARD.
\item
  \(f^{ABC}\): The structure constants of \(SU(3)\),
  \([t^A, t^B] = \mathrm{i}f^{ABC}t^C\). Unit: dimensionless. Status:
  STANDARD.
\item
  \(B\): A colour (adjoint SU(3)) index (as \(A\)). Status: STANDARD.
\item
  \(C\): A colour (adjoint SU(3)) index (as \(A\)). Status: STANDARD.
\item
  \(\epsilon^{IJK}\): The Levi-Civita symbol of three indices (the
  structure constants of \(SU(2)\)), \(\epsilon^{123} = 1\). Unit:
  dimensionless. Status: STANDARD.
\item
  \(J\): A weak-isospin (adjoint SU(2)) index (as \(I\)). Status:
  STANDARD.
\item
  \(K\): A weak-isospin (adjoint SU(2)) index (as \(I\)). Status:
  STANDARD.
\end{itemize}

The gauge fields of the Standard Model in the bulk of eight dimensions,
with the couplings \(\mathsf{g}_3\), \(\mathsf{g}_2\), \(\mathsf{g}_1\)
of eight dimensions. Their massless modes have the profile 1 along
\(\sigma\) and the kinetic weight
\(\mathrm{e}^{-3H\lvert\sigma\rvert}\), whose integral over the hidden
direction is \(2/(3H)\): they are four-dimensional gauge fields with the
couplings \(g_k^2 = 3H\mathsf{g}_k^2/(2V_*)\) (\(k = 1, 2, 3\)), the
measured couplings of the Standard Model, and they travel at the speed
of light at every depth (DERIVED, \texttt{um4\_bulk.py}, B2 and B3).
Their massive modes form a continuum from zero mass upwards; they modify
Coulomb's law at distances of the order of \(1/H\), add contact
interactions between the currents of the fold, and give the W and Z
bosons widths into the bulk, which bound \(H\) from below by about 2.5
TeV (DERIVED, ESTIMATE, \texttt{um4\_bulk.py}, B10, and
\texttt{um4\_propagators.py}). The gauge fields have no components along
the extra times and no dependence on them (ASSUMED). Status: AUTHOR (the
photon in the bulk, by the order), ASSUMED (all gauge fields in the
bulk), STANDARD (the Yang-Mills action).

\hypertarget{the-fold}{%
\subsection{The fold}\label{the-fold}}

\[
\begin{aligned}
&\int\mathrm{d}^8X\,\sqrt{\lvert g\rvert}\;\delta_{\mathrm{D}}(\sigma)\,\frac{1}{V_*}\,\mathbb{F}
 = \int\mathrm{d}^4x\,\sqrt{-\hat h}\;\mathcal{Z}\,\mathbb{F}\big\vert_{\sigma=0},\qquad \mathrm{d}^4x = \mathrm{d}x_1\,\mathrm{d}x_2\,\mathrm{d}x_3\,\mathrm{d}x_4,\\
&\hat h_{\mu\nu} = \mathrm{diag}\big(\mathrm{e}^{2a_4}, \mathrm{e}^{2a_4}, \mathrm{e}^{2a_4}, -1\big),\qquad
 \sqrt{-\hat h} = \mathrm{e}^{3a_4},\qquad
 \mathcal{Z} = \frac{V_{\mathrm{t}}\big(\mathrm{e}^{-a_4}b_4\big)^3}{V_*},\qquad V_{\mathrm{t}} = \int\mathrm{d}x_5\,\mathrm{d}x_6\,\mathrm{d}x_7
\end{aligned} \tag{IN-8}
\] \textbf{Symbols in (IN-8).}

\begin{itemize}
\tightlist
\item
  \(\mathrm{d}\): The differential: \(\mathrm{d}\) in front of a
  quantity is an infinitesimal change of it (as \(\mathrm{d}\sigma\), a
  small step of \(\sigma\)); the quotient of two such changes is a
  derivative; \(\mathrm{d}\) with an exponent, in front of several
  variables, is their volume element, the product of their
  differentials. Status: STANDARD.
\item
  \(X = \{x_0, \dots, x_7\}\): The author's coordinates of the
  eight-dimensional spacetime, \(X = \{x_0, x_1, \dots, x_7\}\) with the
  hidden direction first (Dirac\_Einstein-Lovelock\_4+4\_Nash.nb,
  section DEFCHART); \(X^M\) is the coordinate with the index \(M\) and
  \(\mathrm{d}^8X = \mathrm{d}x_0\,\mathrm{d}x_1\cdots\mathrm{d}x_7\)
  the coordinate volume element. Unit: length (each coordinate). Status:
  AUTHOR.
\item
  \(g_{MN}\): The metric tensor of the eight-dimensional spacetime
  (signature \(+,+,+,+,-,-,-,-\) in the order \(x_0, \dots, x_7\)):
  \(g_{MN}\,\mathrm{d}X^M\mathrm{d}X^N\) is the squared interval,
  \(g^{MN}\) its inverse, \(\lvert g\rvert\) the absolute value of its
  determinant (\(\sqrt{\lvert g\rvert}\,\mathrm{d}^8X\) is the proper
  volume element). Unit: dimensionless. Status: AUTHOR.
\item
  \(\delta_{\mathrm{D}}\): The Dirac delta function of \(\sigma\): zero
  for \(\sigma \ne 0\), with the integral 1 over any range that contains
  \(\sigma = 0\); the integral of \(\delta_{\mathrm{D}}(\sigma)\) times
  a smooth function of \(\sigma\) is the value of that function at
  \(\sigma = 0\). Unit: 1/length (for \(\delta_{\mathrm{D}}(\sigma)\)).
  Status: STANDARD.
\item
  \(\sigma\): The proper coordinate of the hidden direction:
  \(\mathrm{d}\sigma = \lvert\cot z\rvert\,\mathrm{d}x_0\),
  \(\sigma = \mathrm{sgn}(\pi/2 - z)\ln(\sin z)/(6H)\) on one period
  \(0 < z < \pi\); \(\sigma = 0\) is the fold \(z = \pi/2\),
  \(\sigma \to \mp\infty\) the tips \(z \to 0, \pi\); one period of
  \(x_0\) is the whole line \(-\infty < \sigma < \infty\). Unit: length.
  Status: DERIVED.
\item
  \(V_*\): The physical volume of the extra times \(x_5, x_6, x_7\) on
  the fold in the late universe,
  \(V_* = V_{\mathrm{t}}(\mathrm{e}^{-a_4}b_4)^3\) (with
  \(V_{\mathrm{t}} = \int\mathrm{d}x_5\,\mathrm{d}x_6\,\mathrm{d}x_7\)
  the coordinate volume of the extra times) at a time when
  \(\mathrm{e}^{-a_4}b_4\) is constant; it converts the
  eight-dimensional normalisation into the four-dimensional one. Unit:
  length\(^3\) (GeV\(^{-3}\)). Status: ASSUMED.
\item
  \(\mathbb{F}\): Any function of the fields at the fold (used to state
  a rule of integration). Status: STANDARD.
\item
  \(x\): The coordinates of the fold: \(x_1, x_2, x_3\) of ordinary
  space, \(x_4\) the time, \(x_5, x_6, x_7\) the extra times; a field of
  the fold is a function of the four observable coordinates
  \(x = (x_1, x_2, x_3, x_4)\), and
  \(\mathrm{d}^4x = \mathrm{d}x_1\,\mathrm{d}x_2\,\mathrm{d}x_3\,\mathrm{d}x_4\).
  Unit: length (each coordinate). Status: AUTHOR.
\item
  \(x_1\): The first coordinate of ordinary 3-space. Unit: length.
  Status: AUTHOR.
\item
  \(x_2\): The second coordinate of ordinary 3-space. Unit: length.
  Status: AUTHOR.
\item
  \(x_3\): The third coordinate of ordinary 3-space. Unit: length.
  Status: AUTHOR.
\item
  \(x_4\): The time; a prime on a function of \(x_4\) alone is
  \(\mathrm{d}/\mathrm{d}x_4\). Unit: length (= time, \(c = 1\)).
  Status: AUTHOR.
\item
  \(\hat h_{\mu\nu}\): The four-dimensional metric of the fold,
  \(\hat h_{\mu\nu} = \mathrm{diag}(\mathrm{e}^{2a_4}, \mathrm{e}^{2a_4}, \mathrm{e}^{2a_4}, -1)\)
  in the coordinates \(x_1, \dots, x_4\) (the induced metric at
  \(\sigma = 0\) without the extra times); \(\hat h^{\mu\nu}\) its
  inverse, \(\hat h\) its determinant,
  \(\sqrt{-\hat h} = \mathrm{e}^{3a_4}\); \(\hat h^{(1)}_{\mu\nu}\) its
  first-order cosmological perturbation. Unit: dimensionless. Status:
  DERIVED.
\item
  \(\mu\): A four-dimensional coordinate index running over the
  observable coordinates \(x_1, x_2, x_3, x_4\) of the fold (signature
  \(+,+,+,-\)); a repeated index is summed. Status: STANDARD.
\item
  \(\nu\): A four-dimensional coordinate index (as \(\mu\)). Status:
  STANDARD.
\item
  \(\mathrm{diag}\): A diagonal matrix with the listed entries on its
  diagonal (in the order of the coordinates \(x_0, x_1, \dots, x_7\) for
  a tensor of eight dimensions, \(x_1, \dots, x_7\) for a tensor of the
  fold, \(x_1, \dots, x_4\) for a four-dimensional tensor). Status:
  STANDARD.
\item
  \(\mathrm{e}\): Euler's number, the base of the natural logarithm;
  \(\mathrm{e}^{X}\) is the exponential of \(X\). Unit: dimensionless.
  Value: 2.718281828. Status: STANDARD.
\item
  \(a_4\): The author's metric function of the time \(x_4\) that scales
  ordinary 3-space: on the fold \(\mathrm{e}^{a_4}\) is the scale factor
  of cosmology and \(a_4'\) the Hubble rate. Unit: dimensionless.
  Status: AUTHOR.
\item
  \(\mathcal{Z}\): The fold factor,
  \(\mathcal{Z}(x_4) = V_{\mathrm{t}}(\mathrm{e}^{-a_4}b_4)^3/V_*\): the
  factor by which the size of the extra times multiplies every term of
  the fold and the kinetic terms of the massless modes of the fields of
  the bulk; \(\mathcal{Z} = 1\) when the extra times are frozen at their
  late value. Unit: dimensionless. Value: 1 in the late universe.
  Status: DERIVED.
\item
  \(V_{\mathrm{t}}\): The coordinate volume of the (compact) extra
  times,
  \(V_{\mathrm{t}} = \int\mathrm{d}x_5\,\mathrm{d}x_6\,\mathrm{d}x_7\).
  Unit: length\(^3\). Status: ASSUMED.
\item
  \(b_4\): The author's metric function of the time \(x_4\) in the
  extra-time components, \(b_4 > 0\), never identically 1 (the author's
  order of 2026-10-24); the scale factor (the factor of the lengths) of
  the extra times is \(\mathrm{e}^{-a_4}b_4\) on the fold and
  \(\mathrm{e}^{-a_4}b_4\,\mathrm{e}^{-H\lvert\sigma\rvert}\) at the
  distance \(\sigma\) from it (the square root of the warp factor times
  \(\mathrm{e}^{-a_4}b_4\)). Unit: dimensionless. Status: AUTHOR.
\item
  \(x_5\): The first of the three extra time-like coordinates. Unit:
  length. Status: AUTHOR.
\item
  \(x_6\): The second of the three extra time-like coordinates. Unit:
  length. Status: AUTHOR.
\item
  \(x_7\): The third of the three extra time-like coordinates. Unit:
  length. Status: AUTHOR.
\end{itemize}

\[
\mathcal{L}_{\mathrm{IG}} = \frac{\hat R}{2\kappa_{\mathrm{ind}}} \tag{IN-9}
\] \textbf{Symbols in (IN-9).}

\begin{itemize}
\tightlist
\item
  \(\mathcal{L}_{\mathrm{IG}}\): The Lagrangian of the Einstein-Hilbert
  term of the four-dimensional metric of the fold that the quantum
  fluctuations of the fields of the fold induce; it adds to the
  Einstein-Hilbert term that the graviton of the bulk gives the fold and
  does not change the speed of gravitational waves. Unit: GeV\(^4\).
  Status: ESTIMATE.
\item
  \(\hat R\): The curvature of the metric \(\hat g_{ab}\) of the fold
  (FD-2): with four frame indices its Riemann tensor, with two its Ricci
  tensor \(\hat R_{ab}\), without indices its scalar curvature
  \(\hat R\); for frozen extra times the scalar curvature equals that of
  the four-dimensional metric \(\hat h_{\mu\nu}\). Unit: 1/length\(^2\).
  Status: STANDARD.
\item
  \(\kappa_{\mathrm{ind}}\): The coupling of the Einstein-Hilbert term
  \(\hat R/(2\kappa_{\mathrm{ind}})\) that the quantum fluctuations of
  the fields of the fold may induce,
  \(\kappa_{\mathrm{ind}} = 8\pi G_{\mathrm{ind}}\); it adds to the
  gravity of the fold,
  \(1/(8\pi G_{\mathrm{N}}) = V_*/\kappa_7 + 1/\kappa_{\mathrm{ind}}\).
  Unit: GeV\(^{-2}\). Value: not fixed by UM-4; negligible for a cutoff
  of the order of \(H\). Status: ESTIMATE.
\end{itemize}

\[
\mathcal{T}_{\mathrm{ff}}{}^a{}_b = \mathrm{diag}\big(p_{\mathrm{ff}}, p_{\mathrm{ff}}, p_{\mathrm{ff}}, -\rho_{\mathrm{ff}},
 q_{\mathrm{ff}}, q_{\mathrm{ff}}, q_{\mathrm{ff}}\big) \tag{IN-10}
\] \textbf{Symbols in (IN-10).}

\begin{itemize}
\tightlist
\item
  \(\mathcal{T}_{\mathrm{ff}}{}^a{}_b\): The energy-momentum tensor of
  the fold fluid (frame components of the fold, four-dimensional
  normalisation). Unit: GeV\(^4\). Status: ASSUMED.
\item
  \(a\): A frame (orthonormal) index running over \(0, 1, \dots, 7\)
  (frame direction \((0)\) along \(\sigma\), \((1), (2), (3)\) along
  ordinary space, \((4)\) along the time, \((5), (6), (7)\) along the
  extra times); raised and lowered with the frame metric
  \(\eta_{ab} = \mathrm{diag}(1, 1, 1, 1, -1, -1, -1, -1)\); on the fold
  (\(a, b \in \{1, \dots, 7\}\)) the frame of the fold; in the tensors
  of the fold written with its coordinates \(x^a\) (its metric
  \(\hat g_{ab}\) and the tensors built from it) a coordinate index of
  the fold, \(a \in \{1, \dots, 7\}\); a repeated index is summed.
  Status: STANDARD.
\item
  \(b\): A frame index (as \(a\)). Status: STANDARD.
\item
  \(\mathrm{diag}\): A diagonal matrix with the listed entries on its
  diagonal (in the order of the coordinates \(x_0, x_1, \dots, x_7\) for
  a tensor of eight dimensions, \(x_1, \dots, x_7\) for a tensor of the
  fold, \(x_1, \dots, x_4\) for a four-dimensional tensor). Status:
  STANDARD.
\item
  \(p_{\mathrm{ff}}\): The pressure of the fold fluid along ordinary
  space. Unit: GeV\(^4\). Value: \(-\rho_{\mathrm{ff}}\) (a tension).
  Status: ASSUMED (the fluid), DERIVED (the value).
\item
  \(\rho_{\mathrm{ff}}\): The energy density of the fold fluid. Unit:
  GeV\(^4\). Value:
  \(\lambda_{\mathrm{F}}V_* + \rho_\Lambda - \rho_{\mathrm{vac}}\):
  together with the vacuum energy \(\rho_{\mathrm{vac}}\) of the fields
  of the fold it makes up the tension \(\lambda_{\mathrm{F}}V_*\) that
  the static block requires plus the dark energy \(\rho_\Lambda\).
  Status: ASSUMED (the fluid), DERIVED (the value).
\item
  \(q_{\mathrm{ff}}\): The pressure of the fold fluid along the extra
  times: minus its energy density (its tension) plus the pressure
  \(q_{\mathrm{hold}}\) that holds the volume of the extra times fixed.
  Unit: GeV\(^4\). Value: \(-\rho_{\mathrm{ff}} + q_{\mathrm{hold}}\).
  Status: ASSUMED (the fluid), DERIVED (the value).
\end{itemize}

The fold: the reduction of the terms of the fold to four-dimensional
fields with the volume \(V_*\) of the extra times and the factor
\(\mathcal{Z}\) (IN-8); the Einstein-Hilbert term that the fields of the
fold induce (IN-9), negligible for a cutoff of the order of \(H\)
because the graviton of the bulk gives the fold its gravity: to the
first order in its curvature over \(H^2\) the fold obeys the Einstein
equations of seven dimensions with \(\kappa_7 = 5\kappa H/2\), which
with the volume of the extra times held fixed are the Friedmann
equations of general relativity with Newton's constant
\(8\pi G_{\mathrm{N}} = \kappa_7/V_*\), and its four-dimensional
gravitational waves travel at the speed of light at every wavelength;
and the fold fluid (IN-10, with the values FE-65), the thin source that
the kink of the warp factor requires: a pure tension
\(\lambda_{\mathrm{F}}\) at the static block,
\(\kappa\lambda_{\mathrm{F}} = 12H\), as for the brane of Randall and
Sundrum, whose excess over \(\lambda_{\mathrm{F}}\) (together with the
vacuum energy of the fields) is the dark energy, and the pressure
\(q_{\mathrm{hold}}\) along the extra times that holds their volume
fixed (without it Newton's constant, the squares of the gauge and Yukawa
couplings and the quartic Higgs coupling would change far beyond the
observed bounds). Status: ASSUMED (the fold fluid and its holding
pressure), DERIVED (the reduction, the gravity of the fold, the values
of the fold fluid), ESTIMATE (the induced term).

\[
\begin{aligned}
&p_{\mathrm{ff}} = -\rho_{\mathrm{ff}},\qquad q_{\mathrm{ff}} = -\rho_{\mathrm{ff}} + q_{\mathrm{hold}},\qquad
 \kappa\lambda_{\mathrm{F}} = 12H\ \ \big(\lambda_{\mathrm{F}}V_* = 30H^2M_{\mathrm{red}}^2\big),\\
&q_{\mathrm{hold}} = -\frac{\rho - 3\mathcal{P}}{2} - q_{\mathrm{SM}} - q_\chi - q_{\mathrm{P}}:\qquad
 \text{radiation } 0,\quad \text{dust } -\frac{\rho}{2},\quad \text{a vacuum energy } -\rho
\end{aligned} \tag{FE-65}
\] \textbf{Symbols in (FE-65).}

\begin{itemize}
\tightlist
\item
  \(p_{\mathrm{ff}}\): The pressure of the fold fluid along ordinary
  space. Unit: GeV\(^4\). Value: \(-\rho_{\mathrm{ff}}\) (a tension).
  Status: ASSUMED (the fluid), DERIVED (the value).
\item
  \(\rho_{\mathrm{ff}}\): The energy density of the fold fluid. Unit:
  GeV\(^4\). Value:
  \(\lambda_{\mathrm{F}}V_* + \rho_\Lambda - \rho_{\mathrm{vac}}\):
  together with the vacuum energy \(\rho_{\mathrm{vac}}\) of the fields
  of the fold it makes up the tension \(\lambda_{\mathrm{F}}V_*\) that
  the static block requires plus the dark energy \(\rho_\Lambda\).
  Status: ASSUMED (the fluid), DERIVED (the value).
\item
  \(q_{\mathrm{ff}}\): The pressure of the fold fluid along the extra
  times: minus its energy density (its tension) plus the pressure
  \(q_{\mathrm{hold}}\) that holds the volume of the extra times fixed.
  Unit: GeV\(^4\). Value: \(-\rho_{\mathrm{ff}} + q_{\mathrm{hold}}\).
  Status: ASSUMED (the fluid), DERIVED (the value).
\item
  \(q_{\mathrm{hold}}\): The pressure along the extra times that the
  fold fluid adds to its tension to hold the volume of the extra times
  fixed (four-dimensional normalisation). Unit: GeV\(^4\). Value:
  \(-(\rho - 3\mathcal{P})/2 - q_{\mathrm{SM}} - q_\chi - q_{\mathrm{P}}\):
  0 for radiation, \(-\rho/2\) for dust, \(-\rho\) for a vacuum energy;
  today \(-0.842\,\rho_{\mathrm{c}}^{(0)}\) if the baryons have no
  pressure along the extra times (ASSUMED). Status: DERIVED (the value),
  ASSUMED (that the fold fluid supplies it).
\item
  \(\kappa\): The eight-dimensional gravitational coupling,
  \(\kappa = 8\pi G_8\). Unit: GeV\(^{-6}\). Status: AUTHOR.
\item
  \(\lambda_{\mathrm{F}}\): The tension that the static block requires
  of the fold, \(\kappa\lambda_{\mathrm{F}} = 12H\): the total vacuum
  energy of the seven-dimensional fold, made of the tension of the fold
  fluid and the vacuum energy of the fields of the fold (both have the
  pressure minus their energy density along ordinary space and along the
  extra times), per unit volume of the six other directions of the fold
  (eight-dimensional normalisation); in the four-dimensional
  normalisation it is \(\lambda_{\mathrm{F}}V_*\), and the excess of the
  total over \(\lambda_{\mathrm{F}}V_*\) is the dark energy
  \(\rho_\Lambda\). Unit: GeV\(^7\). Value: \(12H/\kappa\);
  \(\lambda_{\mathrm{F}}V_* = 30H^2M_{\mathrm{red}}^2 = 1.1\times10^{45}\)
  GeV\(^4\) for \(H = 2.5\) TeV. Status: DERIVED.
\item
  \(H\): The author's constant (an inverse length) that fixes the period
  \(\pi/(6H)\) of \(x_0\), the warp factor
  \(\mathrm{e}^{-2H\lvert\sigma\rvert}\) of the metric, the cosmological
  constant \(\Lambda = -21H^2\) and the tension
  \(\kappa\lambda_{\mathrm{F}} = 12H\) of the fold; in UM-4 a free
  parameter, bounded from below by the propagation of the gauge fields
  in the bulk. Unit: GeV (1/length). Value: \(H \gtrsim 2.5\) TeV (from
  the invisible width of the Z boson). Status: AUTHOR.
\item
  \(V_*\): The physical volume of the extra times \(x_5, x_6, x_7\) on
  the fold in the late universe,
  \(V_* = V_{\mathrm{t}}(\mathrm{e}^{-a_4}b_4)^3\) (with
  \(V_{\mathrm{t}} = \int\mathrm{d}x_5\,\mathrm{d}x_6\,\mathrm{d}x_7\)
  the coordinate volume of the extra times) at a time when
  \(\mathrm{e}^{-a_4}b_4\) is constant; it converts the
  eight-dimensional normalisation into the four-dimensional one. Unit:
  length\(^3\) (GeV\(^{-3}\)). Status: ASSUMED.
\item
  \(M_{\mathrm{red}}\): The reduced Planck mass,
  \(M_{\mathrm{red}} = (8\pi G_{\mathrm{N}})^{-1/2}\). Unit: GeV. Value:
  \(2.435\times10^{18}\) GeV. Status: STANDARD.
\item
  \(\rho\): The total energy density of the matter of the fold (all
  species); \(\Delta\rho\) its perturbation. Unit: GeV\(^4\). Status:
  STANDARD.
\item
  \(\mathcal{P}\): The total pressure of the matter of the fold;
  \(\Delta\mathcal{P}\) its perturbation. Unit: GeV\(^4\). Status:
  STANDARD.
\item
  \(q_{\mathrm{SM}}\): The pressure of the Standard Model along the
  extra times at the fold: the mean on-shell Lagrangian density of the
  Standard-Model fields (the extra-time component of the energy-momentum
  tensor of every term of the fold is its Lagrangian, because the fields
  of the fold do not depend on the extra times). Unit: GeV\(^4\).
  Status: DERIVED (for radiation and free fields), ASSUMED (for the
  baryons).
\item
  \(q_\chi\): The pressure of the particles \(\chi\) of dirac16complex
  along the extra times, integrated across the fold (their on-shell
  Lagrangian, the self-interaction energy per volume; zero for a dilute
  gas). Unit: GeV\(^4\). Status: DERIVED.
\item
  \(q_{\mathrm{P}}\): The pressure of the portals along the extra times:
  their mean on-shell Lagrangian density. Unit: GeV\(^4\). Status:
  DERIVED.
\end{itemize}

The values of the fold fluid with frozen extra times: its energy density
and its pressure along ordinary space are those of a tension; the
tension that the static block requires is
\(\kappa\lambda_{\mathrm{F}} = 12H\), in the four-dimensional
normalisation \(\lambda_{\mathrm{F}}V_* = 30H^2M_{\mathrm{red}}^2\); the
pressure along the extra times adds \(q_{\mathrm{hold}}\), which makes
up what the matter of the fold lacks of the pressure
\(-(\rho - 3\mathcal{P})/2\) that the radion condition requires (today
\(-0.842\) times the critical density if the baryons have no pressure
along the extra times, ASSUMED). The observed dark energy needs the
tension tuned to one part in \(10^{92}\) for \(H = 2.5\) TeV, the
cosmological-constant problem of the brane of Randall and Sundrum.
Status: DERIVED (the values), ASSUMED (that the fold fluid supplies
them; the pressure of the baryons), OPEN (its origin, the tuning).

\hypertarget{the-higgs-doublet-with-the-ad-hoc-vacuum}{%
\subsection{The Higgs doublet with the ad hoc
vacuum}\label{the-higgs-doublet-with-the-ad-hoc-vacuum}}

\[
\mathcal{L}_{\mathrm{H}} = -\hat h^{\mu\nu}\big(D_\mu\Phi\big)^\dagger D_\nu\Phi - V(\Phi),\qquad
V(\Phi) = \mu^2\,\Phi^\dagger\Phi + \lambda\big(\Phi^\dagger\Phi\big)^2,\qquad \lambda > 0,\quad \mu^2 < 0 \tag{IN-11}
\] \textbf{Symbols in (IN-11).}

\begin{itemize}
\tightlist
\item
  \(\mathcal{L}_{\mathrm{H}}\): The Lagrangian of the Higgs doublet on
  the fold. Unit: GeV\(^4\). Status: AUTHOR.
\item
  \(\hat h_{\mu\nu}\): The four-dimensional metric of the fold,
  \(\hat h_{\mu\nu} = \mathrm{diag}(\mathrm{e}^{2a_4}, \mathrm{e}^{2a_4}, \mathrm{e}^{2a_4}, -1)\)
  in the coordinates \(x_1, \dots, x_4\) (the induced metric at
  \(\sigma = 0\) without the extra times); \(\hat h^{\mu\nu}\) its
  inverse, \(\hat h\) its determinant,
  \(\sqrt{-\hat h} = \mathrm{e}^{3a_4}\); \(\hat h^{(1)}_{\mu\nu}\) its
  first-order cosmological perturbation. Unit: dimensionless. Status:
  DERIVED.
\item
  \(\mu\): A four-dimensional coordinate index running over the
  observable coordinates \(x_1, x_2, x_3, x_4\) of the fold (signature
  \(+,+,+,-\)); a repeated index is summed. Status: STANDARD.
\item
  \(\nu\): A four-dimensional coordinate index (as \(\mu\)). Status:
  STANDARD.
\item
  \(D_M,\ D_\mu\): The covariant derivative: on dirac16complex
  \(D_M = \partial_M + \Omega_M\) (with
  \(D_M\bar\Psi = \partial_M\bar\Psi - \bar\Psi\Omega_M\)); on the
  fields of the fold the gauge and spin covariant derivative \(D_\mu\)
  (the partial derivative plus the spin connection of the fold and the
  gauge fields times their couplings and generators);
  \(D^\mu = \hat h^{\mu\nu}D_\nu\); on an adjoint field
  \((D_\mu X)^A = \partial_\mu X^A + g_3f^{ABC}G^B_\mu X^C\) (likewise
  for \(SU(2)\) with \(g_2\epsilon^{IJK}\)). Unit: 1/length. Status:
  STANDARD.
\item
  \(\Phi\): The Higgs doublet (two complex components, hypercharge
  \(1/2\)), a field of the fold. Unit: GeV. Status: AUTHOR.
\item
  \({}^\dagger\): Hermitian conjugation: the complex conjugate of the
  transpose of a matrix or spinor; for an operator on the space of
  states, its adjoint with respect to the scalar product of that space
  (for the quantised dirac16complex the indefinite Krein scalar
  product). Status: STANDARD.
\item
  \(V(\Phi)\): The Higgs potential,
  \(V = \mu^2\Phi^\dagger\Phi + \lambda(\Phi^\dagger\Phi)^2\). Unit:
  GeV\(^4\). Status: AUTHOR.
\item
  \(\mu^2\): The mass parameter of the Higgs potential, \(\mu^2 < 0\)
  (it enters only as \(\mu^2\)). Unit: GeV\(^2\). Value:
  \(\mu^2 = -\lambda v^2 = -m_h^2/2 = -(88.6\ \mathrm{GeV})^2\). Status:
  AUTHOR.
\item
  \(\lambda\): The quartic coupling of the Higgs potential,
  \(\lambda > 0\). Unit: dimensionless. Value:
  \(\lambda = m_h^2/(2v^2) = 0.129\). Status: AUTHOR.
\end{itemize}

\[
\Phi = \frac{1}{\sqrt2}\begin{pmatrix}\mathcal{A}(x_4)\\ v + h\end{pmatrix},\qquad
\mathcal{A}(x_4) = v_{\mathcal{A}}\,\mathrm{e}^{-2a_{40}x_4},\qquad a_{40} > 0,\quad x_4 \ge 0,\qquad v^2 = -\frac{\mu^2}{\lambda} \tag{IN-12}
\] \textbf{Symbols in (IN-12).}

\begin{itemize}
\tightlist
\item
  \(\Phi\): The Higgs doublet (two complex components, hypercharge
  \(1/2\)), a field of the fold. Unit: GeV. Status: AUTHOR.
\item
  \(\mathcal{A}\): The author's ad hoc upper component of the Higgs
  vacuum, almightyAdHoc,
  \(\mathcal{A}(x_4) = v_{\mathcal{A}}\mathrm{e}^{-2a_{40}x_4}\). Unit:
  GeV. Status: AUTHOR.
\item
  \(x_4\): The time; a prime on a function of \(x_4\) alone is
  \(\mathrm{d}/\mathrm{d}x_4\). Unit: length (= time, \(c = 1\)).
  Status: AUTHOR.
\item
  \(v\): The vacuum expectation value of the lower component of the
  Higgs doublet, \(v^2 = -\mu^2/\lambda\). Unit: GeV. Value: 246.22 GeV
  \(= (\sqrt2\,G_{\mathrm{F}})^{-1/2}\) with the Fermi constant
  \(G_{\mathrm{F}} = 1.1663788\times10^{-5}\) GeV\(^{-2}\) of the muon
  lifetime (Particle Data Group). Status: STANDARD.
\item
  \(h\): The Higgs field: the fluctuation of the lower component of the
  Higgs doublet about \(v\). Unit: GeV. Value: \(m_h = 125.25\) GeV.
  Status: STANDARD.
\item
  \(v_{\mathcal{A}}\): The amplitude of almightyAdHoc at \(x_4 = 0\);
  the author's almightyAdHoc, \(\mathrm{e}^{-2a_{40}x_4}\) (written
  Exp{[}-2 a40 x4{]} in his notebook), has \(v_{\mathcal{A}} = 1\) in
  the units of his notebook. Unit: GeV. Value: 1 in the units of the
  author's notebook; UM-4 keeps it free (the examples take
  \(v_{\mathcal{A}} = v\) and \(10v\)). Status: AUTHOR (the form),
  ASSUMED (a free amplitude).
\item
  \(\mathrm{e}\): Euler's number, the base of the natural logarithm;
  \(\mathrm{e}^{X}\) is the exponential of \(X\). Unit: dimensionless.
  Value: 2.718281828. Status: STANDARD.
\item
  \(a_{40}\): The decay constant of almightyAdHoc, \(a_{40} > 0\):
  \(\mathcal{A}\) falls by \(\mathrm{e}\) in the time \(1/(2a_{40})\).
  Unit: GeV (1/length). Status: AUTHOR.
\item
  \(\mu^2\): The mass parameter of the Higgs potential, \(\mu^2 < 0\)
  (it enters only as \(\mu^2\)). Unit: GeV\(^2\). Value:
  \(\mu^2 = -\lambda v^2 = -m_h^2/2 = -(88.6\ \mathrm{GeV})^2\). Status:
  AUTHOR.
\item
  \(\lambda\): The quartic coupling of the Higgs potential,
  \(\lambda > 0\). Unit: dimensionless. Value:
  \(\lambda = m_h^2/(2v^2) = 0.129\). Status: AUTHOR.
\end{itemize}

The Higgs doublet with the author's potential, expanded about the
author's ad hoc vacuum, whose upper component almightyAdHoc decays as
\(\mathrm{e}^{-2a_{40}x_4}\) (the author writes
\(\mathrm{Exp}[-2a_{40}x_4]\); UM-4 gives it the amplitude
\(v_{\mathcal{A}}\) so that it has the dimension of the Higgs field).
Status: AUTHOR.

\hypertarget{the-gauge-fields-at-the-fold-the-quarks-and-the-leptons}{%
\subsection{The gauge fields at the fold, the quarks and the
leptons}\label{the-gauge-fields-at-the-fold-the-quarks-and-the-leptons}}

\[
\begin{aligned}
D_\mu &= \partial_\mu + \hat\omega_\mu - \mathrm{i}g_3\,G^A_\mu\,t^A - \mathrm{i}g_2\,W^I_\mu\,\frac{\tau^I}{2}
 - \mathrm{i}g_1\,Y_f\,B_\mu,\\
g_3\,G^A_\mu(x) &= G^A_\mu(\sigma = 0, x),\qquad g_2\,W^I_\mu(x) = W^I_\mu(\sigma = 0, x),\qquad g_1\,B_\mu(x) = B_\mu(\sigma = 0, x)
\end{aligned} \tag{IN-13}
\] \textbf{Symbols in (IN-13).}

\begin{itemize}
\tightlist
\item
  \(D_M,\ D_\mu\): The covariant derivative: on dirac16complex
  \(D_M = \partial_M + \Omega_M\) (with
  \(D_M\bar\Psi = \partial_M\bar\Psi - \bar\Psi\Omega_M\)); on the
  fields of the fold the gauge and spin covariant derivative \(D_\mu\)
  (the partial derivative plus the spin connection of the fold and the
  gauge fields times their couplings and generators);
  \(D^\mu = \hat h^{\mu\nu}D_\nu\); on an adjoint field
  \((D_\mu X)^A = \partial_\mu X^A + g_3f^{ABC}G^B_\mu X^C\) (likewise
  for \(SU(2)\) with \(g_2\epsilon^{IJK}\)). Unit: 1/length. Status:
  STANDARD.
\item
  \(\mu\): A four-dimensional coordinate index running over the
  observable coordinates \(x_1, x_2, x_3, x_4\) of the fold (signature
  \(+,+,+,-\)); a repeated index is summed. Status: STANDARD.
\item
  \(\partial\): A partial derivative: the derivative with respect to one
  variable at fixed values of the others (\(\partial_\sigma\) with
  respect to \(\sigma\); \(\partial_M\) and \(\partial_\mu\) with
  respect to the coordinate that carries the same index, of the eight
  dimensions or of the fold). Status: STANDARD.
\item
  \(\hat\omega_\mu\): The spin connection of the four-dimensional metric
  \(\hat h_{\mu\nu}\) of the fold acting on four-dimensional spinors.
  Unit: 1/length. Status: STANDARD.
\item
  \(\mathrm{i}\): The imaginary unit, \(\mathrm{i}^2 = -1\). Unit:
  dimensionless. Status: STANDARD.
\item
  \(g_3\): The \(SU(3)_C\) gauge coupling of the four-dimensional world
  (the strong coupling \(g_s\)): the coupling of the massless mode of
  the gluon field of the bulk, \(g_3^2 = 3H\mathsf{g}_3^2/(2V_*)\).
  Unit: dimensionless. Value: \(g_3(m_Z) = 1.2177\)
  (\(g_3^2/(4\pi) = 0.1180\)). Status: STANDARD.
\item
  \(G^A_M,\ G^A_\mu,\ G^A_{MN}\): The gluon field of \(SU(3)_C\)
  (\(A = 1, \dots, 8\)): in the bulk \(G^A_M\), a field of eight
  dimensions in its geometric normalisation (the covariant derivative
  \(\partial - \mathrm{i}G\)), and on the fold the canonically
  normalised four-dimensional field \(G^A_\mu(x)\) of its massless mode,
  \(g_3G^A_\mu(x) = G^A_\mu(\sigma = 0, x)\); \(G^A_{MN}\) and
  \(G^A_{\mu\nu}\) their field strengths. Unit: GeV (field), GeV\(^2\)
  (field strength). Status: AUTHOR.
\item
  \(A\): A colour (adjoint SU(3)) index running over \(1, \dots, 8\); a
  repeated index is summed. Status: STANDARD.
\item
  \(t^A\): The generators of \(SU(3)\) on colour triplets, \(t^A\): half
  the Gell-Mann matrices, with \(\mathrm{tr}(t^At^B) = \delta^{AB}/2\).
  Unit: dimensionless. Status: STANDARD.
\item
  \(g_2\): The \(SU(2)_L\) gauge coupling of the four-dimensional world
  (usually \(g\)): the coupling of the massless mode of the \(SU(2)_L\)
  field of the bulk, \(g_2^2 = 3H\mathsf{g}_2^2/(2V_*)\). Unit:
  dimensionless. Value: \(g_2(m_Z) = 0.6517\). Status: STANDARD.
\item
  \(W^I_M,\ W^I_\mu,\ W^I_{MN}\): The gauge field of \(SU(2)_L\)
  (\(I = 1, 2, 3\)): in the bulk \(W^I_M\) (geometric normalisation), on
  the fold the canonically normalised four-dimensional field
  \(W^I_\mu(x)\) of its massless mode,
  \(g_2W^I_\mu(x) = W^I_\mu(\sigma = 0, x)\); \(W^I_{MN}\),
  \(W^I_{\mu\nu}\) their field strengths;
  \(W^\pm_\mu = (W^1_\mu \mp \mathrm{i}W^2_\mu)/\sqrt2\). Unit: GeV
  (field), GeV\(^2\) (field strength). Status: AUTHOR.
\item
  \(I\): A weak-isospin (adjoint SU(2)) index running over \(1, 2, 3\);
  a repeated index is summed. Status: STANDARD.
\item
  \(\tau^I\): The Pauli matrices acting on weak doublets
  (\(\tau^1, \tau^2, \tau^3\); \(\tau^2\) is the second one, not a
  square). Unit: dimensionless. Status: STANDARD.
\item
  \(g_1\): The \(U(1)_Y\) gauge coupling of the four-dimensional world
  (usually \(g'\); not normalised for grand unification): the coupling
  of the massless mode of the hypercharge field of the bulk,
  \(g_1^2 = 3H\mathsf{g}_1^2/(2V_*)\). Unit: dimensionless. Value:
  \(g_1(m_Z) = 0.3574\). Status: STANDARD.
\item
  \(Y_f\): The hypercharge of a field, the charge of its coupling
  \(g_1Y_fB_\mu\) to the hypercharge field in the covariant derivative
  (\(1/6\), \(2/3\), \(-1/3\), \(-1/2\), \(-1\), \(0\) for
  \(Q, u_{\mathrm{R}}, d_{\mathrm{R}}, \ell, e_{\mathrm{R}}, \nu_{\mathrm{R}}\),
  \(1/2\) for \(\Phi\)). Unit: dimensionless. Status: STANDARD.
\item
  \(f\): The label of a kind of Standard-Model multiplet:
  \(f \in \{Q, u, d, \ell, e, \nu\}\) (quark doublets, up-quark
  singlets, down-quark singlets, lepton doublets, charged-lepton
  singlets, right-handed neutrinos); a sum over \(f\) runs over the six.
  Status: STANDARD.
\item
  \(B_M,\ B_\mu,\ B_{MN}\): The hypercharge gauge field of \(U(1)_Y\):
  in the bulk \(B_M\) (geometric normalisation), on the fold the
  canonically normalised four-dimensional field \(B_\mu(x)\) of its
  massless mode, \(g_1B_\mu(x) = B_\mu(\sigma = 0, x)\); \(B_{MN}\),
  \(B_{\mu\nu}\) their field strengths. Unit: GeV (field), GeV\(^2\)
  (field strength). Status: AUTHOR.
\item
  \(\sigma\): The proper coordinate of the hidden direction:
  \(\mathrm{d}\sigma = \lvert\cot z\rvert\,\mathrm{d}x_0\),
  \(\sigma = \mathrm{sgn}(\pi/2 - z)\ln(\sin z)/(6H)\) on one period
  \(0 < z < \pi\); \(\sigma = 0\) is the fold \(z = \pi/2\),
  \(\sigma \to \mp\infty\) the tips \(z \to 0, \pi\); one period of
  \(x_0\) is the whole line \(-\infty < \sigma < \infty\). Unit: length.
  Status: DERIVED.
\item
  \(x\): The coordinates of the fold: \(x_1, x_2, x_3\) of ordinary
  space, \(x_4\) the time, \(x_5, x_6, x_7\) the extra times; a field of
  the fold is a function of the four observable coordinates
  \(x = (x_1, x_2, x_3, x_4)\), and
  \(\mathrm{d}^4x = \mathrm{d}x_1\,\mathrm{d}x_2\,\mathrm{d}x_3\,\mathrm{d}x_4\).
  Unit: length (each coordinate). Status: AUTHOR.
\end{itemize}

\[
\begin{aligned}
&\mathcal{L}_{\mathrm{f}} = \sum_{\psi}\mathrm{i}\,\bar\psi\,\hat\gamma^\mu D_\mu\psi,\qquad
 \psi \in \{Q, u_{\mathrm{R}}, d_{\mathrm{R}}, \ell, e_{\mathrm{R}}, \nu_{\mathrm{R}}\},\\
&\{\hat\gamma^\mu, \hat\gamma^\nu\} = -2\hat h^{\mu\nu}\mathbf{1},\qquad \bar\psi = \psi^\dagger\beta,\qquad
 \beta = \hat\gamma^{(4)},\qquad \hat\gamma_5 = \mathrm{i}\,\hat\gamma^{(4)}\hat\gamma^{(1)}\hat\gamma^{(2)}\hat\gamma^{(3)},\\
&\hat P_{\mathrm{L}} = \frac{\mathbf{1} - \hat\gamma_5}{2},\quad \hat P_{\mathrm{R}} = \frac{\mathbf{1} + \hat\gamma_5}{2},\qquad
 \hat P_{\mathrm{L}}Q = Q,\quad \hat P_{\mathrm{L}}\ell = \ell,\quad \hat P_{\mathrm{R}}u_{\mathrm{R}} = u_{\mathrm{R}},\quad
 \hat P_{\mathrm{R}}d_{\mathrm{R}} = d_{\mathrm{R}},\quad \hat P_{\mathrm{R}}e_{\mathrm{R}} = e_{\mathrm{R}},\quad
 \hat P_{\mathrm{R}}\nu_{\mathrm{R}} = \nu_{\mathrm{R}}
\end{aligned} \tag{IN-14}
\] \textbf{Symbols in (IN-14).}

\begin{itemize}
\tightlist
\item
  \(\mathcal{L}_{\mathrm{f}}\): The kinetic Lagrangian of the quarks and
  leptons on the fold. Unit: GeV\(^4\). Status: STANDARD.
\item
  \(\psi\): The label of a Standard-Model fermion field in a sum:
  \(\psi\) runs over
  \(Q, u_{\mathrm{R}}, d_{\mathrm{R}}, \ell, e_{\mathrm{R}}, \nu_{\mathrm{R}}\).
  Status: STANDARD.
\item
  \(\mathrm{i}\): The imaginary unit, \(\mathrm{i}^2 = -1\). Unit:
  dimensionless. Status: STANDARD.
\item
  \(\psi,\ \bar\psi\): A Standard-Model fermion field (a
  four-dimensional Weyl field of the fold with three generations,
  written as a four-component spinor of fixed chirality);
  \(\bar\psi = \psi^\dagger\beta\); \(\psi_f\) the fields of the kind
  \(f\). Unit: GeV\(^{3/2}\). Status: STANDARD.
\item
  \(\hat\gamma^\mu\): The four-dimensional complex \(4 \times 4\) Dirac
  matrices of the fold: \(\hat\gamma^{(a)}\) along the frame directions
  of the fold and \(\hat\gamma^\mu\) along its coordinates (the frame
  matrices contracted with the inverse frame field of the fold),
  \(\{\hat\gamma^\mu, \hat\gamma^\nu\} = -2\hat h^{\mu\nu}\mathbf{1}\);
  \(\hat\gamma_\mu = \hat h_{\mu\nu}\hat\gamma^\nu\). Unit:
  dimensionless. Status: STANDARD.
\item
  \(\mu\): A four-dimensional coordinate index running over the
  observable coordinates \(x_1, x_2, x_3, x_4\) of the fold (signature
  \(+,+,+,-\)); a repeated index is summed. Status: STANDARD.
\item
  \(D_M,\ D_\mu\): The covariant derivative: on dirac16complex
  \(D_M = \partial_M + \Omega_M\) (with
  \(D_M\bar\Psi = \partial_M\bar\Psi - \bar\Psi\Omega_M\)); on the
  fields of the fold the gauge and spin covariant derivative \(D_\mu\)
  (the partial derivative plus the spin connection of the fold and the
  gauge fields times their couplings and generators);
  \(D^\mu = \hat h^{\mu\nu}D_\nu\); on an adjoint field
  \((D_\mu X)^A = \partial_\mu X^A + g_3f^{ABC}G^B_\mu X^C\) (likewise
  for \(SU(2)\) with \(g_2\epsilon^{IJK}\)). Unit: 1/length. Status:
  STANDARD.
\item
  \(Q,\ \bar Q\): The left-handed quark doublets of the three
  generations, \(Q = (u_{\mathrm{L}}, d_{\mathrm{L}})\) (colour
  triplets, hypercharge \(1/6\)). Unit: GeV\(^{3/2}\). Status: AUTHOR.
\item
  \(u_{\mathrm{R}}\): The right-handed up-type quarks (up, charm and
  top; singlets of \(SU(2)_L\), hypercharge \(2/3\)). Unit:
  GeV\(^{3/2}\). Status: STANDARD.
\item
  \(d_{\mathrm{R}}\): The right-handed down-type quarks (down, strange
  and bottom; hypercharge \(-1/3\)). Unit: GeV\(^{3/2}\). Status:
  STANDARD.
\item
  \(\ell,\ \bar\ell\): The left-handed lepton doublets of the three
  generations, \(\ell = (\nu_{\mathrm{L}}, e_{\mathrm{L}})\)
  (hypercharge \(-1/2\)). Unit: GeV\(^{3/2}\). Status: AUTHOR.
\item
  \(e_{\mathrm{R}}\): The right-handed electron, muon and tau (singlets,
  hypercharge \(-1\)). Unit: GeV\(^{3/2}\). Status: AUTHOR.
\item
  \(\nu_{\mathrm{R}}\): The three right-handed neutrinos (gauge
  singlets). Unit: GeV\(^{3/2}\). Status: ASSUMED.
\item
  \(\nu\): A four-dimensional coordinate index (as \(\mu\)). Status:
  STANDARD.
\item
  \(\hat h_{\mu\nu}\): The four-dimensional metric of the fold,
  \(\hat h_{\mu\nu} = \mathrm{diag}(\mathrm{e}^{2a_4}, \mathrm{e}^{2a_4}, \mathrm{e}^{2a_4}, -1)\)
  in the coordinates \(x_1, \dots, x_4\) (the induced metric at
  \(\sigma = 0\) without the extra times); \(\hat h^{\mu\nu}\) its
  inverse, \(\hat h\) its determinant,
  \(\sqrt{-\hat h} = \mathrm{e}^{3a_4}\); \(\hat h^{(1)}_{\mu\nu}\) its
  first-order cosmological perturbation. Unit: dimensionless. Status:
  DERIVED.
\item
  \(\mathbf{1},\ \mathbf{1}_8\): The identity matrix of the size fixed
  by the context (\(16 \times 16\) on dirac16complex, \(4 \times 4\) on
  a four-dimensional spinor, \(3 \times 3\) in generation or flavour
  space); \(\mathbf{1}_8\) is the \(8 \times 8\) identity matrix. Unit:
  dimensionless. Status: STANDARD.
\item
  \({}^\dagger\): Hermitian conjugation: the complex conjugate of the
  transpose of a matrix or spinor; for an operator on the space of
  states, its adjoint with respect to the scalar product of that space
  (for the quantised dirac16complex the indefinite Krein scalar
  product). Status: STANDARD.
\item
  \(\beta\): The matrix of the four-dimensional Dirac adjoint,
  \(\beta = \hat\gamma^{(4)}\) (Hermitian, \(\beta^2 = \mathbf{1}\)).
  Unit: dimensionless. Status: STANDARD.
\item
  \(\hat\gamma_5\): The four-dimensional chirality matrix,
  \(\hat\gamma_5 = \mathrm{i}\hat\gamma^{(4)}\hat\gamma^{(1)}\hat\gamma^{(2)}\hat\gamma^{(3)}\),
  \(\hat\gamma_5^2 = \mathbf{1}\). Unit: dimensionless. Status:
  STANDARD.
\item
  \(\hat P_{\mathrm{L}}\): The four-dimensional left-handed projector,
  \(\hat P_{\mathrm{L}} = (\mathbf{1} - \hat\gamma_5)/2\). Unit:
  dimensionless. Status: STANDARD.
\item
  \(\hat P_{\mathrm{R}}\): The four-dimensional right-handed projector,
  \(\hat P_{\mathrm{R}} = (\mathbf{1} + \hat\gamma_5)/2\). Unit:
  dimensionless. Status: STANDARD.
\end{itemize}

\[
\mathcal{L}_{\mathrm{Y}} = -\big(\bar Q\,\mathcal{Y}^{u}\,\tilde\Phi\,u_{\mathrm{R}} + \bar Q\,\mathcal{Y}^{d}\,\Phi\,d_{\mathrm{R}}
 + \bar\ell\,\mathcal{Y}^{e}\,\Phi\,e_{\mathrm{R}} + \bar\ell\,\mathcal{Y}^{\nu}\,\tilde\Phi\,\nu_{\mathrm{R}}\big)
 + \mathrm{h.c.},\qquad \tilde\Phi = \mathrm{i}\tau^2\Phi^* \tag{IN-15}
\] \textbf{Symbols in (IN-15).}

\begin{itemize}
\tightlist
\item
  \(\mathcal{L}_{\mathrm{Y}}\): The Yukawa Lagrangian of the quarks and
  leptons with the Higgs doublet. Unit: GeV\(^4\). Status: STANDARD.
\item
  \(Q,\ \bar Q\): The left-handed quark doublets of the three
  generations, \(Q = (u_{\mathrm{L}}, d_{\mathrm{L}})\) (colour
  triplets, hypercharge \(1/6\)). Unit: GeV\(^{3/2}\). Status: AUTHOR.
\item
  \(\mathcal{Y}^{u},\ \mathcal{Y}^{d},\ \mathcal{Y}^{e},\ \mathcal{Y}^{\nu}\):
  The Yukawa matrices (\(3 \times 3\) complex matrices in generation
  space) of the up quarks, the down quarks, the charged leptons and the
  neutrinos. Unit: dimensionless. Status: STANDARD.
\item
  \(\tilde\Phi\): The conjugate Higgs doublet,
  \(\tilde\Phi = \mathrm{i}\tau^2\Phi^*\) (hypercharge \(-1/2\)). Unit:
  GeV. Status: STANDARD.
\item
  \(u_{\mathrm{R}}\): The right-handed up-type quarks (up, charm and
  top; singlets of \(SU(2)_L\), hypercharge \(2/3\)). Unit:
  GeV\(^{3/2}\). Status: STANDARD.
\item
  \(\Phi\): The Higgs doublet (two complex components, hypercharge
  \(1/2\)), a field of the fold. Unit: GeV. Status: AUTHOR.
\item
  \(d_{\mathrm{R}}\): The right-handed down-type quarks (down, strange
  and bottom; hypercharge \(-1/3\)). Unit: GeV\(^{3/2}\). Status:
  STANDARD.
\item
  \(\ell,\ \bar\ell\): The left-handed lepton doublets of the three
  generations, \(\ell = (\nu_{\mathrm{L}}, e_{\mathrm{L}})\)
  (hypercharge \(-1/2\)). Unit: GeV\(^{3/2}\). Status: AUTHOR.
\item
  \(e_{\mathrm{R}}\): The right-handed electron, muon and tau (singlets,
  hypercharge \(-1\)). Unit: GeV\(^{3/2}\). Status: AUTHOR.
\item
  \(\nu_{\mathrm{R}}\): The three right-handed neutrinos (gauge
  singlets). Unit: GeV\(^{3/2}\). Status: ASSUMED.
\item
  \(\mathrm{h.c.}\): The Hermitian conjugate of the preceding terms.
  Status: STANDARD.
\item
  \(\mathrm{i}\): The imaginary unit, \(\mathrm{i}^2 = -1\). Unit:
  dimensionless. Status: STANDARD.
\item
  \(\tau^I\): The Pauli matrices acting on weak doublets
  (\(\tau^1, \tau^2, \tau^3\); \(\tau^2\) is the second one, not a
  square). Unit: dimensionless. Status: STANDARD.
\item
  \({}^*\): Complex conjugation of a number, or of each entry of a
  matrix or column. Status: STANDARD.
\end{itemize}

The gauge fields of the bulk act on the fields of the fold through their
values at the fold (IN-13: for the massless modes, the canonically
normalised four-dimensional gauge fields with the measured couplings),
the three generations of quarks and leptons with three right-handed
neutrinos (IN-14), and their Yukawa couplings to the Higgs doublet
(IN-15). With the right-handed neutrinos the gauge anomalies still
cancel (DERIVED, \texttt{um4\_fields.py}), and the neutrino Yukawa
couplings give the neutrinos Dirac masses, whose values are chosen to
reproduce the measured mass differences with the lightest mass zero
(ASSUMED). Status: STANDARD (UM-4 adds the right-handed neutrinos and
puts the gauge fields in the bulk).

\hypertarget{the-portals-of-dirac16complex}{%
\subsection{The portals of
dirac16complex}\label{the-portals-of-dirac16complex}}

\[
\begin{aligned}
\mathcal{L}_{\mathrm{P}} = {}&-c_\Phi\,\big(\Phi^\dagger\Phi\big)\,\bar\Psi P_{\mathrm{L}}\Psi
 - \sum_{X}\Big[\frac14\big(c^{G}_X\,G^A_{\mu\nu}G^{A\mu\nu} + c^{W}_X\,W^I_{\mu\nu}W^{I\mu\nu}
 + c^{B}_X\,B_{\mu\nu}B^{\mu\nu}\big)\,\bar\Psi P_X\Psi\\
&\qquad + \frac{d_X}{2}\,B_{\mu\nu}\,\mathcal{D}_X^{\mu\nu} + \sum_{f} c^{f}_X\,J_f^\mu\,j_{X\mu}\Big],\\
\mathcal{D}_X^{\mu\nu} = {}&\mathrm{i}\,\bar\Psi\,\Sigma^{\mu\nu}P_X\Psi,\qquad
j_X^\mu = \mathrm{i}\big(\bar\Psi P_X D^\mu\Psi - (D^\mu\bar\Psi)P_X\Psi\big),\qquad
J_f^\mu = \bar\psi_f\,\hat\gamma^\mu\psi_f
\end{aligned} \tag{IN-16}
\] \textbf{Symbols in (IN-16).}

\begin{itemize}
\tightlist
\item
  \(\mathcal{L}_{\mathrm{P}}\): The Lagrangian of the portals of
  dirac16complex at the fold. Unit: GeV\(^4\). Status: INVENTED.
\item
  \(c_\Phi\): The Higgs-portal coupling of the left-handed half of
  dirac16complex (the right-handed half has none). Unit: GeV\(^{-5}\).
  Status: INVENTED.
\item
  \(\Phi\): The Higgs doublet (two complex components, hypercharge
  \(1/2\)), a field of the fold. Unit: GeV. Status: AUTHOR.
\item
  \({}^\dagger\): Hermitian conjugation: the complex conjugate of the
  transpose of a matrix or spinor; for an operator on the space of
  states, its adjoint with respect to the scalar product of that space
  (for the quantised dirac16complex the indefinite Krein scalar
  product). Status: STANDARD.
\item
  \(\Psi,\ \bar\Psi\): dirac16complex: the author's 16-component complex
  spinor field of eight dimensions;
  \(\bar\Psi = \Psi^\dagger\sigma_{16}\) is its Dirac adjoint; in a term
  of the fold, \(\Psi\) is its value at \(\sigma = 0\). Unit:
  GeV\(^{7/2}\). Status: AUTHOR.
\item
  \(P_{\mathrm{L}}\): The author's left-handed chiral projector,
  \(P_{\mathrm{L}} = (\mathbf{1} - \Gamma_8)/2\). Unit: dimensionless.
  Status: AUTHOR.
\item
  \(X\): The label of a half of dirac16complex:
  \(X \in \{\mathrm{L}, \mathrm{R}\}\) (left-handed or right-handed
  half); a sum over \(X\) runs over both. Status: STANDARD.
\item
  \(c^{G}_X\): The gluon-kinetic portal coupling of the half \(X\) of
  dirac16complex. Unit: GeV\(^{-7}\). Status: INVENTED.
\item
  \(G^A_M,\ G^A_\mu,\ G^A_{MN}\): The gluon field of \(SU(3)_C\)
  (\(A = 1, \dots, 8\)): in the bulk \(G^A_M\), a field of eight
  dimensions in its geometric normalisation (the covariant derivative
  \(\partial - \mathrm{i}G\)), and on the fold the canonically
  normalised four-dimensional field \(G^A_\mu(x)\) of its massless mode,
  \(g_3G^A_\mu(x) = G^A_\mu(\sigma = 0, x)\); \(G^A_{MN}\) and
  \(G^A_{\mu\nu}\) their field strengths. Unit: GeV (field), GeV\(^2\)
  (field strength). Status: AUTHOR.
\item
  \(A\): A colour (adjoint SU(3)) index running over \(1, \dots, 8\); a
  repeated index is summed. Status: STANDARD.
\item
  \(\mu\): A four-dimensional coordinate index running over the
  observable coordinates \(x_1, x_2, x_3, x_4\) of the fold (signature
  \(+,+,+,-\)); a repeated index is summed. Status: STANDARD.
\item
  \(\nu\): A four-dimensional coordinate index (as \(\mu\)). Status:
  STANDARD.
\item
  \(c^{W}_X\): The \(W\)-kinetic portal coupling of the half \(X\).
  Unit: GeV\(^{-7}\). Status: INVENTED.
\item
  \(W^I_M,\ W^I_\mu,\ W^I_{MN}\): The gauge field of \(SU(2)_L\)
  (\(I = 1, 2, 3\)): in the bulk \(W^I_M\) (geometric normalisation), on
  the fold the canonically normalised four-dimensional field
  \(W^I_\mu(x)\) of its massless mode,
  \(g_2W^I_\mu(x) = W^I_\mu(\sigma = 0, x)\); \(W^I_{MN}\),
  \(W^I_{\mu\nu}\) their field strengths;
  \(W^\pm_\mu = (W^1_\mu \mp \mathrm{i}W^2_\mu)/\sqrt2\). Unit: GeV
  (field), GeV\(^2\) (field strength). Status: AUTHOR.
\item
  \(I\): A weak-isospin (adjoint SU(2)) index running over \(1, 2, 3\);
  a repeated index is summed. Status: STANDARD.
\item
  \(c^{B}_X\): The \(B\)-kinetic portal coupling of the half \(X\).
  Unit: GeV\(^{-7}\). Status: INVENTED.
\item
  \(B_M,\ B_\mu,\ B_{MN}\): The hypercharge gauge field of \(U(1)_Y\):
  in the bulk \(B_M\) (geometric normalisation), on the fold the
  canonically normalised four-dimensional field \(B_\mu(x)\) of its
  massless mode, \(g_1B_\mu(x) = B_\mu(\sigma = 0, x)\); \(B_{MN}\),
  \(B_{\mu\nu}\) their field strengths. Unit: GeV (field), GeV\(^2\)
  (field strength). Status: AUTHOR.
\item
  \(P_X\): The chiral projector of the half \(X\):
  \(P_X = P_{\mathrm{L}}\) for \(X = \mathrm{L}\), \(P_{\mathrm{R}}\)
  for \(X = \mathrm{R}\). Unit: dimensionless. Status: AUTHOR.
\item
  \(d_X\): The dipole-portal coupling of the half \(X\) of
  dirac16complex with the hypercharge field strength. Unit:
  GeV\(^{-5}\). Status: INVENTED.
\item
  \(\mathcal{D}_X^{\mu\nu}\): The dipole (antisymmetric tensor) density
  of the half \(X\) of dirac16complex at the fold,
  \(\mathcal{D}_X^{\mu\nu} = \mathrm{i}\bar\Psi\Sigma^{\mu\nu}P_X\Psi\)
  (real, DERIVED). Unit: GeV\(^7\). Status: DERIVED.
\item
  \(f\): The label of a kind of Standard-Model multiplet:
  \(f \in \{Q, u, d, \ell, e, \nu\}\) (quark doublets, up-quark
  singlets, down-quark singlets, lepton doublets, charged-lepton
  singlets, right-handed neutrinos); a sum over \(f\) runs over the six.
  Status: STANDARD.
\item
  \(c^{f}_X\): The current-portal coupling of the half \(X\) of
  dirac16complex with the fermion-number current of the multiplets of
  kind \(f\) (flavour blind); \(c^{\ell}_X\) and \(c^{\nu}_X\) are those
  of the lepton doublets and of the right-handed neutrinos. Unit:
  GeV\(^{-7}\). Status: INVENTED.
\item
  \(J_f^\mu\): The fermion-number current of the Standard-Model
  multiplets of kind \(f\), \(J_f^\mu = \bar\psi_f\hat\gamma^\mu\psi_f\)
  summed over the three generations and the gauge components; with other
  labels the currents that are the sources of the gauge fields. Unit:
  GeV\(^3\). Status: STANDARD.
\item
  \(j_X^\mu\): The convective current of the half \(X\) of
  dirac16complex,
  \(j_X^\mu = \mathrm{i}(\bar\Psi P_XD^\mu\Psi - (D^\mu\bar\Psi)P_X\Psi)\)
  (real and joining the half with itself, DERIVED). Unit: GeV\(^8\).
  Status: DERIVED.
\item
  \(\mathrm{i}\): The imaginary unit, \(\mathrm{i}^2 = -1\). Unit:
  dimensionless. Status: STANDARD.
\item
  \(\Sigma^{ab},\ \Sigma^{\mu\nu}\): The generators of Spin(4,4) on
  dirac16complex, \(\Sigma^{ab} = \frac14[T^a, T^b]\) (the scaled
  commutators of the author's matrices);
  \(\Sigma^{\mu\nu} = \mathbf{e}_a{}^\mu\mathbf{e}_b{}^\nu\Sigma^{ab}\)
  along the fold. Unit: dimensionless. Status: AUTHOR.
\item
  \(D_M,\ D_\mu\): The covariant derivative: on dirac16complex
  \(D_M = \partial_M + \Omega_M\) (with
  \(D_M\bar\Psi = \partial_M\bar\Psi - \bar\Psi\Omega_M\)); on the
  fields of the fold the gauge and spin covariant derivative \(D_\mu\)
  (the partial derivative plus the spin connection of the fold and the
  gauge fields times their couplings and generators);
  \(D^\mu = \hat h^{\mu\nu}D_\nu\); on an adjoint field
  \((D_\mu X)^A = \partial_\mu X^A + g_3f^{ABC}G^B_\mu X^C\) (likewise
  for \(SU(2)\) with \(g_2\epsilon^{IJK}\)). Unit: 1/length. Status:
  STANDARD.
\item
  \(\psi,\ \bar\psi\): A Standard-Model fermion field (a
  four-dimensional Weyl field of the fold with three generations,
  written as a four-component spinor of fixed chirality);
  \(\bar\psi = \psi^\dagger\beta\); \(\psi_f\) the fields of the kind
  \(f\). Unit: GeV\(^{3/2}\). Status: STANDARD.
\item
  \(\hat\gamma^\mu\): The four-dimensional complex \(4 \times 4\) Dirac
  matrices of the fold: \(\hat\gamma^{(a)}\) along the frame directions
  of the fold and \(\hat\gamma^\mu\) along its coordinates (the frame
  matrices contracted with the inverse frame field of the fold),
  \(\{\hat\gamma^\mu, \hat\gamma^\nu\} = -2\hat h^{\mu\nu}\mathbf{1}\);
  \(\hat\gamma_\mu = \hat h_{\mu\nu}\hat\gamma^\nu\). Unit:
  dimensionless. Status: STANDARD.
\end{itemize}

The new interactions (the same as in UM-3): every Standard-Model field
couples at the fold to the left-handed and to the right-handed half of
dirac16complex, each with its own coupling, through gauge-invariant
operators: the Higgs portal (left-handed half only), the gauge-kinetic
portals of the gluons, \(W\) and \(B\), a dipole portal to \(B\), and
current portals between the fermion-number current of every multiplet of
the Standard Model and the convective currents of the two halves. Each
is invariant under the gauge group of the Standard Model and contains
dirac16complex only in bilinears, so the charge of dirac16complex is
conserved and its lightest states cannot decay into the Standard Model
alone; in UM-4 they decay slowly into the waves of dirac16complex in the
bulk (the resonances). The charge is indefinite (a Krein form);
quantised with positive energies, every quantum carries a positive
energy, so the vacuum is stable (whether a gas of the particles is
classically stable is not decided); but the modes of positive Krein
charge are ghosts (negative norm), and every coupling of the fold that
distinguishes the two chiralities, the Higgs portal of the left-handed
half among them, creates them together with physical particles: with
these portals UM-4 is not a unitary quantum theory, and a unitary UM-4
would need couplings of the fold that do not distinguish the chiralities
(necessary, not shown to be sufficient). Status: INVENTED (the portals),
DERIVED (the stability of the vacuum and the ghosts,
\texttt{um4\_krein.py}), OPEN (the stability of the gas, a unitary
coupling).

\hypertarget{the-parameters}{%
\section{The parameters}\label{the-parameters}}

\begin{longtable}[]{@{}
  >{\raggedright\arraybackslash}p{(\columnwidth - 6\tabcolsep) * \real{0.2500}}
  >{\raggedright\arraybackslash}p{(\columnwidth - 6\tabcolsep) * \real{0.2500}}
  >{\raggedright\arraybackslash}p{(\columnwidth - 6\tabcolsep) * \real{0.2500}}
  >{\raggedright\arraybackslash}p{(\columnwidth - 6\tabcolsep) * \real{0.2500}}@{}}
\toprule\noalign{}
\begin{minipage}[b]{\linewidth}\raggedright
parameter
\end{minipage} & \begin{minipage}[b]{\linewidth}\raggedright
meaning
\end{minipage} & \begin{minipage}[b]{\linewidth}\raggedright
value
\end{minipage} & \begin{minipage}[b]{\linewidth}\raggedright
status
\end{minipage} \\
\midrule\noalign{}
\endhead
\bottomrule\noalign{}
\endlastfoot
\(H\) & the author's constant: the period \(\pi/(6H)\) of \(x_0\), the
warp factor, the curvature \(-H^2\) of the static block & free, bounded
from below: \(H \gtrsim 2.5\) TeV by the invisible width of the Z boson
(its decay into the bulk), 1.3 TeV by contact interactions, 0.97 TeV by
the parameter rho of the electroweak fits, \(7.7\times10^{-3}\) eV by
Newton's law at 52 micrometres & AUTHOR (the constant), DERIVED and
ESTIMATE (the bounds) \\
\(a_4(x_4)\), \(b_4(x_4)\) & the author's functions of the time & frozen
extra times, \(\mathrm{e}^{-a_4}b_4\) constant, with the expansion
\(a_4(x_4)\) of general relativity (the metric FE-27 solves the field
equations exactly only for the static block; around an expanding fold
the bulk metric is corrected by terms of the order of the curvature of
the fold divided by \(H^2\), which vanish at the fold); the author's
vacuum \(a_4 = Hx_4\) with constant \(b_4\) is not a solution & AUTHOR,
DERIVED \\
\(w_{(p)}\), \(\Lambda\) & the Lovelock weights and the cosmological
constant of the bulk & \(w = (1, 0, 0)\) (Einstein gravity),
\(\Lambda = -21H^2\): gravitons on the light cone on every background
(the author's original values \((1, 1/32, 1/1152)\) and \(-10H^2\);
UM-3: \((1, 1/20, 1/360)\) and \(-7H^2\)) & AUTHOR (the form), DERIVED
(the values that the order selects) \\
\(\kappa = 8\pi G_8\) & the gravitational coupling of the bulk & free;
Newton's constant fixes \(5\kappa H/(2V_*)\) & ASSUMED \\
\(\kappa_7\) & the gravitational coupling of the fold & \(5\kappa H/2\);
\(8\pi G_{\mathrm{N}} = \kappa_7/V_*\) (plus \(1/\kappa_{\mathrm{ind}}\)
in \(1/(8\pi G_{\mathrm{N}})\)) & DERIVED \\
\(\kappa_{\mathrm{ind}}\) & the coupling of the induced Einstein-Hilbert
term of the fold & negligible for a cutoff of the order of \(H\); not
fixed by UM-4 & ESTIMATE \\
\(V_*\) & the physical volume of the extra times at the fold & free;
Newton's constant fixes \(\kappa H/V_*\) & ASSUMED \\
\(\mathsf{g}_1\), \(\mathsf{g}_2\), \(\mathsf{g}_3\) & the gauge
couplings of eight dimensions & \(\mathsf{g}_k^2 = 2V_*g_k^2/(3H)\) from
the measured couplings & ASSUMED (the place), DERIVED (the relation) \\
\(\xi_{\mathrm{H}}\) & the coupling of the Higgs doublet to the
curvature of the fold & 0: the Higgs Lagrangian of UM-4 (IN-11) has no
such term; about \(1.6\times10^4\) if the Higgs field drove an inflation
of the fold in an extension of the action (ESTIMATE), which would need
\(H > 1.6\times10^{14}\) GeV & not in the action; ESTIMATE (the
extension) \\
\(\rho_{\mathrm{ff}}\), \(p_{\mathrm{ff}}\), \(q_{\mathrm{ff}}\) & the
fold fluid (energy density, pressure along ordinary space, pressure
along the extra times) & a tension:
\(p_{\mathrm{ff}} = -\rho_{\mathrm{ff}}\); with the vacuum energy
\(\rho_{\mathrm{vac}}\) of the fields,
\(\rho_{\mathrm{ff}} + \rho_{\mathrm{vac}} = \lambda_{\mathrm{F}}V_* + \rho_\Lambda\),
\(\lambda_{\mathrm{F}}V_* = 30H^2M_{\mathrm{red}}^2\),
\(\rho_\Lambda = (2.2\ \mathrm{meV})^4\) the dark energy;
\(q_{\mathrm{ff}} = -\rho_{\mathrm{ff}} + q_{\mathrm{hold}}\),
\(q_{\mathrm{hold}} = -0.842\) times the critical density today (with no
pressure of the baryons along the extra times) & ASSUMED (its values
DERIVED) \\
\(m_\Psi\) & the bulk mass of dirac16complex & free; the particles
\(\chi\) of dark matter live longer than the universe only if
\(m_\Psi/H\) is large (above about \(1.1\times10^3\) for
\(q_{\mathrm{b}} = m_\Psi/2\)) & RECORD \\
\(\lambda_\Psi\) & the self-coupling of dirac16complex & free; at
\(m_\Psi/H = 10^4\) a gas of the particles \(\chi\) lives longer than
the universe only if its dimensionless strength is below 0.147 &
RECORD \\
\(c_\Phi\), \(c^{G}_X\), \(c^{W}_X\), \(c^{B}_X\), \(d_X\), \(c^{f}_X\)
& the portal couplings (\(X\) = L, R) & free; direct detection and the
effective number of neutrino species \(N_{\mathrm{eff}}\) (the energy
density of the relativistic species besides the photons, in units of one
neutrino species) constrain them; a binding of the particles \(\chi\) to
the fold by the Higgs portal alone is excluded by direct detection
together with their lifetime & INVENTED \\
\(v\), \(\lambda\), \(\mu^2\) & the Higgs vacuum and potential &
\(v = 246.22\) GeV, \(\lambda = 0.129\), \(\mu^2 = -\lambda v^2\) &
AUTHOR, STANDARD \\
\(v_{\mathcal{A}}\), \(a_{40}\) & amplitude and rate of the ad hoc
vacuum \(\mathcal{A} = v_{\mathcal{A}}\mathrm{e}^{-2a_{40}x_4}\) & the
author's amplitude is 1 in the units of his notebook; UM-4 keeps
\(v_{\mathcal{A}}\) free (the examples take \(v\) and \(10v\));
\(a_{40} > 0\); nucleosynthesis needs
\(a_{40} > \ln(v_{\mathcal{A}}/(0.1v))/(2\ \text{s})\) & AUTHOR, ASSUMED
(a free amplitude) \\
\(g_1\), \(g_2\), \(g_3\) & the four-dimensional gauge couplings & the
measured values & STANDARD \\
\(\mathcal{Y}^{u}\), \(\mathcal{Y}^{d}\), \(\mathcal{Y}^{e}\),
\(\mathcal{Y}^{\nu}\) & the Yukawa matrices & the measured masses and
mixings (Particle Data Group); neutrinos (the global fit NuFIT 5.2 of
the oscillation data, November 2022, normal ordering, the fit without
the Super-Kamiokande atmospheric data): the differences of the squared
masses of the three mass states \(\Delta m^2_{21} = 7.41\times10^{-5}\)
eV\(^2\) and \(\Delta m^2_{31} = 2.511\times10^{-3}\) eV\(^2\), the
mixing angles \(\theta_{12} = 33.40^\circ\),
\(\theta_{13} = 8.54^\circ\), \(\theta_{23} = 49.14^\circ\); the
lightest mass 0 and the Dirac phase 0. The newer fit NuFIT 6.0 (2024)
changes the results by at most 1.5 per cent; the Dirac phase (212
degrees in NuFIT 6.0) drops out of every result of UM-4; DESI DR2 with
the CMB and NuFIT 6.0 bounds the lightest mass below 23 meV & STANDARD,
ASSUMED (the lightest mass; the phase is irrelevant, DERIVED) \\
\(A_{\mathrm{s}}\), \(n_{\mathrm{S}}\), \(\tau_{\mathrm{reion}}\) & the
primordial spectrum and the optical depth & Planck 2018:
\(\ln(10^{10}A_{\mathrm{s}}) = 3.044\), \(n_{\mathrm{S}} = 0.9649\),
\(\tau_{\mathrm{reion}} = 0.0544\); they need an inflation of the fold
with frozen extra times, in the regime of general relativity only if its
rate is below \(H/10\) (its inflaton is not specified; the Higgs field
only with a curvature coupling that the action does not contain) &
ASSUMED (inputs; the inflaton and the holding pressure OPEN) \\
\end{longtable}

\hypertarget{further-symbols}{%
\section{Further symbols}\label{further-symbols}}

\textbf{Further symbols.} These symbols appear in the definitions of the
symbols listed after the equations of this document.

\begin{itemize}
\tightlist
\item
  \(M_{\mathrm{red}}\): The reduced Planck mass,
  \(M_{\mathrm{red}} = (8\pi G_{\mathrm{N}})^{-1/2}\). Unit: GeV. Value:
  \(2.435\times10^{18}\) GeV. Status: STANDARD.
\item
  \(k\): A wave number: the momentum of a plane wave along \(x_3\),
  \(\mathrm{e}^{\mathrm{i}(kx_3 - \omega x_4)}\); in the cosmological
  perturbations the comoving wave number of a Fourier mode,
  \(k = \lvert\vec k\rvert\). Unit: GeV (1/length). Status: STANDARD.
\item
  \(G_{\mathrm{F}}\): The Fermi constant,
  \(G_{\mathrm{F}} = 1/(\sqrt2\,v^2)\) in the Standard-Model vacuum (its
  measured value); in the equations of the neutrino flavours it stands
  for the local Fermi constant of the ad hoc vacuum. Unit: GeV\(^{-2}\).
  Value: \(1.1664\times10^{-5}\) GeV\(^{-2}\). Status: STANDARD.
\item
  \(G_{\mathrm{ind}}\): The Newton constant of the induced
  Einstein-Hilbert term of the fold,
  \(\kappa_{\mathrm{ind}} = 8\pi G_{\mathrm{ind}}\). Unit: GeV\(^{-2}\).
  Status: ESTIMATE.
\item
  \(\Delta\): A change, difference or small variation of the quantity
  that follows (for example an interval of the time, a perturbation of a
  density, a difference of two squared masses, or an addition to an
  effective number of species). Status: STANDARD.
\item
  \(\kappa_7\): The gravitational coupling of the fold: to the first
  order in its curvature over \(H^2\) the fold obeys the Einstein
  equations of seven dimensions
  \(\hat G^a{}_b = \kappa_7\mathsf{t}^a{}_b\) with
  \(\kappa_7 = 5\kappa H/2\), the coupling of the massless graviton of
  the bulk bound to the fold; four-dimensional gravity on the fold has
  \(8\pi G_{\mathrm{N}} = \kappa_7/V_*\) (plus the induced term, which
  is negligible). Unit: GeV\(^{-5}\). Value: \(5\kappa H/2\). Status:
  DERIVED.
\item
  \(\mathrm{tr}\): The trace of a matrix: the sum of its diagonal
  entries. Status: STANDARD.
\item
  \(\nu_{\mathrm{L}}\): The upper (neutral, left-handed) components of
  the lepton doublets: the left-handed neutrinos. Unit: GeV\(^{3/2}\).
  Status: STANDARD.
\item
  \(d_{\mathrm{L}}\): The lower (down-type, left-handed) components of
  the quark doublets. Unit: GeV\(^{3/2}\). Status: STANDARD.
\item
  \(e_{\mathrm{L}}\): The lower (charged, left-handed) components of the
  lepton doublets. Unit: GeV\(^{3/2}\). Status: STANDARD.
\item
  \(m\): A mass: \(m_i\) the mass of the neutrino mass state \(i\);
  \(\Delta m_{ij}^2 = m_i^2 - m_j^2\) the squared-mass difference of the
  states \(i\) and \(j\) (\(\Delta m_{21}^2\), \(\Delta m_{31}^2\),
  \(\Delta m_{32}^2\)); \(\Delta m^2\) without labels the difference of
  the two states of a resonance, \(m_2^2 - m_1^2\), multiplied by
  \(v_{\mathrm{loc}}^2/v^2\) where this differs from 1; \(m\) without a
  label the mass of a mode of a field of the bulk as the fold sees it,
  \(m^2 = \omega^2 - k^2\); for the W and Z bosons their complex pole
  mass, \(m^2 = m_0^2 - \mathrm{i}m_0\Gamma_{\mathrm{bulk}}\) to the
  leading order; in the relic abundance of the particles \(\chi\) their
  mass. Unit: GeV. Status: STANDARD.
\item
  \(m_h\): The mass of the Higgs boson. Unit: GeV. Value: 125.25 GeV
  (Particle Data Group 2022). Status: STANDARD.
\item
  \(u_{\mathrm{L}}\): The upper (up-type, left-handed) components of the
  quark doublets. Unit: GeV\(^{3/2}\). Status: STANDARD.
\item
  \(\Gamma_{\mathrm{bulk}}\): The width of the W or the Z boson into the
  bulk (its decay into the massive modes of the gauge fields of the
  bulk), \(\Gamma_{\mathrm{bulk}} = m_0(m_0/H)^3/3\). Unit: GeV. Value:
  1.4 MeV for the Z and 0.9 MeV for the W at \(H = 2.5\) TeV. Status:
  DERIVED.
\item
  \(\mathsf{t}^a{}_b,\ \mathsf{t}\): The matter of the fold above its
  tension: the thin source minus the tension,
  \(\mathcal{S}^a{}_b = -\lambda_{\mathrm{F}}\delta^a{}_b + \mathsf{t}^a{}_b\)
  (eight-dimensional normalisation); \(\mathsf{t} = \mathsf{t}^c{}_c\)
  its trace; with frozen extra times
  \(\mathsf{t}^a{}_b = \mathrm{diag}(P_7, P_7, P_7, -\rho_7, q_7, q_7, q_7)\).
  Unit: GeV\(^7\). Status: DERIVED.
\item
  \(i\): A generation index (\(i = 1, 2, 3\)), which also labels the
  three neutrino mass states; in a frame component such as
  \(\mathcal{S}^{(i)}{}_{(i)}\) one of the three directions of ordinary
  space (the three components are equal), and in a tensor of ordinary
  space such as \(\mathsf{h}_{ij}\) one of the directions
  \(x_1, x_2, x_3\); a repeated index is not summed unless a sum sign is
  written. Status: STANDARD.
\item
  \(m_0\): The mass of the W or the Z boson without the massive modes of
  the bulk (the mass of the massless mode with the Higgs vacuum),
  \(m_0^2 = \frac32HM_{\mathrm{loc}}^2\);
  \(m_0^2 = g_2^2(v^2 + \mathcal{A}^2)/4\) for the W. Unit: GeV. Status:
  DERIVED.
\item
  \(m_u, \dots, m_3\): The masses of the quarks (\(m_u, m_c, m_t\);
  \(m_d, m_s, m_b\)), of the charged leptons (\(m_e, m_\mu, m_\tau\))
  and of the three neutrino mass states (\(m_1, m_2, m_3\), Dirac
  masses). Unit: GeV. Value: measured values;
  \(\Delta m_{21}^2 = 7.41\times10^{-5}\) eV\(^2\),
  \(\lvert\Delta m_{31}^2\rvert = 2.511\times10^{-3}\) eV\(^2\) in the
  computations of UM-4 (the global fit NuFIT 5.2 of the oscillation
  data, November 2022), \(7.49\times10^{-5}\) and \(2.513\times10^{-3}\)
  eV\(^2\) in NuFIT 6.0 (2024); the lightest mass \(m_1 = 0\) in the
  computations, below 23 meV by DESI DR2 with the CMB and NuFIT 6.0.
  Status: STANDARD.
\item
  \(v_{\mathrm{loc}}\): The local value of the electroweak vacuum,
  \(v_{\mathrm{loc}}^2 = 2\langle\Phi^\dagger\Phi\rangle\): it sets the
  local masses of the quarks and leptons (each mass times
  \(v_{\mathrm{loc}}/v\)) and the local Fermi constant;
  \(v_{\mathrm{loc}} = v\) in the Standard-Model vacuum. Unit: GeV.
  Value: \(v\) in the late universe outside dense dirac16complex.
  Status: DERIVED.
\item
  \(M_{\mathrm{loc}}\): The mass term that the Higgs vacuum on the fold
  gives a gauge field of the bulk:
  \(M_{\mathrm{loc}}^2\delta_{\mathrm{D}}(\sigma)\) in its profile
  equation; for the W,
  \(M_{\mathrm{loc}}^2 = \mathsf{g}_2^2(v^2 + \mathcal{A}^2)/(4V_*)\).
  Unit: GeV\(^{1/2}\) (\(M_{\mathrm{loc}}^2\) in GeV). Status: DERIVED.
\item
  \(\mathsf{h}_{ij}\): The transverse traceless perturbation of the
  metric of ordinary space by a gravitational wave as the
  four-dimensional world sees it:
  \(\hat h_{ij} = \mathrm{e}^{2a_4}(\delta_{ij} + \mathsf{h}_{ij})\)
  (\(i, j = 1, 2, 3\)). Unit: dimensionless. Status: STANDARD.
\item
  \(\rho_7\): The energy density of the matter of the fold above its
  tension in the seven-dimensional normalisation (minus the component
  \(\mathsf{t}^{(4)}{}_{(4)}\): an energy per unit volume of the six
  other directions of the fold); \(V_*\rho_7\) in the four-dimensional
  normalisation. Unit: GeV\(^7\). Status: DERIVED.
\item
  \(q_7\): The pressure along the extra times of the matter of the fold
  above its tension in the seven-dimensional normalisation (the
  component \(\mathsf{t}^{(r)}{}_{(r)}\)). Unit: GeV\(^7\). Status:
  DERIVED.
\item
  \(A_{\mathrm{s}}\): The amplitude of the primordial curvature spectrum
  at the pivot \(k_*\). Unit: dimensionless. Value:
  \(2.10\times10^{-9}\) (Planck 2018). Status: ASSUMED.
\item
  \(N_{\mathrm{eff}}\): The effective number of neutrino species: the
  energy density of the relativistic particles other than photons in
  units of that of one neutrino species whose temperature is
  \((4/11)^{1/3}\) times that of the photons. Unit: dimensionless.
  Value: 3.044 (Standard Model); measured \(2.99 \pm 0.17\) (Planck
  2018). Status: STANDARD.
\item
  \(\mathcal{T}\): The energy-momentum tensor of any one term of the
  fold (in \(\mathcal{T}^r{}_r = \mathcal{L}_{\mathrm{fold}}\)). Unit:
  GeV\(^4\). Status: STANDARD.
\item
  \(\mathcal{W}\): The warp factor: in a warped metric, the function of
  the position along one direction that multiplies the squared steps of
  other directions (its square root multiplies their lengths); in UM-4
  the function of \(\sigma\) that multiplies all seven directions of the
  fold, ordinary space, the time and the extra times,
  \(\mathcal{W} = \mathrm{e}^{-2H\lvert\sigma\rvert} = (\sin^2z)^{1/6}\),
  equal to 1 on the fold. Unit: dimensionless. Status: DERIVED.
\item
  \(n_{\mathrm{S}}\): The spectral index of the primordial curvature
  spectrum. Unit: dimensionless. Value: 0.965 (Planck 2018). Status:
  ASSUMED.
\item
  \(\theta_{12},\ \theta_{13},\ \theta_{23}\): The three mixing angles
  of \(U_{\mathrm{PMNS}}\) (its fourth parameter is the Dirac phase,
  which does not enter the electron-flavour results of UM-4). Unit:
  dimensionless. Value: \(33.40^\circ\), \(8.54^\circ\), \(49.14^\circ\)
  in the computations of UM-4 (\(\sin^2\theta_{12} = 0.303\),
  \(\sin^2\theta_{13} = 0.02203\), \(\sin^2\theta_{23} = 0.572\): NuFIT
  5.2, November 2022, normal ordering, the fit without the
  Super-Kamiokande atmospheric data); \(33.71^\circ\), \(8.56^\circ\),
  \(43.28^\circ\) in NuFIT 6.0 (2024, with the Super-Kamiokande
  atmospheric data; \(\sin^2\theta_{12} = 0.308\),
  \(\sin^2\theta_{13} = 0.02215\), \(\sin^2\theta_{23} = 0.470\)).
  Status: STANDARD.
\item
  \(\upsilon,\ \upsilon_{\mathrm{b}}\): The argument of the tunnelling
  exponent \(F\), a ratio of an energy to the bulk mass \(m_\Psi\);
  \(\upsilon_{\mathrm{b}} = \omega_{\mathrm{b}}/m_\Psi\) that of the
  particle \(\chi\). Unit: dimensionless. Status: DERIVED.
\item
  \(\xi_{\mathrm{H}}\): The coupling of the Higgs doublet to the scalar
  curvature of the fold, the term
  \(-\xi_{\mathrm{H}}\hat R\,\Phi^\dagger\Phi\) in the Lagrangian of the
  fold: an extension of UM-4, whose Higgs Lagrangian has no such term
  (\(\xi_{\mathrm{H}} = 0\)). Unit: dimensionless. Value: 0 in UM-4;
  about \(1.6\times10^4\) if the Higgs field drove inflation (ESTIMATE,
  an extension of the action). Status: ESTIMATE (an extension of the
  action).
\item
  \(F(\upsilon)\): The tunnelling exponent of the width of the particles
  \(\chi\) per unit of \(m_\Psi/H\): the integral of the local decay
  rate \(q\) over the region inside the turning point, divided by
  \(m_\Psi/H\), as a function of the ratio \(\upsilon\) of the energy to
  the bulk mass. Unit: dimensionless. Value: 0.049 for
  \(\upsilon = 0.866\); \((q_{\mathrm{b}}/\omega_{\mathrm{b}})^3/3\) for
  weak binding. Status: DERIVED.
\item
  \(U_{\mathrm{PMNS}}\): The lepton mixing matrix
  (Pontecorvo-Maki-Nakagawa-Sakata) of the three Dirac neutrinos; its
  entries \((U_{\mathrm{PMNS}})_{\alpha i}\) connect the flavour state
  \(\nu_\alpha\) with the mass state \(\nu_i\). Unit: dimensionless.
  Value: \(\theta_{12} = 33.40^\circ\), \(\theta_{23} = 49.14^\circ\),
  \(\theta_{13} = 8.54^\circ\) (NuFIT 5.2, November 2022, normal
  ordering, the fit without the Super-Kamiokande atmospheric data); the
  Dirac phase is set to zero (ASSUMED; the fit prefers \(197^\circ\)).
  Status: STANDARD, ASSUMED (the phase).
\item
  \(\mathcal{L}_{\mathrm{fold}}\): The Lagrangian of any one term of the
  fold (four-dimensional normalisation). Unit: GeV\(^4\). Status:
  STANDARD.
\item
  \(k_*\): The pivot wave number of the primordial spectrum. Unit: GeV
  (1/length). Value: 0.05 Mpc\(^{-1}\). Status: STANDARD.
\item
  \(q\): The local decay rate of a profile of dirac16complex along
  \(\sigma\) at the frequency \(\omega\),
  \(q = \sqrt{m_\Psi^2 - \omega^2\mathrm{e}^{2H\lvert\sigma\rvert}}\)
  (imaginary beyond the turning point, where the profile oscillates).
  Unit: GeV. Status: DERIVED.
\end{itemize}

\end{document}
